% Production d’écrits utilitaires : rédaction d’un CV — Français Tle Tech — Cours signé OnBuch+
% Programme officiel MINESEC. Compiler avec : tectonic N12-production-ecrits-utilitaires-redaction-cv.tex
\newcommand{\DOCMATIERE}{Français}
\newcommand{\DOCNIVEAU}{Tle Tech}
\newcommand{\DOCMODULE}{Français – classe de Terminale technique}
\newcommand{\DOCLECON}{N12}
\newcommand{\DOCTITRE}{Production d’écrits utilitaires : rédaction d’un CV}
% OnBuch+ — préambule « cours signés » ENRICHI (compilé avec tectonic / XeLaTeX)
% Système visuel « Pop » d'OnBuch+ : fond crème (jamais blanc), bordures encre
% épaisses, ombres « dures » décalées, titres Archivo Black en capitales.
% Version MATHÉMATIQUES : graphes (pgfplots/tikz), boîtes pédagogiques
% (définition, propriété, méthode, attention, le savais-tu, exercice résolu,
% exercice, TP, bilan), page de garde et signature OnBuch+ ancrée en bas.
%
% Macros attendues dans main.tex AVANT \input{preamble} :
%   \DOCMATIERE  \DOCNIVEAU  \DOCMODULE  \DOCLECON (numéro)  \DOCTITRE
\documentclass[11pt]{article}

\usepackage{fontspec}
\usepackage[french]{babel}
\usepackage[a4paper,margin=2cm,top=3.4cm,bottom=2.8cm,headheight=1.4cm,footskip=1.3cm]{geometry}
\usepackage{amsmath,amssymb,amsfonts}
\usepackage{mathtools} % \xmapsto, \coloneqq... souvent utilisés par les modèles
\usepackage{cancel} % simplifications barrées (bilans de forces)
% \abs{z} (module, valeur absolue) et \norm{u} (norme), utilisés par les modèles
\providecommand{\abs}{}\let\abs\relax\DeclarePairedDelimiter{\abs}{\lvert}{\rvert}
\providecommand{\norm}{}\let\norm\relax\DeclarePairedDelimiter{\norm}{\lVert}{\rVert}
\usepackage[table]{xcolor} % [table] : fournit aussi \rowcolors (alternance de lignes)
\usepackage{pagecolor}
\usepackage{tikz}
\usetikzlibrary{arrows.meta,positioning,calc,shapes.geometric,shapes.misc,shapes.arrows,decorations.pathmorphing,decorations.pathreplacing,decorations.markings,patterns,shadows,mindmap,trees,fit,backgrounds,matrix,intersections,angles,quotes}
\usepackage{pgfplots}
\usepackage[version=4]{mhchem} % équations nucléaires \ce{^{235}_{92}U}
\usepackage[european,siunitx]{circuitikz} % schémas électriques
\pgfplotsset{compat=1.18}
\usepgfplotslibrary{fillbetween,groupplots} % aires entre courbes (calcul intégral)
\usepackage{enumitem}
\usepackage{fancyhdr}
% [most] charge toutes les bibliothèques tcolorbox (listings, minted, raster,
% documentation...) et gonfle inutilement la mémoire TeX sur un document long
% (~300 boîtes) : on ne charge que ce qui est réellement utilisé.
\usepackage[skins,breakable]{tcolorbox}
\usepackage{graphicx}
\usepackage{colortbl}
\usepackage{booktabs}
\usepackage{tabularx}
\usepackage{multirow}
\usepackage{diagbox} % case d'en-tête diagonale des tableaux à double entrée
% Colonnes extensibles centrée / gauche / droite (C, L, R), souvent utilisées par les modèles
\newcolumntype{C}{>{\centering\arraybackslash}X}
\newcolumntype{L}{>{\raggedright\arraybackslash}X}
\newcolumntype{R}{>{\raggedleft\arraybackslash}X}
\usepackage{array}
\usepackage{multicol}
\usepackage{siunitx}
% parse-numbers=false : accepte aussi \SI{2,5 \times 10^{-3}}{...}
\sisetup{locale=FR,output-decimal-marker={,},group-minimum-digits=4,parse-numbers=false}
\DeclareSIUnit\ohm{\ensuremath{\symup{\Omega}}} % U+2126 absent de la police
\DeclareSIUnit\division{div} % réglages d'oscilloscope (ms/div, V/div)
\DeclareSIUnit\tr{tr} % tours (vitesse de rotation, ex. \SI{1500}{\tr\per\minute})
\DeclareSIUnit\molar{M} % concentration molaire (ex. \SI{3}{\molar}, \SI{1}{\micro\molar})
\DeclareSIUnit\an{an} % année (le modèle confond parfois avec \year, un compteur TeX) — ex. \SI{5}{\kilo\meter\per\an}
\DeclareSIUnit\FCFA{FCFA} % franc CFA (ex. \SI{500}{\FCFA\per\kilo\gram})
\DeclareSIUnit\degreeBrix{\textdegree Brix} % teneur en sucre d'un sirop/jus (réfractométrie)
\DeclareSIUnit\month{mois} % le modèle utilise parfois \month, un compteur TeX, comme unité de durée
\DeclareSIUnit\microliter{\micro\liter} % le modèle écrit parfois \microliter en un seul token
\DeclareSIUnit\million{millions} % dénombrement (ex. \SI{15}{\million\per\mL} spermatozoïdes)
\usepackage{qrcode}
% \degree : commande fréquemment utilisée par les modèles pour un angle ou une
% température en dehors de \SI{}{} (non fournie par les paquets chargés).
\providecommand{\degree}{^{\circ}}
% \qed (fin de démonstration), utilisé par les modèles sans amsthm
\providecommand{\qed}{\hfill$\square$}
% \barre : double barre des valeurs interdites dans un tableau de variations (tkz-tab)
\providecommand{\barre}{\ensuremath{\|}}
% environnement proof (amsthm non chargé) utilisé par les modèles
\ifdefined\proof\else\newenvironment{proof}[1][Démonstration]{\par\noindent\textit{#1.}\ }{\qed\par}\fi
\usepackage{lastpage}
\usepackage{hyperref}
\hypersetup{hidelinks}

% ============================================================
% Palette « Pop » OnBuch+ — mêmes hex que la classe P (plus_ui.dart)
% ============================================================
\definecolor{poporange}{HTML}{F59321}
\definecolor{popdark}{HTML}{E07A0C}
\definecolor{popcream}{HTML}{FFF6E8}
\definecolor{popink}{HTML}{1C1714}
\definecolor{poppurple}{HTML}{7A5AE0}
\definecolor{popgreen}{HTML}{2E9E6A}
\definecolor{poppink}{HTML}{E0457B}
\definecolor{popblue}{HTML}{2F7DE1}
\definecolor{popgold}{HTML}{C98A12}
\definecolor{popmuted}{HTML}{8A7F76}
\definecolor{popline}{HTML}{EADFD2}
\definecolor{popgray}{HTML}{8A7F76}
\colorlet{popgrey}{popgray}
% Alias tolérants (le modèle nomme parfois les couleurs différemment)
\colorlet{popwhite}{white}
\colorlet{popblack}{popink}
\colorlet{popred}{poppink}
\colorlet{popyellow}{popgold}
% Teintes claires (fonds de tableaux/encadrés) — le modèle nomme parfois
% les couleurs avec un suffixe L (ex. popblueL) sans les avoir définies.
\colorlet{poporangeL}{poporange!20}
\colorlet{popdarkL}{popdark!20}
\colorlet{poppurpleL}{poppurple!20}
\colorlet{popgreenL}{popgreen!20}
\colorlet{poppinkL}{poppink!20}
\colorlet{popblueL}{popblue!20}
\colorlet{popgoldL}{popgold!20}
\colorlet{popredL}{popred!20}
\colorlet{popgrayL}{popgray!20}
% Teintes claires pour fonds de tableaux / zones
\colorlet{popblueL}{popblue!10!white}
\colorlet{poporangeL}{poporange!14!white}
\colorlet{popgreenL}{popgreen!12!white}
\colorlet{poppurpleL}{poppurple!12!white}
\colorlet{poppinkL}{poppink!12!white}
\colorlet{popgoldL}{popgold!14!white}

\pagecolor{popcream}

% ============================================================
% Typographie
% ============================================================
\defaultfontfeatures{Ligatures=TeX}
\setmainfont{PlusJakartaSans-Medium.ttf}[Path=../../../fonts/,BoldFont=PlusJakartaSans-Bold.ttf]
% Maths en sans-serif (Fira Math) avec chiffres et lettres droites de Plus
% Jakarta, pour que les formules se fondent dans le texte.
\usepackage[math-style=ISO,bold-style=ISO]{unicode-math}
\setmathfont{FiraMath-Regular.otf}[Path=../../../fonts/]
\setmathfont{PlusJakartaSans-Medium.ttf}[Path=../../../fonts/,range={up/num,up/{Latin,latin},\mathup}]
\usepackage{icomma} % virgule décimale sans espace en mode math
% Caractères mathématiques Unicode absents des polices de texte (Plus Jakarta,
% Archivo Black) : rendus par leur équivalent mathématique, en texte comme en
% formule — sinon ils s'affichent en carré vide (ex. « Arithmétique dans ℤ »).
\usepackage{newunicodechar}
\newunicodechar{ℝ}{\ensuremath{\mathbb{R}}}
\newunicodechar{ℤ}{\ensuremath{\mathbb{Z}}}
\newunicodechar{ℂ}{\ensuremath{\mathbb{C}}}
\newunicodechar{ℕ}{\ensuremath{\mathbb{N}}}
\newunicodechar{ℚ}{\ensuremath{\mathbb{Q}}}
\newunicodechar{𝔻}{\ensuremath{\mathbb{D}}}
\newunicodechar{α}{\ensuremath{\alpha}}
\newunicodechar{β}{\ensuremath{\beta}}
\newunicodechar{γ}{\ensuremath{\gamma}}
\newunicodechar{δ}{\ensuremath{\delta}}
\newunicodechar{ε}{\ensuremath{\varepsilon}}
\newunicodechar{θ}{\ensuremath{\theta}}
\newunicodechar{λ}{\ensuremath{\lambda}}
\newunicodechar{μ}{\ensuremath{\mu}}
\newunicodechar{σ}{\ensuremath{\sigma}}
\newunicodechar{τ}{\ensuremath{\tau}}
\newunicodechar{φ}{\ensuremath{\varphi}}
\newunicodechar{ψ}{\ensuremath{\psi}}
\newunicodechar{ω}{\ensuremath{\omega}}
\newunicodechar{Σ}{\ensuremath{\Sigma}}
% majuscules produites par \MakeUppercase dans les titres (α-aminés -> Α)
\newunicodechar{Α}{\ensuremath{\alpha}}
\newunicodechar{Β}{\ensuremath{\beta}}
\newunicodechar{Γ}{\ensuremath{\Gamma}}
\newunicodechar{Δ}{\ensuremath{\Delta}}
\newunicodechar{Ε}{\ensuremath{\varepsilon}}
\newunicodechar{Θ}{\ensuremath{\Theta}}
\newunicodechar{Λ}{\ensuremath{\Lambda}}
\newunicodechar{Μ}{\ensuremath{\mu}}
\newunicodechar{Τ}{\ensuremath{\tau}}
\newunicodechar{Φ}{\ensuremath{\Phi}}
\newunicodechar{Ψ}{\ensuremath{\Psi}}
\newunicodechar{Ω}{\ensuremath{\Omega}}
\newunicodechar{↦}{\ensuremath{\mapsto}}
\newunicodechar{∈}{\ensuremath{\in}}
\newunicodechar{∉}{\ensuremath{\notin}}
\newunicodechar{∧}{\ensuremath{\wedge}}
\newunicodechar{⇔}{\ensuremath{\Leftrightarrow}}
\newunicodechar{⟺}{\ensuremath{\Longleftrightarrow}}
\newunicodechar{⇒}{\ensuremath{\Rightarrow}}
\newunicodechar{⟹}{\ensuremath{\Longrightarrow}}
\newunicodechar{∘}{\ensuremath{\circ}}
\newunicodechar{∪}{\ensuremath{\cup}}
\newunicodechar{∩}{\ensuremath{\cap}}
\newunicodechar{≡}{\ensuremath{\equiv}}
\newunicodechar{⊂}{\ensuremath{\subset}}
\newunicodechar{⊥}{\ensuremath{\perp}}
\newunicodechar{∃}{\ensuremath{\exists}}
\newunicodechar{∀}{\ensuremath{\forall}}
\newunicodechar{∛}{\ensuremath{\sqrt[3]{\,}}}
\newunicodechar{∜}{\ensuremath{\sqrt[4]{\,}}}
\newunicodechar{⃗}{\ensuremath{{}^{\to}}}
\newunicodechar{ⁿ}{\ifmmode^{n}\else\textsuperscript{\ensuremath{n}}\fi}
\newunicodechar{ˣ}{\ifmmode^{x}\else\textsuperscript{\ensuremath{x}}\fi}
\newunicodechar{ᵇ}{\ifmmode^{b}\else\textsuperscript{\ensuremath{b}}\fi}
\newunicodechar{ʳ}{\ifmmode^{r}\else\textsuperscript{\ensuremath{r}}\fi}
\newunicodechar{ᵘ}{\ifmmode^{u}\else\textsuperscript{\ensuremath{u}}\fi}
\newunicodechar{ʸ}{\ifmmode^{y}\else\textsuperscript{\ensuremath{y}}\fi}
\newunicodechar{ᵃ}{\ifmmode^{a}\else\textsuperscript{\ensuremath{a}}\fi}
\newunicodechar{ᵉ}{\ifmmode^{e}\else\textsuperscript{\ensuremath{e}}\fi}
\newunicodechar{ᵏ}{\ifmmode^{k}\else\textsuperscript{\ensuremath{k}}\fi}
\newunicodechar{ᵖ}{\ifmmode^{p}\else\textsuperscript{\ensuremath{p}}\fi}
\newunicodechar{ᴺ}{\ifmmode^{N}\else\textsuperscript{\ensuremath{N}}\fi}
\newunicodechar{ᶜ}{\ifmmode^{c}\else\textsuperscript{\ensuremath{c}}\fi}
\newunicodechar{ʰ}{\ifmmode^{h}\else\textsuperscript{\ensuremath{h}}\fi}
\newunicodechar{ᵠ}{\ifmmode^{q}\else\textsuperscript{\ensuremath{q}}\fi}
\newunicodechar{ₙ}{\ifmmode_{n}\else\textsubscript{\ensuremath{n}}\fi}
\newunicodechar{ₐ}{\ifmmode_{a}\else\textsubscript{\ensuremath{a}}\fi}
\newunicodechar{ᵢ}{\ifmmode_{i}\else\textsubscript{\ensuremath{i}}\fi}
\newunicodechar{ᵣ}{\ifmmode_{r}\else\textsubscript{\ensuremath{r}}\fi}
\newunicodechar{ₖ}{\ifmmode_{k}\else\textsubscript{\ensuremath{k}}\fi}
\newunicodechar{ᵦ}{\ifmmode_{\beta}\else\textsubscript{\ensuremath{\beta}}\fi}
\newunicodechar{ₓ}{\ifmmode_{x}\else\textsubscript{\ensuremath{x}}\fi}
\newfontfamily\popdisplay{ArchivoBlack-Regular.ttf}[Path=../../../fonts/]
\newfontfamily\poplogo{SpaceGrotesk-Bold.ttf}[Path=../../../fonts/]
\newfontfamily\popsemi{PlusJakartaSans-SemiBold.ttf}[Path=../../../fonts/]
\newfontfamily\popxbold{PlusJakartaSans-ExtraBold.ttf}[Path=../../../fonts/]
\newfontfamily\popxboldls{PlusJakartaSans-ExtraBold.ttf}[Path=../../../fonts/,LetterSpace=3.0]

\setlist[enumerate]{leftmargin=1.7em,itemsep=5pt,topsep=4pt}
\setlist[itemize]{leftmargin=1.7em,itemsep=4pt,topsep=4pt,label={\color{poporange}\textbullet}}
\setlength{\parindent}{0pt}
\setlength{\parskip}{5pt}
\renewcommand{\arraystretch}{1.25}

% pgfplots : style Pop par défaut
\pgfplotsset{
  every axis/.append style={axis line style={popink,line width=1pt},tick style={popink},
    grid style={popline,line width=0.5pt},label style={font=\small},tick label style={font=\footnotesize}},
  popcurve/.style={poporange,line width=1.8pt,no markers,smooth},
}
% Même style disponible dans un \draw TikZ simple (hors pgfplots)
\tikzset{popcurve/.style={poporange,line width=1.8pt}}
\tikzset{popgrid/.style={popline,very thin}}
\colorlet{popcurve}{poporange} % le nom du style est parfois utilisé comme couleur

% ============================================================
% Pill / squiggle
% ============================================================
\newcommand{\poppill}[3]{%
  \tikz[baseline=-0.55ex]\node[draw=popink,line width=1.2pt,rounded corners=2.6pt,%
    fill=#2,inner xsep=6.5pt,inner ysep=3pt]{%
    \popxboldls\fontsize{7}{7}\selectfont\color{#3}\MakeUppercase{#1}};%
}
\newcommand{\popsquiggle}[1]{%
  \tikz[baseline=-0.4ex]{
    \draw[#1,line width=2.6pt,line cap=round]
      plot [smooth] coordinates {(0,0.09) (0.34,0.01) (0.68,0.21) (1.02,0.01) (1.36,0.10)};
  }%
}
% Mot-clé surligné (marqueur orange sous le texte)
\newcommand{\cle}[1]{{\popxbold\color{popdark}#1}}

% ============================================================
% En-tête / pied
% ============================================================
\pagestyle{fancy}
\fancyhf{}
\renewcommand{\headrulewidth}{0pt}
\renewcommand{\footrulewidth}{0pt}
\lhead{\raisebox{-0.15ex}{\poplogo\fontsize{13}{13}\selectfont\color{popink}OnBuch\color{poporange}+}}
\rhead{\poppill{\DOCMATIERE\ \textbullet\ \DOCNIVEAU\ \textbullet\ Leçon \DOCLECON}{white}{popink}}
\lfoot{\popxbold\fontsize{7.6}{7.6}\selectfont\color{popmuted}\MakeUppercase{Cours signé OnBuch+}\ \textbullet\ {\popsemi\DOCTITRE}}
\rfoot{\tikz[baseline=-0.6ex]\node[rounded corners=8.5pt,fill=popink,minimum height=17pt,minimum width=17pt,inner xsep=5pt,inner sep=0]{\popxbold\fontsize{8}{8}\selectfont\color{popcream}\thepage\,/\,\pageref*{LastPage}};}

% ============================================================
% Titres
% ============================================================
\newcommand{\chaptitre}[2]{%
  \begin{tcolorbox}[enhanced,colback=poporange,colframe=popink,boxrule=2.6pt,arc=11pt,
      drop shadow=popink,left=15pt,right=15pt,top=12pt,bottom=11pt,before skip=3pt,after skip=18pt]
    \poppill{#2}{popink}{popcream}\par\vspace{9pt}
    {\raggedright\hyphenpenalty=10000\exhyphenpenalty=10000\popdisplay\fontsize{22}{24}\selectfont\color{popcream}\MakeUppercase{#1}\par}
  \end{tcolorbox}
}
% Section numérotée : pastille orange avec le numéro + titre Archivo
\newcounter{popsec}
\newcommand{\coursec}[1]{%
  \stepcounter{popsec}\setcounter{popsub}{0}%
  \par\vspace{16pt}\noindent
  \begin{minipage}{\linewidth}
  \tikz[baseline=-0.7ex]\node[draw=popink,line width=1.6pt,fill=poporange,rounded corners=4pt,minimum size=22pt,inner sep=0]{\popdisplay\fontsize{12}{12}\selectfont\color{popcream}\thepopsec};%
  \hspace{8pt}{\popdisplay\fontsize{15}{17}\selectfont\color{popink}\MakeUppercase{#1}}\par
  \vspace{4pt}\noindent{\color{poporange}\rule{36pt}{3.6pt}}
  \end{minipage}\par\nopagebreak\vspace{8pt}
}
\newcounter{popsub}
\newcommand{\courssub}[1]{%
  \stepcounter{popsub}%
  \par\vspace{9pt}\noindent{\popxbold\fontsize{11.5}{13}\selectfont\color{popink}{\color{poporange}\thepopsec.\thepopsub}\hspace{6pt}#1}\par\nopagebreak\vspace{4pt}
}

% ============================================================
% Boîtes pédagogiques — même recette : fond blanc, bordure épaisse colorée,
% ombre dure encre, badge plein. Titre optionnel [#1].
% ============================================================
\tcbset{popbase/.style={breakable,enhanced,colback=white,boxrule=2.3pt,arc=9pt,
  shadow={3pt}{-3pt}{0pt}{popink},left=14pt,right=14pt,top=11pt,bottom=11pt,before skip=11pt,after skip=15pt,
  fontupper=\popsemi\fontsize{10.6}{14.5}\selectfont\color{popink},
  fontlower=\popsemi\fontsize{10.6}{14.5}\selectfont\color{popink}}}
% Coche dessinée (le glyphe ✓ n'existe pas dans les polices du document)
\newcommand{\popcheck}[1]{\tikz[baseline=-0.5ex]\draw[#1,line width=1.6pt,line cap=round,line join=round](-0.13,0.0)--(-0.04,-0.09)--(0.13,0.1);}
\providecommand{\coche}{\popcheck{popgreen}}
\providecommand{\resultat}[1]{\par\smallskip\noindent\fcolorbox{popgreen}{popgreenL}{\parbox{\dimexpr\linewidth-2\fboxsep-2\fboxrule\relax}{\textbf{Résultat.} #1}}\par\smallskip}
\newcommand{\popboxhead}[3]{\poppill{#1}{#2}{white}\ifx\relax#3\relax\else\hspace{6pt}{\popxbold\fontsize{10.6}{12}\selectfont\color{popink}#3}\fi\par\vspace{7pt}}

\newtcolorbox{aretenir}[1][]{popbase,colframe=popgreen,before upper={\popboxhead{À retenir}{popgreen}{#1}}}
\newtcolorbox{exemplebox}[1][]{popbase,colframe=poporange,before upper={\popboxhead{Exemple}{poporange}{#1}}}
\newtcolorbox{definition}[1][]{popbase,colframe=popblue,before upper={\popboxhead{Définition}{popblue}{#1}}}
\newtcolorbox{propriete}[1][]{popbase,colframe=poppurple,before upper={\popboxhead{Propriété}{poppurple}{#1}}}
\newtcolorbox{methode}[1][]{popbase,colframe=popgold,before upper={\popboxhead{Méthode}{popgold}{#1}}}
\newtcolorbox{attention}[1][]{popbase,colframe=poppink,before upper={\popboxhead{Attention}{poppink}{#1}}}
\newtcolorbox{experience}[1][]{popbase,colframe=popblue,colback=popblueL,before upper={\popboxhead{Expérience}{popblue}{#1}}}
% « Le savais-tu ? » : carte violette pleine (contexte, culture, Cameroun)
\newtcolorbox{savaistu}[1][]{popbase,colback=poppurple,colframe=popink,
  fontupper=\popsemi\fontsize{10.4}{14.2}\selectfont\color{white},
  before upper={\poppill{Le savais-tu\,?}{white}{poppurple}\ifx\relax#1\relax\else\hspace{6pt}{\popxbold\color{white}#1}\fi\par\vspace{7pt}}}
% Exercice résolu : énoncé en haut, \tcblower puis la solution sur fond crème
\newcounter{popexo}\newcounter{popexr}
\newtcolorbox{exoresolu}[1][]{popbase,colframe=poporange,colbacklower=popcream,
  segmentation style={popink,line width=1.2pt,dashed},
  before upper={\stepcounter{popexr}\popboxhead{Exercice résolu \thepopexr}{poporange}{#1}},
  before lower={\poppill{Solution}{popink}{popcream}\par\vspace{6pt}}}
% Exercice (non résolu en place) — la correction est dans la partie Corrigés
\newtcolorbox{exercice}[1][]{popbase,colframe=popink,
  before upper={\stepcounter{popexo}\popboxhead{Exercice \thepopexo}{popink}{#1}}}
% Corrigé
\newtcolorbox{corrige}[1][]{popbase,colframe=popgreen,colback=popgreenL,
  before upper={\popboxhead{Corrigé}{popgreen}{#1}}}
% Objectifs / prérequis / bilan
\newtcolorbox{objectifs}{popbase,colframe=popink,colback=poporangeL,
  before upper={\popboxhead{Objectifs de la leçon}{popink}{}}}
\newtcolorbox{prerequis}{popbase,colframe=popink,colback=popblueL,
  before upper={\popboxhead{Prérequis}{popblue}{}}}
\newtcolorbox{bilan}{popbase,colframe=popink,colback=white,boxrule=2.8pt,
  before upper={\popboxhead{Fiche bilan}{poporange}{L'essentiel à connaître pour l'examen}}}
% Figure encadrée avec légende
\newtcolorbox{popfigure}[1][]{enhanced,breakable=false,colback=white,colframe=popline,boxrule=1.4pt,arc=8pt,
  left=10pt,right=10pt,top=10pt,bottom=8pt,before skip=10pt,after skip=12pt,halign=center,
  lower separated=false,halign lower=center,
  fontlower=\popsemi\fontsize{9}{11.5}\selectfont\color{popmuted},
  #1}
\newcommand{\legende}[1]{\par\vspace{6pt}{\popsemi\fontsize{9}{11.5}\selectfont\color{popmuted}#1}}

% ============================================================
% Signature OnBuch+
% ============================================================
\newcommand{\poplogoinline}[1]{{\poplogo\fontsize{#1}{#1}\selectfont\color{popink}OnBuch\color{poporange}+}}
\newcommand{\signatureonbuch}{%
  \begin{tcolorbox}[enhanced,colback=popink,colframe=popink,boxrule=2.2pt,arc=10pt,
      drop shadow=poporange,left=14pt,right=14pt,top=10pt,bottom=10pt,before skip=0pt,after skip=0pt]
    \tikz[baseline=-0.7ex]\node[circle,fill=popgreen,draw=popcream,line width=1.3pt,minimum size=22pt,inner sep=0]{\popcheck{white}};%
    \hspace{9pt}\begin{minipage}[c]{0.62\linewidth}
      {\popxbold\fontsize{12}{13}\selectfont\color{popcream}Signé\ \textbullet\ {\poplogo OnBuch\color{poporange}+}}\par\vspace{2pt}
      {\raggedright\popsemi\fontsize{8}{10.5}\selectfont\color{popline}\MakeUppercase{Équipe pédagogique OnBuch+}\\\MakeUppercase{Conforme au programme officiel MINESEC}\par}
    \end{minipage}\hfill
    {\poplogo\fontsize{22}{22}\selectfont\color{popcream}OnBuch\color{poporange}+}
  \end{tcolorbox}%
}

% ============================================================
% Page de garde
% #1 titre · #2 matière · #3 niveau · #4 module · #5 n° leçon · #6 durée ·
% #7 description / objectifs (texte ou itemize)
% ============================================================
\newcommand{\pagegarde}[7]{%
  \thispagestyle{empty}
  \newgeometry{a4paper,margin=1.7cm,top=1.6cm,bottom=1.4cm}%
  \begin{tikzpicture}[remember picture,overlay]
    \fill[poporange!18] ([xshift=-3.2cm,yshift=-3.6cm]current page.north east) circle (4.6cm);
    \fill[poppurple!14] ([xshift=2.4cm,yshift=7.6cm]current page.south west) circle (3.8cm);
  \end{tikzpicture}%
  {\poplogo\fontsize{30}{30}\selectfont\color{popink}OnBuch\color{poporange}+}\hfill
  \raisebox{0.6ex}{\poppill{Cours signé \textbullet\ Premium}{popink}{popcream}}\par
  \vspace{1pt}\popsquiggle{poppurple}\par
  \vspace{8pt}
  \begin{tcolorbox}[enhanced,colback=poporange,colframe=popink,boxrule=2.8pt,arc=13pt,
      drop shadow=popink,left=18pt,right=18pt,top=12pt,bottom=13pt,after skip=10pt]
    \poppill{#2 \textbullet\ #3}{popink}{popcream}\hspace{5pt}\poppill{#4}{white}{popink}\par\vspace{8pt}
    {\popdisplay\fontsize{14}{14}\selectfont\color{popink}LEÇON #5}\par\vspace{4pt}
    {\raggedright\hyphenpenalty=10000\exhyphenpenalty=10000\popdisplay\fontsize{25}{27}\selectfont\color{popcream}\MakeUppercase{#1}\par}
    \vspace{8pt}
    {\popxbold\fontsize{9.5}{9.5}\selectfont\color{popink}\MakeUppercase{Durée conseillée : #6}}
  \end{tcolorbox}
  \begin{tcolorbox}[enhanced,colback=white,colframe=popink,boxrule=2.2pt,arc=10pt,
      drop shadow=popink,left=16pt,right=16pt,top=10pt,bottom=10pt,after skip=8pt]
    \poppill{Dans cette leçon}{poppurple}{white}\par\vspace{5pt}
    \popsemi\fontsize{9.3}{12.6}\selectfont\color{popink}#7
  \end{tcolorbox}
  \noindent\begin{minipage}[t]{0.487\linewidth}
    \begin{tcolorbox}[enhanced,colback=popgreen,colframe=popink,boxrule=2.2pt,arc=10pt,
        drop shadow=popink,left=11pt,right=11pt,top=10pt,bottom=10pt,height=3.1cm,valign=center]
      \tcbox[colback=white,colframe=popink,boxrule=1.3pt,arc=5pt,boxsep=0pt,nobeforeafter,
        left=3pt,right=3pt,top=3pt,bottom=3pt,valign=center]{\qrcode[height=1.9cm]{https://play.google.com/store/apps/details?id=com.luvvix.onbuch&hl=fr}}%
      \hspace{8pt}\begin{minipage}[c]{0.5\linewidth}
        {\popdisplay\fontsize{11}{12.5}\selectfont\color{white}\MakeUppercase{Télécharge OnBuch}}\par\vspace{4pt}
        {\raggedright\popsemi\fontsize{8.4}{11}\selectfont\color{white}Gratuit \textbullet\ Google Play\\Tuteur IA, annales, concours, résultats\par}
      \end{minipage}
    \end{tcolorbox}
  \end{minipage}\hfill
  \begin{minipage}[t]{0.487\linewidth}
    \begin{tcolorbox}[enhanced,colback=poppurple,colframe=popink,boxrule=2.2pt,arc=10pt,
        drop shadow=popink,left=11pt,right=11pt,top=10pt,bottom=10pt,height=3.1cm,valign=center]
      \tikz[baseline=-0.7ex]\node[circle,fill=white,draw=popink,line width=1.3pt,minimum size=1.6cm,inner sep=0]{\popdisplay\fontsize{24}{24}\selectfont\color{poppurple}+};%
      \hspace{8pt}\begin{minipage}[c]{0.6\linewidth}
        {\popdisplay\fontsize{11}{12.5}\selectfont\color{white}\MakeUppercase{Passe à OnBuch+}}\par\vspace{4pt}
        {\raggedright\popsemi\fontsize{8.4}{11}\selectfont\color{white}Tous les cours signés, Tuteur IA illimité, fiches PDF — onglet OnBuch+ de l'app\par}
      \end{minipage}
    \end{tcolorbox}
  \end{minipage}
  \vspace{8pt}
  \popcontenu
  \vfill
  \signatureonbuch
  \restoregeometry
  \newpage
}

% Rangée « Ce cours contient » (page de garde)
\newcommand{\poptile}[3]{% #1 couleur #2 gros texte #3 légende
  \begin{tcolorbox}[enhanced,colback=#1,colframe=popink,boxrule=2pt,arc=9pt,shadow={2.5pt}{-2.5pt}{0pt}{popink},
      left=6pt,right=6pt,top=9pt,bottom=9pt,halign=center,valign=center,height=1.75cm,width=\linewidth,nobeforeafter]
    {\popdisplay\fontsize{12.5}{13}\selectfont\color{white}\MakeUppercase{#2}}\par\vspace{3pt}
    {\centering\popsemi\fontsize{7.8}{9.5}\selectfont\color{white}#3\par}
  \end{tcolorbox}}
\newcommand{\popcontenu}{%
  \par\noindent{\popxboldls\fontsize{7.5}{7.5}\selectfont\color{popmuted}CE COURS SIGNÉ CONTIENT}\par\vspace{7pt}
  \noindent\begin{minipage}[t]{0.235\linewidth}\poptile{popblue}{Cours}{complet, illustré, pas à pas}\end{minipage}\hfill
  \begin{minipage}[t]{0.235\linewidth}\poptile{poporange}{Méthodes}{et exercices résolus}\end{minipage}\hfill
  \begin{minipage}[t]{0.235\linewidth}\poptile{poppink}{Exercices}{type examen + corrigés}\end{minipage}\hfill
  \begin{minipage}[t]{0.235\linewidth}\poptile{popgreen}{Fiche bilan}{l'essentiel à retenir}\end{minipage}\par}

% Sommaire
\newtcolorbox{sommaire}{enhanced,colback=white,colframe=popink,boxrule=2.2pt,arc=10pt,
  drop shadow=popink,left=18pt,right=18pt,top=13pt,bottom=13pt,before skip=6pt,after skip=20pt,
  before upper={\poppill{Sommaire}{poppurple}{white}\par\vspace{9pt}\popsemi\fontsize{10.6}{14.5}\selectfont\color{popink}}}

% Signature de fin de cours — poussée tout en bas de la dernière page
\newcommand{\findecours}{%
  \par\vfill\nopagebreak
  \begin{center}\popsquiggle{poporange}\end{center}\vspace{4pt}
  \signatureonbuch
}
\AtBeginDocument{\def\llbracket{\mathopen{[\mkern-3.5mu[}}\def\rrbracket{\mathclose{]\mkern-3.5mu]}}} % intervalles d'entiers (glyphe ⟦ absent de la police)
% Glyphes absents de Fira Math : redéfinis à partir de glyphes disponibles
\AtBeginDocument{%
  \def\setminus{\mathbin{\backslash}}\let\smallsetminus\setminus
  \def\vdots{\mathord{\vbox{\baselineskip3.2pt\lineskiplimit0pt\kern4pt\hbox{.}\hbox{.}\hbox{.}}}}%
  \def\ddots{\mathinner{\mkern1mu\raise6.5pt\hbox{.}\mkern2mu\raise3.6pt\hbox{.}\mkern2mu\raise0.7pt\hbox{.}\mkern1mu}}%
  \def\triangle{\mathord{\scalebox{0.9}{$\symup{\Delta}$}}}%
  \def\nabla{\mathord{\raisebox{\depth}{\scalebox{1}[-1]{$\symup{\Delta}$}}}}%
  \def\checkmark{\popcheck{popdark}}%
  \def\star{\mathbin{\ast}}\def\bigstar{\mathord{\ast}}%
  \def\diamond{\mathbin{\scalebox{0.6}{$\Diamond$}}}%
  \def\bigcup{\mathop{\mathchoice{\raisebox{-0.3em}{\scalebox{1.7}{$\cup$}}}{\scalebox{1.25}{$\cup$}}{\cup}{\cup}}\displaylimits}%
  \def\bigcap{\mathop{\mathchoice{\raisebox{-0.3em}{\scalebox{1.7}{$\cap$}}}{\scalebox{1.25}{$\cap$}}{\cap}{\cap}}\displaylimits}%
  \def\lesseqqgtr{\mathrel{\lessgtr}}\def\bigoplus{\mathop{\oplus}\displaylimits}%
}
\providecommand{\ssi}{si et seulement si} % abréviation fréquente
\AtBeginDocument{\providecommand{\wideparen}{\overparen}} % arc de cercle

%% FIGFIX-BEGIN v5 — correctifs globaux des figures (docs/onbuch-generation/tools/figfix)
\usepackage{adjustbox}\usepackage{etoolbox}\tcbuselibrary{raster}
% toute figure TikZ/pgfplots plus large que la ligne est réduite à la largeur disponible
\BeforeBeginEnvironment{tikzpicture}{\begin{adjustbox}{max width=\linewidth}}
\AfterEndEnvironment{tikzpicture}{\end{adjustbox}}
% légendes pgfplots lisibles par-dessus les courbes
\pgfplotsset{every axis/.append style={legend style={fill=white,fill opacity=0.92,text opacity=1,draw=black!15,inner sep=3pt}}}
% étiquettes d'arêtes (syntaxe edge["..."]) avec fond blanc
\tikzset{every edge quotes/.append style={fill=white,inner sep=1.2pt}}
%% FIGFIX-END
\begin{document}

\pagegarde{Production d’écrits utilitaires : rédaction d’un CV}{Français}{Tle Tech}{Français – classe de Terminale technique}{N12}{12 séances (fiche de progression harmonisée)}{Cette leçon de production d'écrits utilitaires apprend aux élèves de Terminale technique à lire, analyser et rédiger un curriculum vitae efficace. À partir de modèles authentiques et de situations d'embauche camerounaises, elle conduit chaque élève à produire, améliorer et adapter son propre CV.
\vspace{4pt}
\begin{multicols}{2}\small
\begin{itemize}
\item À la fin de cette leçon, je sais identifier les caractéristiques d'un texte fonctionnel et d'un texte iconique (annonce d'offre d'emploi, modèle de CV).
\item À la fin de cette leçon, je sais décrire la structure externe d'un CV (titre, état civil, formation, expérience, compétences, centres d'intérêt, références).
\item À la fin de cette leçon, je sais respecter la structure interne d'un CV : ordre logique, concision, cohérence, mise en page aérée.
\item À la fin de cette leçon, je sais employer correctement les valeurs de l'impératif et de l'infinitif dans les consignes et les descriptions de tâches.
\item À la fin de cette leçon, je sais rédiger un CV complet à partir de ma propre situation scolaire et professionnelle.
\item À la fin de cette leçon, je sais adapter mon CV à une offre d'embauche précise en sélectionnant les informations pertinentes.
\end{itemize}\end{multicols}}

\begin{sommaire}
\begin{multicols}{2}
\begin{enumerate}
\item Découverte : lire des textes fonctionnels et iconiques
\item La structure externe du CV
\item La structure interne du CV
\item Langue : les valeurs de l'impératif et de l'infinitif
\item Production guidée : rédiger et améliorer son CV
\item Adapter le CV à une situation d'embauche
\item Compte rendu, remédiation et évaluation
\item Activité d'intégration
\item Méthodes et exercices résolus
\item Exercices
\item Corrigés des exercices
\item Fiche bilan
\end{enumerate}
\end{multicols}
\end{sommaire}

\begin{objectifs}
\begin{itemize}
\item À la fin de cette leçon, je sais identifier les caractéristiques d'un texte fonctionnel et d'un texte iconique (annonce d'offre d'emploi, modèle de CV).
\item À la fin de cette leçon, je sais décrire la structure externe d'un CV (titre, état civil, formation, expérience, compétences, centres d'intérêt, références).
\item À la fin de cette leçon, je sais respecter la structure interne d'un CV : ordre logique, concision, cohérence, mise en page aérée.
\item À la fin de cette leçon, je sais employer correctement les valeurs de l'impératif et de l'infinitif dans les consignes et les descriptions de tâches.
\item À la fin de cette leçon, je sais rédiger un CV complet à partir de ma propre situation scolaire et professionnelle.
\item À la fin de cette leçon, je sais adapter mon CV à une offre d'embauche précise en sélectionnant les informations pertinentes.
\item À la fin de cette leçon, je sais analyser une situation de communication (destinataire, objectif, contexte) avant de rédiger.
\item À la fin de cette leçon, je sais confronter ma production à celle des autres, la corriger et l'améliorer.
\item À la fin de cette leçon, je sais rédiger et justifier une démarche, une réponse ou une conclusion dans un écrit utilitaire.
\end{itemize}
\end{objectifs}

\begin{prerequis}
\begin{itemize}
\item Les types de textes et les textes fonctionnels étudiés en classe de Première
\item La phrase simple et la phrase complexe (morphosyntaxe de base)
\item Les modes verbaux déjà rencontrés : indicatif, subjonctif, conditionnel
\item La notion de situation de communication (émetteur, récepteur, message, contexte)
\item La lettre administrative et la lettre de motivation (écrits utilitaires antérieurs)
\end{itemize}
\end{prerequis}

\newpage

\coursec{Découverte : lire des textes fonctionnels et iconiques}

\courssub{Situation d'accroche et problématique}

À Douala, une entreprise de BTP publie une offre d'emploi de technicien en génie civil : 300 candidatures arrivent, mais le recruteur ne retient que 10 CV en moins de cinq minutes. Pourquoi certains CV finissent-ils directement à la corbeille alors que les candidats ont les mêmes diplômes ? Quelles informations faut-il mettre en avant, dans quel ordre et avec quelles formulations pour convaincre un employeur camerounais ? Cette leçon nous apprend à rédiger un curriculum vitae clair, structuré et adapté à une situation d'embauche réelle.

Avant d'écrire, il faut savoir lire. Toute la leçon repose sur une idée simple : le CV du candidat répond à l'offre d'emploi de l'employeur. Nous allons donc d'abord apprendre à lire ces deux familles de textes : les \cle{textes fonctionnels} (l'offre d'emploi, le CV) et les \cle{textes iconiques} (les CV dont la mise en page porte une partie du message).

\courssub{Le texte fonctionnel : un écrit qui sert à agir}

Dans la vie quotidienne, nous lisons et écrivons sans cesse des textes qui ne sont ni des romans ni des poèmes : une annonce sur un panneau d'affichage, une facture d'électricité, un formulaire d'inscription à l'examen, une convocation, un CV. Tous ces écrits ont un point commun : ils ne cherchent pas à faire rêver, ils cherchent à \emph{faire agir} ou à accompagner une action concrète.

\begin{definition}[Texte fonctionnel]
Un \cle{texte fonctionnel} (ou texte utilitaire) est un écrit qui accompagne une action de la vie quotidienne et répond à un besoin précis : informer, demander, prouver, vendre, recruter. Il se reconnaît à trois traits :
\begin{itemize}
\item un \textbf{objectif pratique} clair (obtenir un emploi, payer une dette, s'inscrire) ;
\item des \textbf{informations exactes et vérifiables} (noms, dates, chiffres, adresses) ;
\item une \textbf{forme codifiée} : titre, rubriques, formules stéréotypées (« Prière de... », « Date limite : ... »).
\end{itemize}
Exemples : l'annonce, l'offre d'emploi, la facture, le formulaire, la lettre administrative, le CV.
\end{definition}

\begin{attention}[Ne pas confondre les deux partenaires]
Erreur fréquente : confondre l'\textbf{offre d'emploi} et le \textbf{CV}. L'offre d'emploi est le texte de \emph{l'employeur} : il décrit le poste et le profil recherché. Le CV est le texte du \emph{candidat} : il présente son parcours. Les deux textes \emph{se répondent} : un bon CV reprend, point par point, les attentes formulées dans l'offre. Celui qui rédige son CV sans avoir lu attentivement l'offre écrit à l'aveugle.
\end{attention}

\courssub{Lecture guidée du texte fonctionnel N° 3 : une offre d'emploi}

Lisons ensemble l'offre suivante, publiée dans \emph{Cameroon Tribune} et affichée au panneau d'une entreprise de Yaoundé.

\begin{exemplebox}[Texte fonctionnel N° 3 : offre d'emploi]
\textbf{SOCIÉTÉ CAMEROUNAISE DE BÂTIMENT (SOCABAT) — RECRUTE}\par
\smallskip
\textbf{UN(E) TECHNICIEN(NE) EN GÉNIE CIVIL}\par
\smallskip
\textbf{Missions :} suivi de chantier, lecture de plans, contrôle des travaux, rédaction de rapports.\par
\textbf{Profil exigé :} être titulaire d'un Baccalauréat technique (série F) ou d'un BTS en génie civil ; avoir au moins deux (02) ans d'expérience sur chantier ; maîtriser les logiciels de dessin (AutoCAD) ; être âgé de 35 ans au plus ; être disponible immédiatement.\par
\textbf{Dossier à fournir :} lettre de motivation, CV détaillé, photocopie des diplômes, attestation de travail.\par
\textbf{Lieu de travail :} Yaoundé, avec déplacements à l'intérieur du pays.\par
\textbf{Date limite de dépôt :} le 30 novembre 2026 à 15 heures, à la Direction du personnel, immeuble SOCABAT, quartier Elig-Essono, Yaoundé.
\end{exemplebox}

Comment lire ce texte efficacement ? Pas comme un roman, de la première à la dernière ligne, mais comme un document à \emph{découper en informations clés}.

\begin{methode}[Lire une offre d'emploi en trois gestes]
\begin{enumerate}
\item \textbf{Souligner le poste} : intitulé exact du métier, missions proposées. C'est la cible de votre candidature.
\item \textbf{Encadrer le profil exigé} : diplômes, expérience, compétences techniques, conditions d'âge ou de disponibilité. C'est la liste de ce que votre CV devra prouver.
\item \textbf{Entourer les modalités pratiques} : lieu de travail, pièces à fournir, adresse de dépôt et \textbf{date limite}. Une candidature déposée le 1\textsuperscript{er} décembre pour l'offre ci-dessus est irrecevable, même si le dossier est excellent.
\end{enumerate}
\end{methode}

Appliquons la méthode au texte N° 3. Le \emph{poste} : technicien(ne) en génie civil, avec quatre missions précises. Le \emph{profil} : Bac technique F ou BTS, deux ans d'expérience, AutoCAD, 35 ans maximum. Les \emph{modalités} : quatre pièces, Yaoundé, date limite du 30 novembre 2026 à 15 heures. Chacune de ces informations devra trouver un écho dans le CV du candidat.

\begin{popfigure}
\begin{center}
\begin{tikzpicture}[font=\small]
\node[draw=popdark, fill=popcream, rounded corners, text width=7.2cm, align=left, inner sep=8pt] (offre) at (0,0)
{\textbf{\textcolor{popblue}{SOCABAT — RECRUTE}}\\
\textbf{UN(E) TECHNICIEN(NE) EN GÉNIE CIVIL}\\[2pt]
Missions : suivi de chantier, lecture de plans, contrôle des travaux.\\[2pt]
Profil : Bac technique F ou BTS ; 2 ans d'expérience ; AutoCAD ; 35 ans max.\\[2pt]
Dossier : lettre de motivation, CV, diplômes, attestation.\\
Lieu : Yaoundé. \quad Date limite : 30/11/2026 à 15 h.};
\node[draw=popblue, fill=popblueL, rounded corners, text width=3.4cm, align=center, inner sep=5pt] (poste) at (-6.2,2.2) {\textbf{1. Le poste}\\ intitulé + missions};
\node[draw=poporange, fill=poporangeL, rounded corners, text width=3.4cm, align=center, inner sep=5pt] (profil) at (6.2,2.2) {\textbf{2. Le profil exigé}\\ diplômes, expérience, compétences};
\node[draw=popgreen, fill=popgreenL, rounded corners, text width=3.4cm, align=center, inner sep=5pt] (cond) at (0,-3.4) {\textbf{3. Les conditions}\\ pièces, lieu, date limite};
\draw[-{Stealth}, thick, popblue] (poste.east) -- (offre.north west);
\draw[-{Stealth}, thick, poporange] (profil.west) -- (offre.north east);
\draw[-{Stealth}, thick, popgreen] (cond.north) -- (offre.south);
\end{tikzpicture}
\end{center}
\legende{Lecture guidée d'une offre d'emploi : trois zones à repérer. 1 : le poste (souligné). 2 : le profil exigé (encadré). 3 : les conditions pratiques et la date limite (entourées).}
\end{popfigure}

\begin{exoresolu}[Repérer les informations clés d'une offre]
\textbf{Énoncé.} Dans l'offre SOCABAT ci-dessus, réponds aux questions : a) Quel diplôme minimum est exigé ? b) Un candidat de 38 ans peut-il postuler ? c) Quelles sont les quatre pièces du dossier ? d) Pourquoi la mention « disponible immédiatement » est-elle importante pour la rédaction du CV ?
\tcblower
\textbf{Solution.} a) Le diplôme minimum est le Baccalauréat technique série F (ou un BTS en génie civil). b) Non : l'offre fixe la limite à 35 ans au plus ; un candidat de 38 ans ne remplit pas la condition d'âge. c) Les quatre pièces sont : la lettre de motivation, le CV détaillé, la photocopie des diplômes et l'attestation de travail. d) « Disponible immédiatement » signifie que l'employeur veut un recrutement rapide ; le candidat doit donc indiquer dans son CV qu'il est libre de tout engagement, par exemple dans la rubrique « Divers » ou « Informations complémentaires ». C'est une information de l'offre à laquelle le CV doit répondre explicitement.
\end{exoresolu}

\courssub{Le texte fonctionnel N° 4 : un premier modèle de CV}

Le second texte fonctionnel de notre leçon est un CV. Avant de le lire en détail (ce sera l'objet des sections suivantes), adoptons la bonne démarche : \textbf{l'observation globale d'abord, la lecture détaillée ensuite}.

\begin{exemplebox}[Texte fonctionnel N° 4 : premier modèle de CV (extrait)]
\textbf{MVOGO ATEBA Jean-Claude}\par
Né le 12 mars 2004 à Ebolowa — Célibataire — Tél. : 6 90 00 00 00 — jeanclaude.mvogo@mail.cm\par
\smallskip
\textbf{DIPLÔMES} : 2025 : Baccalauréat technique F4 (Génie civil), Lycée technique d'Ebolowa. 2022 : Probatoire technique F4.\par
\textbf{EXPÉRIENCE PROFESSIONNELLE} : 2025-2026 : technicien stagiaire, Entreprise ETS Nkoulou BTP, Yaoundé (suivi de chantier, métrés). 2024 : stage ouvrier, mairie d'Ebolowa.\par
\textbf{COMPÉTENCES} : AutoCAD, lecture de plans, Word, Excel. Langues : français, anglais (niveau scolaire), ewondo.\par
\textbf{DIVERS} : disponible immédiatement ; permis B ; sport : football.
\end{exemplebox}

L'observation globale nous apprend déjà beaucoup : le texte est organisé en \textbf{rubriques} (état civil, diplômes, expérience, compétences, divers) ; chaque rubrique contient des informations \textbf{brèves et datées} ; aucune phrase longue, aucun récit. Le CV est bien un texte fonctionnel : son objectif pratique est d'obtenir un entretien d'embauche, ses informations sont vérifiables (dates, diplômes, employeurs), sa forme est codifiée.

\begin{savaistu}[Moins d'une minute pour convaincre]
Des études sur le recrutement montrent qu'un recruteur consacre en moyenne \textbf{moins d'une minute} — souvent trente à quarante secondes — à la première lecture d'un CV. Il ne lit pas : il \emph{balaye} le document du regard, à la recherche du diplôme, de l'expérience et des compétences annoncés dans son offre. C'est dire que la clarté visuelle est décisive : une information noyée dans un paragraphe dense est une information perdue. Au Cameroun comme ailleurs, le CV qui survit au premier tri est celui qui se comprend d'un coup d'œil.
\end{savaistu}

\courssub{Le texte iconique : quand la forme porte le sens}

Observons maintenant deux CV qui contiennent \emph{exactement les mêmes informations} sur le même candidat, mais présentés différemment : l'un en une seule colonne continue, l'autre en deux colonnes avec encadrés, puces et photo. Le second se lit plus vite, met en valeur les rubriques, donne une impression de sérieux et d'organisation. La différence ne tient pas aux mots : elle tient à l'\textbf{image}.

\begin{definition}[Texte iconique]
Un \cle{texte iconique} est un texte dont la signification repose en partie sur l'organisation visuelle et les éléments graphiques : disposition en colonnes, encadrés, puces, taille et graisse des caractères, couleurs, photo, pictogrammes. La mise en page n'est pas un décor : elle \emph{classe} les informations, elle indique au lecteur ce qui est important, elle guide son regard. Le CV, l'affiche, le formulaire et la page de publicité sont des textes iconiques.
\end{definition}

\courssub{Lecture des textes iconiques N° 3 et N° 4 : deux présentations du CV}

Il existe deux grandes façons d'organiser les informations d'un CV.

\textbf{Texte iconique N° 3 : le CV chronologique.} Les diplômes et les expériences sont rangés par \textbf{ordre de dates}, de la plus récente à la plus ancienne (on parle d'ordre antichronologique). C'est le modèle du texte N° 4 ci-dessus : « 2025-2026 : technicien stagiaire... 2024 : stage ouvrier... ». Le recruteur suit la progression du candidat comme une ligne de temps.

\textbf{Texte iconique N° 4 : le CV thématique (ou par compétences).} Les informations sont regroupées par \textbf{domaines de compétences} : « Conduite de chantier », « Dessin technique », « Gestion administrative », chaque domaine réunissant des expériences et des savoir-faire quelle que soit leur date. Ce modèle met en avant ce que le candidat \emph{sait faire} plutôt que son parcours.

\begin{center}
\begin{tabularx}{\linewidth}{|l|X|X|}
\hline
\rowcolor{popblueL} \textbf{Critère} & \textbf{CV chronologique} & \textbf{CV thématique} \\
\hline
Principe & Rubriques organisées par dates, de la plus récente à la plus ancienne & Informations regroupées par domaines de compétences \\
\hline
Avantages & Parcours clair et facile à vérifier ; montre la progression ; modèle attendu par la plupart des employeurs camerounais & Met en valeur les savoir-faire ; utile quand le parcours comporte des trous ou des changements de métier \\
\hline
Inconvénients & Révèle immédiatement les périodes sans emploi ; répétitif si les postes se ressemblent & Peut sembler cacher quelque chose ; le recruteur peine à situer les expériences dans le temps \\
\hline
À conseiller à... & Un jeune diplômé au parcours régulier (cas le plus fréquent en Terminale technique) & Un candidat expérimenté ayant exercé des métiers variés \\
\hline
\end{tabularx}
\end{center}

\courssub{Comparaison : l'effet de la mise en page sur le lecteur}

Comparons maintenant l'\emph{effet produit} par les deux textes iconiques. Dans le CV chronologique, la mise en valeur porte sur les \textbf{dates} : alignées en marge, elles créent une colonne visuelle que l'œil parcourt en premier ; le lecteur retient « ce candidat progresse régulièrement ». Dans le CV thématique, la mise en valeur porte sur les \textbf{titres de compétences} : en gras, dans des encadrés, ils captent le regard ; le lecteur retient « ce candidat sait faire ceci et cela ».

La photo, les puces, les encadrés et la couleur jouent aussi un rôle : une photo d'identité correcte personnalise la candidature ; les puces rendent les listes lisibles ; un encadré « Compétences » bien placé permet au recruteur pressé de vérifier en cinq secondes que le profil correspond à l'offre. À l'inverse, une page surchargée, des fautes d'orthographe dans les titres ou une photo inadaptée produisent une impression négative \emph{avant même} que le contenu soit lu.

\begin{attention}[La forme ne remplace pas le fond]
Un texte iconique réussi ne doit jamais sacrifier l'exactitude à l'esthétique. Un CV magnifique mais truffé de fautes, ou dont les dates ne correspondent pas aux diplômes, sera éliminé. Règle d'or : la mise en page sert le contenu, elle ne le maquille pas. Et toute information du CV doit pouvoir être prouvée par une pièce du dossier.
\end{attention}

\courssub{Bilan de la découverte}

Cette première étape de la leçon nous a permis de dégager l'idée directrice de toute l'unité : \textbf{le CV est à la fois un texte fonctionnel et un texte iconique}.

\begin{aretenir}[L'essentiel de la découverte]
\begin{itemize}
\item Un \cle{texte fonctionnel} est un écrit utilitaire qui répond à un besoin précis de la vie quotidienne (offre d'emploi, facture, formulaire, CV).
\item L'\textbf{offre d'emploi} est le texte de l'employeur ; le \textbf{CV} est le texte du candidat ; les deux se répondent point par point.
\item Lire une offre, c'est repérer trois zones : le \textbf{poste}, le \textbf{profil exigé}, les \textbf{conditions} (pièces, lieu, date limite).
\item Un \cle{texte iconique} porte une partie de son sens par la mise en page : colonnes, encadrés, puces, photo.
\item Le CV peut être \textbf{chronologique} (par dates) ou \textbf{thématique} (par compétences) ; le jeune diplômé choisit en général le modèle chronologique.
\item Le recruteur lit un CV en moins d'une minute : la clarté visuelle est aussi importante que la richesse du contenu.
\end{itemize}
\end{aretenir}

\begin{exercice}[À vous de jouer]
Voici une offre extraite d'un panneau d'affichage : « Atelier mécanique NJOH \& FILS, Douala-Bonabéri, recherche un tôlier-formeur. Profil : CAP ou Bac technique en chaudronnerie, 1 an d'expérience souhaité. Dossier : CV + photocopie du diplôme. Dépôt avant le 15 décembre 2026, 12 h. »
\begin{enumerate}
\item Souligne le poste, encadre le profil exigé, entoure les conditions pratiques.
\item Indique si cette offre est un texte fonctionnel ou iconique, et justifie.
\item Explique pourquoi le CV que tu rédigeras pour cette offre sera, lui, un texte à la fois fonctionnel et iconique.
\end{enumerate}
\end{exercice}

\begin{corrige}[Exercice : offre NJOH \& FILS]
1. Poste (à souligner) : « un tôlier-formeur ». Profil exigé (à encadrer) : « CAP ou Bac technique en chaudronnerie, 1 an d'expérience souhaité ». Conditions (à entourer) : « CV + photocopie du diplôme », « Douala-Bonabéri », « avant le 15 décembre 2026, 12 h ».
2. Cette offre est un texte fonctionnel : elle répond à un besoin précis (recruter), donne des informations exactes et vérifiables (diplôme, date, lieu) et utilise une forme codifiée (rubriques « Profil », « Dossier », « Dépôt »).
3. Le CV est fonctionnel parce qu'il vise une action concrète : obtenir le poste de tôlier-formeur ; il est iconique parce que sa mise en page (rubriques, puces, ordre des informations) guidera le regard du recruteur et mettra en valeur le diplôme de chaudronnerie et l'expérience exigés par l'offre.
\end{corrige}

La prochaine étape consistera à ouvrir le CV comme on ouvre une machine : nous en étudierons la \emph{structure externe}, c'est-à-dire l'organisation matérielle de la page.

\coursec{La structure externe du CV}

Dans la section précédente, nous avons lu et observé plusieurs textes fonctionnels et iconiques : des offres d'emploi, des annonces, des CV réels. Nous avons constaté que tous les CV, malgré leurs différences de présentation, se ressemblent : on y retrouve toujours les mêmes grandes parties, disposées à peu près dans le même ordre. Cette organisation visible, ce « découpage en blocs », c'est exactement ce que nous allons étudier maintenant : la \cle{structure externe} du CV. Savoir la reconnaître, c'est savoir où placer chaque information ; savoir la construire, c'est déjà faire la moitié du travail de rédaction.

\courssub{Définition : qu'est-ce que la structure externe ?}

\begin{definition}[Structure externe du CV]
La \cle{structure externe} d'un curriculum vitae est l'ensemble des \textbf{rubriques} (obligatoires ou facultatives) qui organisent le contenu du document : titre, état civil, formation, expérience professionnelle, compétences, centres d'intérêt, références. Elle concerne la \emph{forme} du document (l'ordre et la disposition des blocs d'information), par opposition à la \emph{structure interne}, qui concerne la manière de rédiger le contenu de chaque rubrique (choix des mots, formulation, style).
\end{definition}

Pour bien comprendre, faisons une analogie avec la construction d'une maison, familière aux élèves de Génie civil : la structure externe du CV, c'est le \emph{plan} de la maison — le nombre de pièces et leur disposition. La structure interne, ce sera la décoration et l'ameublement de chaque pièce. Un plan confus décourage le visiteur avant même qu'il n'entre ; de même, un CV mal structuré décourage le recruteur avant même qu'il ne lise.

\begin{attention}[Pourquoi la structure externe est-elle décisive ?]
Un employeur qui reçoit cent candidatures consacre en moyenne moins d'une minute à chaque CV lors du premier tri. Il ne lit pas : il \emph{repère}. Si l'information qu'il cherche (le diplôme, le téléphone) n'est pas à sa place habituelle, le CV est écarté. La structure externe est donc une \textbf{convention sociale} qu'il faut respecter pour être lu.
\end{attention}

\courssub{Observation : découper un CV en blocs}

\begin{experience}[Démarche d'investigation : du document observé au schéma]
\textbf{Matériel :} deux ou trois CV anonymisés (fournis par l'enseignant ou apportés par les élèves), un surligneur, une feuille blanche.

\textbf{Protocole :}
\begin{enumerate}
\item Entourer sur chaque CV les blocs d'information homogènes.
\item Nommer chaque bloc d'un seul mot ou d'une courte expression.
\item Comparer les découpages obtenus entre les différents CV et entre les groupes d'élèves.
\item Reporter sur la feuille blanche la maquette commune qui se dégage.
\end{enumerate}

\textbf{Observations :} quels que soient le métier visé et la mise en page, on retrouve presque toujours les mêmes blocs, dans un ordre voisin : identité en haut, diplômes et expérience au centre, compétences et loisirs en bas.

\textbf{Interprétation :} le CV obéit à un modèle social stabilisé. Ce modèle peut être schématisé par une maquette vierge, que tout candidat pourra ensuite remplir avec ses propres informations.
\end{experience}

\begin{popfigure}
\begin{center}
\begin{tikzpicture}[font=\small]
% Cadre de la page
\draw[thick,popdark] (0,0) rectangle (8.2,11.4);
% Titre
\fill[popblue!25] (0.3,10.5) rectangle (7.9,11.1);
\node at (4.1,10.8) {\textbf{1. TITRE : Curriculum Vitae (+ poste visé)}};
% Photo
\draw[thick,poppurple,fill=poppurple!15] (6.1,8.3) rectangle (7.9,10.3);
\node[align=center,popdark] at (7,9.3) {9. Photo\\d'identité};
% État civil
\fill[popgreen!20] (0.3,8.3) rectangle (5.9,10.3);
\node[align=left,anchor=north west] at (0.45,10.2) {\textbf{2. ÉTAT CIVIL}\\Nom, prénom\\Date et lieu de naissance\\Nationalité, situation matrimoniale\\Adresse, téléphone, e-mail};
% Formation
\fill[poporange!20] (0.3,6.2) rectangle (7.9,8.1);
\node[align=left,anchor=north west] at (0.45,8.0) {\textbf{3. FORMATION / DIPLÔMES}\\(du plus récent au plus ancien : année -- diplôme -- établissement)};
% Expérience
\fill[popblue!15] (0.3,4.1) rectangle (7.9,6.0);
\node[align=left,anchor=north west] at (0.45,5.9) {\textbf{4. EXPÉRIENCE PROFESSIONNELLE}\\(stages, emplois : dates -- poste -- lieu -- tâches)};
% Compétences
\fill[poppink!20] (0.3,2.4) rectangle (7.9,3.9);
\node[align=left,anchor=north west] at (0.45,3.8) {\textbf{5. COMPÉTENCES}\\(techniques : informatique, conduite\ldots\ / linguistiques : français, anglais\ldots)};
% Centres d'intérêt
\fill[popgold!25] (0.3,1.3) rectangle (7.9,2.2);
\node[align=left,anchor=north west] at (0.45,2.1) {\textbf{6. CENTRES D'INTÉRÊT}};
% Références + signature
\fill[popcream] (0.3,0.2) rectangle (5.9,1.1);
\node[align=left,anchor=north west] at (0.45,1.0) {\textbf{7. RÉFÉRENCES}};
\draw[thick,popmuted] (6.1,0.2) rectangle (7.9,1.1);
\node at (7,0.65) {10. Signature};
\end{tikzpicture}
\end{center}
\legende{Maquette d'un CV vierge découpé en blocs légendés : 1. titre ; 2. état civil ; 3. formation ; 4. expérience professionnelle ; 5. compétences ; 6. centres d'intérêt ; 7. références ; 9. photo ; 10. signature. L'ordre des blocs 3 et 4 peut être inversé selon le profil du candidat.}
\end{popfigure}

\courssub{Le titre du CV}

Tout CV commence par un \cle{titre}. Deux usages coexistent :

\begin{itemize}
\item la mention classique \emph{« Curriculum Vitae »} (expression latine signifiant « déroulement de la vie »), souvent suivie du nom du candidat ;
\item de plus en plus, l'intitulé du \textbf{poste visé} : « Curriculum Vitae — Candidat au poste de technicien en froid et climatisation ».
\end{itemize}

Mentionner le poste visé est une excellente habitude : elle montre que le CV a été pensé pour une embauche précise, et non photocopié pour être distribué au hasard. Nous y reviendrons dans la section sur l'adaptation du CV à une situation d'embauche.

\courssub{L'état civil : qui êtes-vous et comment vous joindre ?}

La rubrique \cle{état civil} (parfois intitulée « Informations personnelles ») se place en haut du document. Elle comprend traditionnellement :

\begin{center}
\begin{tabularx}{\linewidth}{|l|X|}\hline
\rowcolor{popblueL} \textbf{Élément} & \textbf{Précisions et conseils} \\ \hline
Nom et prénom(s) & Dans l'ordre officiel de l'acte de naissance ; le nom de famille peut être écrit en capitales : \emph{MBALLA Jean-Claude}. \\ \hline
Date et lieu de naissance & Ex. : « Né le 14 mars 2007 à Bafoussam ». \\ \hline
Nationalité & Ex. : « Camerounaise ». \\ \hline
Situation matrimoniale & Célibataire, marié(e), divorcé(e), veuf(ve) ; nombre d'enfants éventuel. \\ \hline
Adresse & Quartier, ville, boîte postale éventuelle : « Quartier Ndogpassi, Douala, BP 4521 ». \\ \hline
Téléphone & Numéro \emph{valide et joignable}, avec l'indicatif si besoin : (+237) 6 90 00 00 00. \\ \hline
E-mail & Adresse \emph{professionnelle et sérieuse}, consultée régulièrement. \\ \hline
\end{tabularx}
\end{center}

\begin{attention}[Erreurs fréquentes qui coûtent cher]
Deux fautes discréditent immédiatement un candidat :
\begin{itemize}
\item \textbf{oublier les coordonnées téléphoniques} ou donner un numéro qui ne fonctionne plus : l'employeur ne peut pas convoquer le candidat, le CV part à la corbeille ;
\item \textbf{donner une adresse e-mail fantaisiste} du type \emph{beaugoss237@...} ou \emph{laprincesse2k@...} : elle donne une image peu sérieuse. Créez une adresse sobre de la forme \emph{prenom.nom@...} et vérifiez-la chaque jour pendant votre recherche d'emploi.
\end{itemize}
\end{attention}

\begin{savaistu}[La situation matrimoniale au Cameroun]
Sur le plan international, la tendance est de supprimer du CV les informations personnelles non professionnelles (âge, situation matrimoniale, photo), afin d'éviter les discriminations. Au Cameroun cependant, beaucoup d'employeurs — administrations comme entreprises privées — \emph{exigent encore} la mention de la situation matrimoniale, parfois même du nombre d'enfants. Le candidat camerounais a donc intérêt à l'indiquer, sauf si l'offre d'emploi précise le contraire : c'est une information \textbf{facultative en droit, mais souvent attendue en pratique}.
\end{savaistu}

\courssub{La rubrique « Formation / Diplômes »}

Cette rubrique présente le parcours scolaire et les diplômes obtenus. Sa règle d'organisation est absolue :

\begin{methode}[Classer ses diplômes : l'ordre antéchronologique]
\begin{enumerate}
\item Classer les diplômes \textbf{du plus récent au plus ancien} (ordre \cle{antéchronologique}) : le recruteur veut d'abord savoir où vous en êtes aujourd'hui.
\item Pour chaque ligne, indiquer \textbf{l'année} (ou la période), \textbf{l'intitulé exact} du diplôme ou de la classe, et \textbf{l'établissement}.
\item Ne rien inventer : tout diplôme annoncé peut être vérifié ; un faux diplôme est un motif de licenciement, voire un délit.
\item Ne remonter en général pas plus bas que le BEPC ou le probatoire : le certificat d'études primaires n'a pas sa place dans le CV d'un bachelier.
\end{enumerate}
\end{methode}

\begin{exemplebox}[Rubrique « Formation » d'un élève de Terminale F4]
\textbf{FORMATION / DIPLÔMES}
\begin{itemize}
\item \textbf{2025-2026 :} classe de Terminale F4 (Génie civil), Lycée technique de Douala-Bassa (baccalauréat technique en préparation) ;
\item \textbf{2024 :} Probatoire technique, série F4 (Génie civil), Lycée technique de Douala-Bassa ;
\item \textbf{2022 :} BEPC, Collège bilingue de Deïdo.
\end{itemize}
Remarquez la présentation : l'année en gras d'abord, puis le diplôme, puis l'établissement. Trois lignes suffisent : clarté et honnêteté.
\end{exemplebox}

\courssub{La rubrique « Expérience professionnelle »}

Cette rubrique recense les \cle{stages}, les \cle{emplois} (même saisonniers ou informels) et les \cle{responsabilités} assumées. Elle suit, elle aussi, l'ordre antéchronologique. Chaque expérience doit préciser :

\begin{itemize}
\item les \textbf{dates} (mois et année de début et de fin) ;
\item le \textbf{poste occupé} et le \textbf{lieu} (entreprise, ville) ;
\item les \textbf{tâches principales} réalisées, formulées brièvement.
\end{itemize}

Un élève de Terminale technique n'a souvent « aucune expérience » : c'est faux. Le stage académique en entreprise, le petit commerce tenu pendant les vacances, l'aide à l'atelier d'un parent, la participation à un chantier de vacances sont des expériences \emph{valables} dès lors qu'on les décrit honnêtement.

\begin{exoresolu}[Rédiger une ligne d'expérience]
\textbf{Énoncé.} Aïcha, élève de Terminale Gestion, a effectué en juillet-août 2025 un stage de six semaines à la SODECOTON de Garoua, au service comptabilité, où elle saisissait les factures et classait les pièces comptables. Rédigez la ligne « Expérience professionnelle » correspondante.

\tcblower
\textbf{Solution détaillée.} On applique le schéma : dates -- poste -- lieu -- tâches.

\textbf{EXPÉRIENCE PROFESSIONNELLE}
\begin{itemize}
\item \textbf{Juillet-août 2025 :} stage académique au service comptabilité de la SODECOTON, Garoua — saisie des factures, classement des pièces comptables.
\end{itemize}

La formulation est brève, factuelle, sans faute ; les verbes d'action (« saisie », « classement ») donnent une image dynamique du candidat. On évite les phrases longues du type « pendant ce stage, j'ai eu l'occasion de\ldots » qui appartiennent à la lettre de motivation, pas au CV.
\end{exoresolu}

\courssub{La rubrique « Compétences »}

Les \cle{compétences} sont ce que le candidat \emph{sait faire}. On distingue deux grandes familles :

\begin{itemize}
\item les \textbf{compétences techniques}, liées à la spécialité : maîtrise de l'outil informatique (Word, Excel, logiciels de DAO comme AutoCAD), conduite d'engins (avec le permis correspondant), tenue d'une comptabilité, soudure à l'arc, installation électrique, secrétariat bureautique\ldots
\item les \textbf{compétences linguistiques} : le niveau en \textbf{français} et en \textbf{anglais} (lu, écrit, parlé), et les \textbf{langues nationales} pratiquées (ewondo, duala, fulfuldé, ghomala\ldots), précieuses dans un pays multilingue comme le Cameroun.
\end{itemize}

\begin{attention}[Ne pas surestimer ses compétences]
Écrire « parfaite maîtrise d'Excel » alors qu'on ne sait pas faire un tableau croisé, ou « anglais courant » alors qu'on ne peut pas tenir une conversation, est un piège : l'employeur vérifie souvent dès l'entretien. Préférez des mentions honnêtes et graduées : « notions », « niveau scolaire », « bon niveau », « courant ».
\end{attention}

\courssub{La rubrique « Centres d'intérêt »}

Les \cle{centres d'intérêt} (ou « Loisirs ») renseignent sur la personnalité du candidat. Ils ne sont pas obligatoires, mais ils peuvent faire la différence entre deux profils identiques. La règle : ne citer que des loisirs \textbf{significatifs et valorisants}, c'est-à-dire qui révèlent une qualité utile au poste.

\begin{center}
\begin{tabularx}{\linewidth}{|X|X|}\hline
\rowcolor{popgreenL} \textbf{Centres d'intérêt valorisants} & \textbf{Ce qu'ils suggèrent} \\ \hline
Football (capitaine de l'équipe du quartier) & Esprit d'équipe, sens des responsabilités \\ \hline
Lecture (romans africains) & Curiosité, maîtrise de la langue \\ \hline
Chorale paroissiale & Assiduité, travail collectif \\ \hline
Informatique (montage de vidéos) & Autonomie numérique \\ \hline
\end{tabularx}
\end{center}

À l'inverse, « regarder la télévision » ou « dormir » n'apportent rien : mieux vaut alors supprimer la rubrique.

\courssub{Les références}

Les \cle{références} sont des personnes (anciens employeurs, chefs de travaux, enseignants, notables) qui acceptent de \textbf{témoigner des qualités} du candidat auprès d'un recruteur. On indique leur nom, leur fonction et leurs coordonnées.

\begin{attention}[La règle d'or des références]
Il faut \textbf{toujours demander l'accord} de la personne avant de la citer comme référence. Un recruteur qui téléphone à quelqu'un de surpris — ou qui ne se souvient pas du candidat — obtient l'effet inverse de celui recherché. Si l'on ne dispose d'aucune référence sûre, on écrit simplement : « Références disponibles sur demande ».
\end{attention}

\courssub{La photo d'identité et la signature}

La \cle{photo d'identité} n'est pas obligatoire, mais elle reste très attendue au Cameroun. Si l'on en met une, elle doit être \textbf{récente}, de \textbf{format d'identité}, sur fond clair, avec une tenue et une expression \emph{professionnelles} ; elle se place en général \textbf{en haut à droite} du CV (voir la maquette). Une photo de fête recadrée ou un autoportrait au téléphone sont à proscrire. La \cle{signature}, enfin, n'est pas exigée, mais elle peut clore le document en bas à droite, éventuellement précédée de la mention « Fait à Douala, le\ldots » : elle donne au CV une valeur d'engagement personnel.

\begin{aretenir}[L'essentiel de la structure externe]
Un CV complet comporte, dans l'ordre usuel : le \textbf{titre}, l'\textbf{état civil} (avec des coordonnées valides et sérieuses), la \textbf{formation} et l'\textbf{expérience} (toutes deux du plus récent au plus ancien), les \textbf{compétences} (techniques et linguistiques), les \textbf{centres d'intérêt} valorisants, les \textbf{références} autorisées, et éventuellement la \textbf{photo} et la \textbf{signature}. Cette architecture fixe est la convention qui permet au recruteur de trouver chaque information en quelques secondes.
\end{aretenir}

\begin{exercice}[À vous de jouer]
\begin{enumerate}
\item Voici les informations de Rodrigue, élève de Terminale F3 : BEPC en 2022 ; classe de Terminale F3 au Lycée technique de Yaoundé en 2025-2026 ; Probatoire F3 en 2024. Rédigez sa rubrique « Formation » dans le bon ordre.
\item Relevez trois erreurs dans cet état civil : « Nom : Petit Prince ; e-mail : championdu237@... ; téléphone : aucun pour le moment ; situation : compliquée ».
\item Rodrigue a aidé son oncle, maçon, sur un chantier à Obala pendant les grandes vacances 2025. Rédigez la ligne « Expérience professionnelle » correspondante.
\end{enumerate}
\end{exercice}

\begin{corrige}[Exercice : structure externe]
\begin{enumerate}
\item \textbf{Formation :} 2025-2026 : classe de Terminale F3 (Génie mécanique), Lycée technique de Yaoundé ; 2024 : Probatoire technique F3, Lycée technique de Yaoundé ; 2022 : BEPC. (Ordre antéchronologique respecté : du plus récent au plus ancien.)
\item Erreurs : « Petit Prince » est un surnom, pas un nom officiel ; l'adresse e-mail est fantaisiste et doit être remplacée par une adresse sobre ; l'absence de téléphone empêche tout contact ; « compliquée » n'est pas une situation matrimoniale recevable (on écrit : célibataire, marié, etc.).
\item \textbf{Expérience professionnelle :} « Juillet-septembre 2025 : aide-maçon sur un chantier de construction, Obala — préparation du mortier, approvisionnement des matériaux. »
\end{enumerate}
\end{corrige}

Maintenant que l'architecture du document est claire — nous savons \emph{quelles} rubriques composent un CV et \emph{dans quel ordre} —, il reste à apprendre \emph{comment remplir} chacune d'elles : quels mots choisir, quelles tournures employer, quel style adopter. Ce sera l'objet de la section suivante, consacrée à la structure interne du CV.

\coursec{La structure interne du CV}

Dans la section précédente, nous avons étudié la \cle{structure externe} du CV : les rubriques qui le composent (état civil, formation, expérience professionnelle, compétences, centres d'intérêt) et leur disposition sur la page. Mais deux CV qui présentent exactement les mêmes rubriques peuvent produire des effets très différents sur un recruteur. Pourquoi ? Parce que ce qui se passe \emph{à l'intérieur} des rubriques — l'ordre des informations, la manière de rédiger, la correction de la langue — est tout aussi décisif. C'est ce que nous appelons la \cle{structure interne} du CV, objet de cette section.

\courssub{Définition de la structure interne}

Commençons par une intuition. Imaginez deux élèves qui racontent le même stage en entreprise. Le premier dit : « J'ai fait un stage, c'était bien, on m'a montré plein de choses, j'ai travaillé avec les ouvriers et aussi un peu au bureau, enfin voilà. » Le second dit : « Stage d'un mois aux Établissements Nkoulou : fabrication de portes et fenêtres, lecture de plans, accueil de la clientèle. » Le contenu est le même, mais la seconde formulation est claire, ordonnée et crédible. La différence ne tient pas aux faits, mais à leur \emph{organisation} et à leur \emph{formulation} : c'est précisément la structure interne.

\begin{definition}[Structure interne du CV]
La \cle{structure interne} du CV désigne l'organisation logique et la qualité rédactionnelle du contenu \emph{à l'intérieur} de chaque rubrique. Elle repose sur quatre exigences :
\begin{itemize}
\item l'\cle{ordre} : les informations sont classées de manière logique (ordre antéchronologique, hiérarchisation) ;
\item la \cle{cohérence} : les rubriques ne se contredisent pas et le parcours ne présente pas de « trou » injustifié ;
\item la \cle{concision} : le CV tient sur une page, en phrases courtes ou en énumérations ;
\item la \cle{correction linguistique} : orthographe irréprochable, vocabulaire précis, mise en page soignée.
\end{itemize}
C'est la structure interne qui donne sa \emph{qualité} au contenu du CV.
\end{definition}

\begin{aretenir}[L'essentiel]
Structure externe = les rubriques et leur disposition (le « squelette »). Structure interne = l'ordre, la cohérence, la concision et la correction du contenu (la « chair »). Un bon CV exige les deux.
\end{aretenir}

\courssub{La concision : une page, pas plus}

Un recruteur qui reçoit deux cents candidatures consacre en moyenne moins d'une minute à chaque CV lors du premier tri. Un document long, dense et bavard est donc presque toujours écarté. La règle est simple :

\begin{propriete}[Règle de concision]
Un CV tient sur \textbf{une page} (deux au maximum pour un cadre très expérimenté). On rédige en \textbf{phrases courtes} ou, mieux, en \textbf{simples énumérations} introduites par des puces. On supprime tout ce qui n'apporte aucune information utile à l'employeur.
\end{propriete}

Comparer deux formulations d'une même expérience permet de mesurer l'efficacité de la concision :

\begin{exemplebox}[Avant et après]
\textbf{Version bavarde (à éviter) :} « Pendant les grandes vacances de l'année 2024, j'ai eu l'occasion d'effectuer un stage dans une menuiserie qui se trouve à Yaoundé et qui appartient aux Établissements Nkoulou, et là-bas j'ai participé à la fabrication de différentes choses et j'ai aussi aidé à l'accueil des clients qui venaient. »

\textbf{Version concise (correcte) :} « 2024 : stagiaire, menuiserie Éts Nkoulou, Yaoundé : fabriquer des portes et fenêtres ; lire des plans ; accueillir la clientèle. »
\end{exemplebox}

La seconde version dit exactement la même chose en trois fois moins de mots, et elle est beaucoup plus lisible.

\courssub{L'ordre logique : l'antéchronologique et la hiérarchisation}

\courssub{L'ordre antéchronologique}

Dans les rubriques \emph{Formation} et \emph{Expérience professionnelle}, les informations se présentent dans l'\cle{ordre antéchronologique} : on commence par la situation la \textbf{plus récente} et on remonte vers la plus ancienne. Pourquoi ? Parce que ce qui intéresse d'abord l'employeur, c'est ce que vous savez faire \emph{aujourd'hui}, pas ce que vous faisiez il y a dix ans.

\begin{exemplebox}[Rubrique « Formation » bien ordonnée]
\begin{itemize}
\item 2025--2026 : classe de Terminale F4 (Génie civil), Lycée technique de Douala — Baccalauréat technique en préparation.
\item 2024 : Probatoire technique F4, Lycée technique de Douala.
\item 2022 : BEPC, Collège de la Cité Verte, Yaoundé.
\end{itemize}
\end{exemplebox}

\begin{attention}[Piège fréquent]
Ne mélangez pas les ordres : si vous choisissez l'antéchronologique (recommandé), appliquez-le à \emph{toutes} les rubriques concernées. Un CV qui commence la formation par le diplôme le plus récent mais l'expérience par le stage le plus ancien donne une impression de désordre.
\end{attention}

\courssub{La hiérarchisation des compétences}

Dans la rubrique \emph{Compétences}, on ne liste pas tout au hasard : on place \textbf{en premier} les compétences les plus utiles pour le poste visé. Un candidat à un poste de dessinateur en bâtiment mettra « maîtrise d'AutoCAD » en tête ; un candidat à un poste de secrétaire mettra « frappe rapide, maîtrise de Word et Excel » en tête. La hiérarchisation est donc déjà une forme d'adaptation du CV à la situation d'embauche, que nous approfondirons dans la section 6.

\courssub{La cohérence : un parcours sans contradiction}

La \cle{cohérence} interne impose deux contrôles :

\begin{enumerate}
\item \textbf{Pas de contradiction entre les rubriques.} Si l'état civil indique une naissance en 2006, la formation ne peut pas mentionner un BEPC obtenu en 2015. Si l'expérience mentionne « technicien de maintenance, 2023--2024 », la formation doit montrer que le candidat était disponible à cette période.
\item \textbf{Pas de « trou » injustifié.} Une année vide entre deux diplômes ou deux emplois intrigue le recruteur. Il faut l'expliquer honnêtement : « 2023--2024 : année de préparation du concours d'entrée à l'ENSET » ou « 2023 : emploi saisonnier, agriculture familiale, Mbanga ». Une période expliquée rassure ; une période vide inquiète.
\end{enumerate}

\courssub{L'honnêteté : ne jamais rien inventer}

Il peut être tentant d'embellir son CV : ajouter un diplôme « en cours », gonfler une durée de stage, s'attribuer un poste jamais occupé. C'est une faute grave, pour trois raisons :

\begin{itemize}
\item les employeurs \textbf{vérifient systématiquement} : appel aux anciens employeurs, demande d'attestations, tests pratiques à l'embauche ;
\item un mensonge découvert \emph{avant} l'embauche élimine définitivement le candidat ;
\item un mensonge découvert \emph{après} l'embauche a des conséquences juridiques.
\end{itemize}

\begin{savaistu}[Le mensonge sur le CV et la loi]
Au Cameroun, comme dans la plupart des pays, un mensonge sur un diplôme ou une expérience découvert après l'embauche constitue une \emph{faute lourde} : c'est un motif \textbf{légal de licenciement} sans préavis ni indemnité. Mieux vaut un CV modeste et vrai qu'un CV brillant et faux.
\end{savaistu}

\begin{attention}[Erreur éliminatoire]
Les fautes d'orthographe dans un CV sont \textbf{éliminatoires} : pour un recruteur, un candidat qui ne relit pas son propre CV ne relira pas non plus un devis, un rapport ou une facture. Faites \textbf{toujours relire votre CV par deux personnes} différentes avant de l'envoyer.
\end{attention}

\courssub{La mise en page et la langue}

\courssub{La mise en page}

La mise en page matérialise visuellement la structure interne. Quatre règles suffisent :

\begin{itemize}
\item \textbf{alignements} réguliers : les dates alignées sur la même colonne, les puces au même retrait ;
\item \textbf{puces} identiques dans tout le document (toutes rondes ou toutes carrées) ;
\item \textbf{caractères gras} réservés aux titres de rubriques et éventuellement aux intitulés de postes ;
\item \textbf{marges régulières} et une seule police de caractères, sobre (Arial, Calibri, Times).
\end{itemize}

\courssub{La langue}

La langue du CV doit être \textbf{irréprochable} : orthographe et accords vérifiés, vocabulaire précis (« assurer la maintenance » plutôt que « faire des trucs de réparation »), et surtout \textbf{aucune abréviation SMS} (« j'ai », et non « g » ; « entreprise », et non « entr »). Le CV est un document professionnel : il obéit aux normes du français écrit formel.

\courssub{La reformulation des tâches : les infinitifs d'action}

C'est le point le plus technique de la structure interne. Pour décrire ce que vous avez fait dans un poste ou un stage, la convention du CV impose d'utiliser des \textbf{verbes d'action à l'infinitif}, sans sujet exprimé. L'infinitif, nous le verrons dans la section 4, a ici une valeur de \emph{consigne générale} et de \emph{description objective} : il présente la tâche comme une fonction, pas comme un récit personnel.

\begin{methode}[Décrire une expérience professionnelle en quatre éléments]
Pour chaque expérience, rédiger une entrée complète en suivant le schéma :
\begin{enumerate}
\item \textbf{Date} (année ou période) ;
\item \textbf{Poste} occupé (stagiaire, aide-comptable, technicien...) ;
\item \textbf{Structure} (nom de l'entreprise et ville) ;
\item \textbf{Deux ou trois tâches} exprimées par des \textbf{infinitifs d'action} : réaliser, assurer, fabriquer, rédiger, accueillir, contrôler, saisir, installer...
\end{enumerate}
Schéma : \emph{Date : poste, structure, ville — infinitif + complément ; infinitif + complément.}
\end{methode}

\begin{exemplebox}[Application de la méthode]
\textbf{2024 : stagiaire, menuiserie Éts Nkoulou, Yaoundé :} fabriquer des portes et fenêtres ; lire des plans ; accueillir la clientèle.

\medskip
On retrouve les quatre éléments : date (2024), poste (stagiaire), structure (menuiserie Éts Nkoulou, Yaoundé), trois tâches à l'infinitif d'action.
\end{exemplebox}

Le tableau comparatif ci-dessous rassemble les principales corrections de structure interne étudiées dans cette section.

\begin{popfigure}
\begin{tikzpicture}[font=\small]
% En-têtes
\fill[popink!85] (-7,0) rectangle (0,0.9);
\fill[popgreen!70!black] (0,0) rectangle (7,0.9);
\node[white,font=\bfseries] at (-3.5,0.45) {CV mal rédigé};
\node[white,font=\bfseries] at (3.5,0.45) {CV corrigé};
% Ligne 1 : concision
\fill[popcream] (-7,-2.6) rectangle (0,0);
\fill[popgreenL] (0,-2.6) rectangle (7,0);
\node[text width=6.5cm,anchor=north west,align=left] at (-6.8,-0.15)
{« Pendant les vacances de 2024, j'ai eu l'occasion de faire un stage dans une menuiserie à Yaoundé où j'ai participé à plein de choses... »};
\node[text width=6.5cm,anchor=north west,align=left] at (0.2,-0.15)
{« 2024 : stagiaire, menuiserie Éts Nkoulou, Yaoundé : fabriquer des portes et fenêtres ; lire des plans. »};
% Ligne 2 : ordre
\fill[popcream!60] (-7,-4.6) rectangle (0,-2.6);
\fill[popgreenL!60] (0,-4.6) rectangle (7,-2.6);
\node[text width=6.5cm,anchor=north west,align=left] at (-6.8,-2.75)
{Formation : 2022 : BEPC ; 2024 : Probatoire ; 2026 : Bac en préparation (ordre chronologique).};
\node[text width=6.5cm,anchor=north west,align=left] at (0.2,-2.75)
{Formation : 2026 : Bac F4 en préparation ; 2024 : Probatoire ; 2022 : BEPC (ordre antéchronologique).};
% Ligne 3 : langue
\fill[popcream] (-7,-6.6) rectangle (0,-4.6);
\fill[popgreenL] (0,-6.6) rectangle (7,-4.6);
\node[text width=6.5cm,anchor=north west,align=left] at (-6.8,-4.75)
{« g travaillé ds 1 entr de btp, je fesais les devis et la compta » (SMS, fautes, imprécision).};
\node[text width=6.5cm,anchor=north west,align=left] at (0.2,-4.75)
{« 2023 : aide-comptable, Éts Mbappé Bâtiment, Douala : rédiger des devis ; saisir les factures. »};
% Contours
\draw[popline,thick] (-7,0.9) rectangle (0,-6.6);
\draw[popline,thick] (0,0.9) rectangle (7,-6.6);
\end{tikzpicture}
\legende{Tableau comparatif : trois extraits fautifs (colonne de gauche) et leurs versions corrigées (colonne de droite) — concision, ordre antéchronologique, langue soignée et infinitifs d'action.}
\end{popfigure}

\courssub{Exercice résolu : corriger une rubrique fautive}

\begin{exoresolu}[Corriger une rubrique « Expérience professionnelle »]
\textbf{Énoncé.} Voici la rubrique « Expérience professionnelle » du CV d'Amina, élève de Terminale F3 (Comptabilité) : « Expérience : en 2022 j'ai aidé ma tante dans sa boutique à Bafoussam, ensuite en 2025 j'ai fait un stage à la SABC de Douala où j'ai fait la saisie des factures et j'ai aussi rangé les archives, c'était très intéressant et j'ai beaucoup appris. » Réécrivez cette rubrique en respectant la structure interne du CV.

\tcblower
\textbf{Solution détaillée.} On applique la méthode en quatre éléments (date + poste + structure + tâches à l'infinitif) et l'ordre antéchronologique (2025 avant 2022). On supprime les commentaires personnels (« c'était très intéressant »), qui n'ont pas leur place dans un CV.

\medskip
\textbf{Expérience professionnelle}
\begin{itemize}
\item 2025 : stagiaire comptable, SABC, Douala : saisir les factures ; classer les archives comptables.
\item 2022 : aide-vendeuse (emploi de vacances), boutique familiale, Bafoussam : accueillir les clients ; tenir la caisse.
\end{itemize}

\medskip
\textbf{Justification de la démarche :} l'expérience la plus récente est placée en premier ; chaque entrée suit le schéma date + poste + structure + infinitifs d'action ; les tâches sont décrites par des verbes précis (saisir, classer, accueillir, tenir) ; les jugements subjectifs ont été supprimés pour gagner en concision et en objectivité.
\end{exoresolu}

\courssub{Bilan de la section}

\begin{aretenir}[Les sept règles d'or de la structure interne]
\begin{enumerate}
\item Une page maximum, phrases courtes ou énumérations (\textbf{concision}).
\item Ordre \textbf{antéchronologique} pour la formation et l'expérience.
\item \textbf{Hiérarchisation} des compétences selon le poste visé.
\item \textbf{Cohérence} : aucune contradiction, aucun trou injustifié.
\item \textbf{Honnêteté} : ne jamais inventer diplôme ni expérience.
\item \textbf{Mise en page} soignée : alignements, puces, gras pour les rubriques.
\item \textbf{Langue} irréprochable et tâches décrites par des \textbf{infinitifs d'action}.
\end{enumerate}
\end{aretenir}

La maîtrise de la structure interne repose en grande partie sur l'emploi correct de l'infinitif et, dans les documents qui accompagnent le CV (consignes, offres d'emploi), de l'impératif. C'est précisément l'objet de la section suivante, consacrée aux valeurs de ces deux modes.

\coursec{Langue : les valeurs de l'impératif et de l'infinitif}

Après avoir analysé la \textbf{structure interne du CV} (rubriques, chronologie, verbes d'action), nous devons maintenant maîtriser les \textbf{outils grammaticaux} qui donnent à ce document son ton professionnel. Le CV et la lettre de motivation relèvent de la communication administrative et professionnelle : le choix entre l'impératif et l'infinitif n'est pas anodin, il détermine la relation entre l'émetteur (le recruteur ou le candidat) et le destinataire.

\courssub{Rappel : les modes personnels et impersonnels}

\begin{definition}[Mode verbal]
Le \cle{mode} est la forme du verbe qui exprime l'attitude du locuteur vis-à-vis de l'action ou de l'état : affirmation (indicatif), ordre (impératif), souhait/hypothèse (subjonctif), ou simple désignation de l'action sans personne ni temps (infinitif, participe).
\end{definition}

En français, on distingue deux grandes catégories :
\begin{itemize}
    \item \textbf{Les modes personnels} (indicatif, subjonctif, conditionnel, impératif) : le verbe est conjugué, il s'accorde avec un sujet (je, tu, il, nous, vous, ils).
    \item \textbf{Les modes impersonnels} (infinitif, participe présent, participe passé) : le verbe n'a \emph{pas de sujet}, il ne marque ni personne, ni nombre, ni temps (seul l'infinitif passé et le participe passé marquent l'antériorité).
\end{itemize}

\begin{aretenir}[Distinction fondamentale pour le CV]
\begin{itemize}
    \item \textbf{Impératif} = mode personnel (2e personne du singulier ou du pluriel, 1re personne du pluriel). Il \textbf{s'adresse à quelqu'un} : « \emph{Envoyez} votre dossier. »
    \item \textbf{Infinitif} = mode impersonnel. Il \textbf{désigne l'action en soi}, sans destinataire précis : « \emph{Envoyer} un dossier complet. »
\end{itemize}
\end{aretenir}

\courssub{L'impératif : formation et valeurs}

\textbf{Formation de l'impératif présent}\par
L'impératif présent se forme sur la base du \textbf{présent de l'indicatif} aux 2e personnes (singulier et pluriel) et à la 1re personne du pluriel, \textbf{sans pronom sujet}.

\begin{exemplebox}[Règles de formation]
\begin{center}
\begin{tabularx}{\linewidth}{|l|X|X|X|}\hline
\textbf{Groupe} & \textbf{2e pers. sg. (tu)} & \textbf{1re pers. pl. (nous)} & \textbf{2e pers. pl. (vous)} \\ \hline
1er groupe (chanter) & Chante & Chantons & Chantez \\ \hline
2e groupe (finir) & Finis & Finissons & Finissez \\ \hline
3e groupe (partir, prendre, être, avoir) & Pars / Prends / Sois / Aie & Partons / Prenons / Soyons / Ayons & Partez / Prenez / Soyez / Ayez \\ \hline
\end{tabularx}
\end{center}
\legende{Rappel : le -s du 2e sg. tombe pour les verbes du 1er groupe et \emph{aller}, \emph{ouvrir}, \emph{offrir}, \emph{cueillir} (sauf devant \emph{en/y} : \emph{Vas-y}, \emph{Envoies-en}).}
\end{exemplebox}

\textbf{Les cinq valeurs de l'impératif dans les textes fonctionnels}\par
L'impératif n'exprime pas seulement l'ordre militaire. Dans une offre d'emploi ou une consigne administrative, il prend des nuances essentielles à identifier.

\begin{popfigure}
\begin{center}
\begin{tikzpicture}[
    cell/.style={rectangle, draw=popline, minimum height=1.1cm, align=center, text width=3.2cm, font=\small},
    head/.style={fill=popblue, text=white, font=\bfseries\small, minimum height=0.9cm},
    val/.style={fill=popblueL!50, font=\itshape\small},
    ex/.style={fill=popcream, font=\small},
    row/.style={minimum height=1.3cm}
]
% En-têtes
\node[head, text width=2.5cm] (A1) at (0,0) {Valeur};
\node[head, text width=3.2cm] (B1) at (2.5,0) {Définition / Intention};
\node[head, text width=3.2cm] (C1) at (5.7,0) {Marqueurs contextuels};
\node[head, text width=3.5cm] (D1) at (8.9,0) {Exemple dans une offre d'emploi / CV};

% Ligne 1 : Ordre
\node[val, row] (A2) at (0,-1.1) {Ordre / Injiction};
\node[ex, row] (B2) at (2.5,-1.1) {Contrainte forte, obligation, autorité hiérarchique};
\node[ex, row] (C2) at (5.7,-1.1) {Verbes : \emph{Envoyez, Déposez, Présentez} ; Délai impératif (\emph{avant le...})};
\node[ex, row] (D2) at (8.9,-1.1) {\textbf{« \emph{Envoyez} votre dossier \emph{avant le 30 juin} sous pli recommandé. »}};

% Ligne 2 : Conseil
\node[val, row] (A3) at (0,-2.4) {Conseil / Recommandation};
\node[ex, row] (B3) at (2.5,-2.4) {Bienveillance, incitation à agir pour son bien};
\node[ex, row] (C3) at (5.7,-2.4) {Verbes : \emph{Privilégiez, Préparez, Soignez, Évitez} ; Adverbes : \emph{vivement, fortement}};
\node[ex, row] (D3) at (8.9,-2.4) {\textbf{« \emph{Soignez} la présentation de votre CV : aéré, sans faute. »}};

% Ligne 3 : Défense / Interdiction
\node[val, row] (A4) at (0,-3.7) {Défense / Interdiction};
\node[ex, row] (B4) at (2.5,-3.7) {Interdiction formelle, sanction possible (négation \emph{ne...pas})};
\node[ex, row] (C4) at (5.7,-3.7) {\emph{Ne... pas / Ne... jamais} + verbe ; \emph{Évitez de...}};
\node[ex, row] (D4) at (8.9,-3.7) {\textbf{« \emph{Ne joignez pas} de photo d'identité si non demandé. »}};

% Ligne 4 : Prière / Demande polie
\node[val, row] (A5) at (0,-5.0) {Prière / Demande polie};
\node[ex, row] (B5) at (2.5,-5.0) {Atténuation de l'ordre, politesse, formule figée};
\node[ex, row] (C5) at (5.7,-5.0) {\emph{Prière de, Veuillez, Merci de, Ayez l'amabilité de} + infinitif};
\node[ex, row] (D5) at (8.9,-5.0) {\textbf{« \emph{Veuillez} adresser votre candidature par email uniquement. »}};

% Ligne 5 : Consigne / Mode d'emploi
\node[val, row] (A6) at (0,-6.3) {Consigne / Mode d'emploi};
\node[ex, row] (B6) at (2.5,-6.3) {Procédure à suivre, étapes successives, ton neutre mais direct};
\node[ex, row] (C6) at (5.7,-6.3) {Verbes d'action : \emph{Remplissez, Cochez, Signez, Datez, Retournez}};
\node[ex, row] (D6) at (8.9,-6.3) {\textbf{« \emph{Remplissez} le formulaire ci-joint, \emph{signez}-le et \emph{retournez}-le avant le 15/07. »}};

\draw[popline, thick] (0,0.45) -- (12.4,0.45);
\draw[popline, thick] (0,-7.4) -- (12.4,-7.4);
\draw[popline, thick] (0,0.45) -- (0,-7.4);
\draw[popline, thick] (2.5,0.45) -- (2.5,-7.4);
\draw[popline, thick] (5.7,0.45) -- (5.7,-7.4);
\draw[popline, thick] (8.9,0.45) -- (8.9,-7.4);
\draw[popline, thick] (12.4,0.45) -- (12.4,-7.4);
\end{tikzpicture}
\end{center}
\legende{Tableau récapitulatif des valeurs de l'impératif avec exemples contextuels (offre d'emploi / consignes).}
\end{popfigure}

\begin{exoresolu}[Identifier la valeur de l'impératif]
\textbf{Énoncé :} Dans l'extrait d'offre ci-dessous, relevez les verbes à l'impératif et précisez leur valeur.
\begin{quote}
\textit{« Candidatez en ligne sur notre site. \textbf{Privilégiez} le format PDF pour votre CV. \textbf{Ne dépassez pas} 2 Mo par fichier. \textbf{Veuillez} indiquer la référence \#REF-2024 dans l'objet. »}
\end{quote}
\tcblower
\textbf{Solution détaillée :}
\begin{enumerate}
    \item \textbf{Candidatez} (2e pl. 1er gr.) $\rightarrow$ \textbf{Consigne} : procédure à suivre (étape 1).
    \item \textbf{Privilégiez} (2e pl. 2e gr.) $\rightarrow$ \textbf{Conseil} : recommandation technique pour le bon traitement du dossier (pas d'obligation stricte, mais forte incitation).
    \item \textbf{Ne dépassez pas} (2e pl. négatif) $\rightarrow$ \textbf{Défense / Interdiction} : contrainte technique stricte (rejet automatique si non respecté).
    \item \textbf{Veuillez indiquer} (Prière de + infinitif) $\rightarrow$ \textbf{Prière / Demande polie} : formule de politesse administrative standard pour une contrainte réelle (obligation de mettre la référence).
\end{enumerate}
\end{exoresolu}

\courssub{L'infinitif : forme, emplois et valeurs}

\textbf{Rappel : l'infinitif, mode impersonnel}\par

\begin{definition}[Infinitif]
L'\cle{infinitif} est la forme nominale du verbe. Il n'a ni personne, ni nombre, ni temps (sauf l'infinitif passé : \emph{avoir fait}). Il fonctionne comme un \textbf{nom verbal} : il peut être sujet, objet, complément du nom, ou former une proposition infinitif (sans sujet explicite).
\end{definition}

Dans les textes fonctionnels (offres d'emploi, formulaires, CV), l'infinitif est roi pour deux raisons :
\begin{enumerate}
    \item \textbf{Impersonnalité :} il s'adresse à \emph{tout le monde} sans viser un « vous » précis (neutralité administrative).
    \item \textbf{Concision :} il permet des listes, des énumérations de tâches ou de pièces à fournir sans répéter le sujet.
\end{enumerate}

\textbf{Les trois valeurs majeures de l'infinitif dans le CV et l'offre d'emploi}\par

\begin{popfigure}
\begin{center}
\begin{tikzpicture}[
    cell/.style={rectangle, draw=popline, minimum height=1.1cm, align=center, text width=3.5cm, font=\small},
    head/.style={fill=poppurple, text=white, font=\bfseries\small, minimum height=0.9cm},
    val/.style={fill=poppurpleL!50, font=\itshape\small},
    ex/.style={fill=popcream, font=\small},
    row/.style={minimum height=1.5cm}
]
% En-têtes
\node[head, text width=3cm] (A1) at (0,0) {Valeur};
\node[head, text width=4cm] (B1) at (3,0) {Fonction syntaxique / Sémantique};
\node[head, text width=4.5cm] (C1) at (7,0) {Exemple dans une offre / CV};

% Ligne 1 : Consigne impersonnelle
\node[val, row] (A2) at (0,-1.1) {Consigne impersonnelle (Règle / Procédure)};
\node[ex, row] (B2) at (3,-1.1) {Remplace l'impératif pour une portée générale. Sujet indéfini : « Il faut... », « Tout candidat doit... »};
\node[ex, row] (C2) at (7,-1.1) {\textbf{Offre :} « \emph{Fournir} un CV + lettre de motivation. \emph{Joindre} les copies des diplômes. »};

% Ligne 2 : Énumération de tâches / Compétences (CV)
\node[val, row] (A3) at (0,-2.6) {Énumération de tâches / Réalisations (CV)};
\node[ex, row] (B3) at (3,-2.6) {Style télégraphique, tête de liste. Verbes d'action. Remplace « J'ai fait / Je fais ».};
\node[ex, row] (C3) at (7,-2.6) {\textbf{CV :} « \emph{Gérer} une équipe de 10 personnes. \emph{Concevoir} des tableaux de bord. \emph{Assurer} le reporting mensuel. »};

% Ligne 3 : Objectif / Finalité (Pour + infinitif)
\node[val, row] (A4) at (0,-4.1) {Objectif / Finalité (Préposition + Infinitif)};
\node[ex, row] (B4) at (3,-4.1) {Introduit par \emph{pour, afin de, en vue de, dans le but de}. Exprime la finalité de l'action principale.};
\node[ex, row] (C4) at (7,-4.1) {\textbf{Lettre / Offre :} « \emph{Pour postuler}, \emph{adresser} votre dossier à l'adresse ci-dessous. »};

% Ligne 4 : Infinitif passé (Antériorité)
\node[val, row] (A5) at (0,-5.6) {Infinitif passé (Antériorité)};
\node[ex, row] (B5) at (3,-5.6) {\emph{Avoir / Être} + Participe passé. Action antérieure à celle du verbe principal ou au moment de l'énonciation.};
\node[ex, row] (C5) at (7,-5.6) {\textbf{CV :} « \emph{Avoir dirigé} un projet international est un atout. » (Condition d'expérience passée).};

\draw[popline, thick] (0,0.45) -- (11.5,0.45);
\draw[popline, thick] (0,-6.7) -- (11.5,-6.7);
\draw[popline, thick] (0,0.45) -- (0,-6.7);
\draw[popline, thick] (3,0.45) -- (3,-6.7);
\draw[popline, thick] (7,0.45) -- (7,-6.7);
\draw[popline, thick] (11.5,0.45) -- (11.5,-6.7);
\end{tikzpicture}
\end{center}
\legende{Tableau récapitulatif des valeurs de l'infinitif dans les textes professionnels.}
\end{popfigure}

\textbf{L'infinitif dans la rubrique « Expériences professionnelles » du CV}\par
C'est l'application la plus critique pour l'élève de Terminale Technique. Le CV moderne (style « Anglo-saxon » ou « Compétences ») exige l'\textbf{infinitif de narration} (ou infinitif de liste) pour décrire les missions.

\begin{propriete}[Règle d'or du CV : Verbes d'action à l'infinitif]
Dans la rubrique \textbf{Expériences professionnelles}, chaque mission commence par un \textbf{verbe d'action à l'infinitif présent} (ou passé pour une mission terminée et ponctuelle).
\begin{itemize}
    \item \textbf{Infinitif présent :} action durable, habituelle, compétence acquise. Ex: \emph{Gérer, Animer, Développer, Maintenance, Superviser}.
    \item \textbf{Infinitif passé :} réalisation ponctuelle, projet achevé, résultat chiffré. Ex: \emph{Avoir mis en place}, \emph{Avoir réduit} les coûts de 15\%.
\end{itemize}
\end{propriete}

\begin{exemplebox}[Transformation « Je fais » $\rightarrow$ « Faire » (Style CV)]
\begin{center}
\begin{tabularx}{\linewidth}{|l|X|X|}\hline
\textbf{Phrase parlée / Lettre (1re pers.)} & \textbf{Style CV (Infinitif)} & \textbf{Verbe d'action choisi} \\ \hline
J'ai géré une équipe de 5 techniciens. & \textbf{Gérer} une équipe de 5 techniciens. & \emph{Gérer} (management) \\ \hline
Je suis responsable de la maintenance. & \textbf{Assurer} la maintenance préventive et curative. & \emph{Assurer} (responsabilité) \\ \hline
J'ai créé un nouveau fichier clients. & \textbf{Concevoir} et \textbf{déployer} un fichier clients. & \emph{Concevoir / Déployer} (création) \\ \hline
Je vends des produits aux clients. & \textbf{Prospecter}, \textbf{négocier} et \textbf{fidéliser} la clientèle. & \emph{Prospecter / Négocier / Fidéliser} (processus commercial) \\ \hline
\end{tabularx}
\end{center}
\legende{L'infinitif gomme le « Je » pour mettre en valeur l'\textbf{action} et la \textbf{compétence}.}
\end{exemplebox}

\courssub{Choix entre Impératif et Infinitif : le paramètre « Destinataire »}

C'est le cœur de la compétence « Adapter sa langue à la situation de communication ».

\begin{propriete}[Impératif vs Infinitif : Le ton de la relation]
\begin{itemize}
    \item \textbf{Impératif (Ton direct, personnel)} $\rightarrow$ \textbf{L'émettrice s'adresse AU destinataire.}
    \begin{itemize}
        \item Recruteur $\rightarrow$ Candidat : « \emph{Envoyez} votre CV. » (Ordre/Consigne directe).
        \item Candidat $\rightarrow$ Recruteur (Lettre) : « \emph{Permettez-moi} de vous présenter... » (Prière).
    \end{itemize}
    \item \textbf{Infinitif (Ton neutre, administratif, impersonnel)} $\rightarrow$ \textbf{L'émettrice énonce une RÈGLE valable pour TOUS.}
    \begin{itemize}
        \item Avis de recrutement (affichage, site web) : « \emph{Fournir} un dossier complet. » (Règle générale).
        \item CV (Rubrique Expériences) : « \emph{Manager} une équipe. » (Description de compétence, pas d'adresse à quelqu'un).
        \item Formulaire / Notice : « \emph{Remplir} en majuscules. » (Mode d'emploi universel).
    \end{itemize}
\end{itemize}
\end{propriete}

\begin{attention}[Piège fréquent au Probatoire / Bac Tech]
\begin{itemize}
    \item \textbf{Erreur 1 (Offre d'emploi) :} Écrire « \emph{Envoyez} votre CV par email » dans un \textbf{avis officiel d'affichage} (destinataire inconnu, pluralité de lecteurs). $\rightarrow$ \textbf{Correction :} « \emph{Envoyer} / \emph{Adresser} / \emph{Fournir} votre CV par email. » (Infinitif de consigne impersonnelle).
    \item \textbf{Erreur 2 (CV - Rubrique Expériences) :} Écrire « \emph{Je gère} les stocks » ou « \emph{Gestion} des stocks » (nominalisation lourde). $\rightarrow$ \textbf{Correction :} « \emph{Gérer} les stocks » (Infinitif d'action standard).
    \item \textbf{Erreur 3 (Lettre de motivation) :} Confondre « \emph{Veuillez} recevoir » (Prière, impératif figé 2e pl.) et « \emph{Recevoir} » (Infinitif). On écrit : « Je vous prie de \emph{recevoir} » (infinitif) OU « \emph{Veuillez recevoir} » (impératif). Pas « \emph{Veuillez de recevoir} ».
\end{itemize}
\end{attention}

\courssub{Méthode : Identifier et justifier la valeur d'un mode}

\begin{methode}[Analyser un verbe dans un texte fonctionnel (Offre / CV / Lettre)]
\begin{enumerate}
    \item \textbf{Identifier la forme :} Conjugué sans sujet ? $\rightarrow$ Impératif. Non conjugué (terminaison -er, -ir, -re, -oir) ? $\rightarrow$ Infinitif.
    \item \textbf{Identifier l'énonciateur (Qui parle ?) :} Le recruteur (offre) ? Le candidat (CV/Lettre) ? L'administration (formulaire) ?
    \item \textbf{Identifier le destinataire (À qui ?) :} Un « vous » précis (candidat ciblé) ? Tout public (avis général) ? Personne (description de soi) ?
    \item \textbf{Déterminer l'intention (Pourquoi ?) :}
    \begin{itemize}
        \item Contraindre / Ordonner $\rightarrow$ \textbf{Ordre} (Impératif).
        \item Aider / Guider $\rightarrow$ \textbf{Conseil} (Impératif).
        \item Interdire $\rightarrow$ \textbf{Défense} (Impératif négatif).
        \item Politesse / Atténuer $\rightarrow$ \textbf{Prière} (Impératif figé \emph{Veuillez/Prière de}).
        \item Expliquer une procédure générale $\rightarrow$ \textbf{Consigne impersonnelle} (Infinitif).
        \item Lister des compétences / tâches $\rightarrow$ \textbf{Énumération} (Infinitif).
        \item Exprimer un but $\rightarrow$ \textbf{Finalité} (\emph{Pour} + Infinitif).
    \end{itemize}
    \item \textbf{Justifier l'effet produit :} « L'usage de l'infinitif \emph{Fournir} rend la consigne impersonnelle et universelle, propre à un affichage public. L'impératif \emph{Envoyez} crée une relation directe et engageante dans un email de relance. »
\end{enumerate}
\end{methode}

\courssub{Entraînement : Transformation et justification}

\begin{exoresolu}[Transformer et justifier]
\textbf{Consigne :} Transformez les phrases suivantes en changeant le mode (Impératif $\leftrightarrow$ Infinitif) et justifiez le changement de ton.

\begin{enumerate}
    \item \textit{(Avis affiché dans le hall de l'entreprise) « \textbf{Déposez} vos candidatures à l'accueil avant midi. »}
    \item \textit{(Email personnel du DRH à un candidat présélectionné) « \textbf{Merci de fournir} vos références professionnelles. »}
    \item \textit{(CV - Rubrique Expérience) « \textbf{J'ai supervisé} la construction du bâtiment B. »}
\end{enumerate}
\tcblower
\textbf{Solution détaillée :}

\begin{enumerate}
    \item \textbf{Transformation :} « \textbf{Déposer} vos candidatures à l'accueil avant midi. » (Infinitif de consigne).
    \textbf{Justification :} L'avis est affiché dans un lieu public (hall). Le destinataire n'est pas un « vous » identifié, mais tout visiteur potentiel. L'infinitif instaure une \textbf{règle générale, impersonnelle et neutre}, adaptée à la signalétique administrative. L'impératif « Déposez » aurait semblé agressif ou personnellement adressé.
    
    \item \textbf{Transformation :} « \textbf{Fournissez} vos références professionnelles. » (Impératif 2e pl. - Valeur : Consigne directe / Prière selon le ton).
    \textbf{Justification :} L'email est un échange \textbf{interpersonnel} (DRH $\rightarrow$ Candidat précis). L'infinitif « Merci de fournir » (figé, politesse) est standard, mais l'impératif « Fournissez » (ou « Veuillez fournir ») marque une \textbf{attente claire, directe et urgente} de la part d'un interlocuteur identifié. C'est une consigne personnelle, pas une règle générale.
    
    \item \textbf{Transformation :} « \textbf{Superviser} la construction du bâtiment B. » (Infinitif d'action / énumération).
    \textbf{Justification :} Dans un CV, on ne raconte pas sa vie à la 1re personne (« J'ai fait »). On \textbf{liste des compétences transférables}. L'infinitif « Superviser » transforme l'action passée en \textbf{compétence actuelle et réutilisable}. Il gomme le sujet « Je » pour mettre le projecteur sur le \textbf{verbe d'action} (le \cle{savoir-faire}).
\end{enumerate}
\end{exoresolu}

\courssub{Exemple chiffré : L'énumération administrative type}

\begin{savaistu}[Le formalisme camerounais : « Pièces à fournir »]
Dans la quasi-totalité des avis de concours (Fonction Publique, Grandes Écoles, Entreprises publiques type SONARA, CAMWATER, CAMTEL) ou d'appels à candidatures au Cameroun, la liste des pièces suit invariablement le modèle de l'\textbf{infinitif impersonnel} (souvent introduit par « \emph{Pièces à fournir :} » ou « \emph{Dossier de candidature composé de :} »).

\textbf{Exemple type (Arrêté d'ouverture de concours) :}
\begin{quote}
\textbf{Pièces à fournir :}
\begin{enumerate}
    \item Une demande manuscrite adressée au Ministre / Directeur Général ;
    \item Un \textbf{CV} détaillé \textbf{daté et signé} (participe passé adjectival) ;
    \item Une \textbf{photocopie certifiée conforme} du diplôme requis ;
    \item Un \textbf{extrait d'acte de naissance} de moins de 3 mois ;
    \item Un \textbf{certificat de nationalité} camerounaise ;
    \item Un \textbf{certificat de visite médicale} délivré par un médecin agréé ;
    \item \textbf{Deux photos 4$\times$4} récentes et identiques ;
    \item Une \textbf{enveloppe timbrée} à l'adresse du candidat.
\end{enumerate}
\end{quote}

\textbf{Analyse linguistique :}
\begin{itemize}
    \item \textbf{Infinitifs (ou équivalents) :} \emph{Fournir} (titre), \emph{Adresser} (1), \emph{Dater/Signer} (implicite dans « daté et signé »), \emph{Photocopier/Certifier} (implicite), \emph{Délivrer} (implicite).
    \item \textbf{Absence totale d'impératif.} Aucune ligne ne commence par « Apportez », « Donnez », « Signez ».
    \item \textbf{Effet :} Neutralité juridique, égalité de traitement des candidats, portée générale (valable pour les 10 000 candidats potentiels).
\end{itemize}
\end{savaistu}

\begin{exercice}[Application : À vous de jouer !]
\textbf{Exercice 1 :} Dans l'offre ci-dessous, entourez les verbes à l'impératif ou à l'infinitif et indiquez leur valeur.
\textit{« \textbf{Postulez} avant le 30/06. \textbf{Joindre} CV + Lettre. \textbf{Privilégier} l'envoi par email. \textbf{Ne pas appeler} le standard. \textbf{Pour tout renseignement, consulter} le site web. »}

\textbf{Exercice 2 :} Réécrivez ces phrases de « style lettre » en « style CV » (Infinitif d'action).
\begin{enumerate}
    \item J'ai géré le budget annuel du service (50 millions FCFA).
    \item Je suis chargé de la formation des nouveaux recrues.
    \item J'ai mis en place une nouvelle procédure de sécurité.
\end{enumerate}

\textbf{Exercice 3 :} Vous êtes DRH. Rédigez une consigne d'affichage (Infinitif) puis un email de relance à un candidat précis (Impératif) pour la même action : « Fournir les originaux des diplômes ».
\end{exercice}

\begin{corrige}[Exercice 1]
\begin{itemize}
    \item \textbf{Postulez} (Impératif) : \textbf{Consigne directe / Ordre} (délai impératif). S'adresse au lecteur identifié.
    \item \textbf{Joindre} (Infinitif) : \textbf{Consigne impersonnelle} (liste de pièces, style administratif).
    \item \textbf{Privilégier} (Infinitif) : \textbf{Conseil / Recommandation} sous forme impersonnelle (plus doux que « Privilégiez »).
    \item \textbf{Ne pas appeler} (Infinitif négatif) : \textbf{Défense / Interdiction} impersonnelle (règle générale pour tous).
    \item \textbf{Consulter} (Infinitif après préposition \emph{Pour}) : \textbf{Finalité / Objectif} (But de l'action).
\end{itemize}
\end{corrige}

\begin{corrige}[Exercice 2]
\begin{enumerate}
    \item \textbf{Gérer} un budget annuel de 50 millions FCFA. (Infinitif présent : compétence durable).
    \item \textbf{Former} / \textbf{Encadrer} les nouveaux recrues. (Infinitif présent : mission habituelle).
    \item \textbf{Avoir mis en place} une nouvelle procédure de sécurité. (Infinitif passé : réalisation ponctuelle, achevée, résultat).
\end{enumerate}
\end{corrige}

\begin{corrige}[Exercice 3]
\begin{itemize}
    \item \textbf{Affichage (Infinitif - Règle générale) :} « \textbf{Fournir} les originaux des diplômes lors de l'entretien. » Ou : « \textbf{Présenter} les originaux des diplômes le jour de l'épreuve. »
    \item \textbf{Email de relance (Impératif - Relation directe) :} « Monsieur X, \textbf{Veuillez apporter} / \textbf{Apportez} les originaux de vos diplômes lors de l'entretien prévu mardi 10h. \textbf{Merci de vous présenter} à l'accueil 15 min à l'avance. »
\end{itemize}
\end{corrige}

\coursec{Production guidée : rédiger et améliorer son CV}

Dans la section précédente, nous avons étudié les valeurs de l'\cle{impératif} et de l'\cle{infinitif} : nous savons désormais qu'une consigne s'exprime à l'impératif (« Rédigez », « Présentez ») et qu'une tâche accomplie se décrit à l'infinitif dans un CV (« Assembler des circuits électriques », « Saisir des factures »). Ces deux modes vont maintenant trouver leur application immédiate : l'impératif dans les consignes que vous donne le professeur, l'infinitif dans les rubriques de votre propre CV. Passons donc de la théorie à la pratique : il est temps de rédiger, relire et améliorer votre curriculum vitae.

\courssub{Vue d'ensemble : les cinq étapes de la production}

Rédiger un bon CV ne se fait pas d'un seul jet. Comme toute production d'écrit au programme de Terminale technique, elle suit une démarche en étapes, que l'on peut comparer au protocole d'un travail pratique : on observe, on trie, on construit, on vérifie, on corrige. Voici la frise des cinq étapes que nous allons suivre.

\begin{popfigure}
\begin{center}
\begin{tikzpicture}[
  etape/.style={rectangle, rounded corners=4pt, draw=popblue, fill=popblueL, text width=2.5cm, minimum height=1.6cm, align=center, font=\small\bfseries},
  fleche/.style={-{Stealth[length=3mm]}, very thick, poporange}
]
\node[etape] (e1) at (0,0) {1. Analyser sa situation};
\node[etape, right=1.1cm of e1] (e2) {2. Collecter et trier les informations};
\node[etape, right=1.1cm of e2] (e3) {3. Rédiger un premier jet};
\node[etape, right=1.1cm of e3] (e4) {4. Relire avec la grille};
\node[etape, right=1.1cm of e4] (e5) {5. Améliorer et finaliser};
\draw[fleche] (e1) -- (e2);
\draw[fleche] (e2) -- (e3);
\draw[fleche] (e3) -- (e4);
\draw[fleche] (e4) -- (e5);
\draw[fleche, dashed, poppurple] (e5.south) to[bend right=25] node[below, font=\small\itshape, text=poppurple]{nouvelle relecture si nécessaire} (e3.south);
\end{tikzpicture}
\end{center}
\legende{Frise des cinq étapes de la production d'un CV : analyse de la situation, tri des informations, rédaction du premier jet, relecture critériée, amélioration. La flèche en pointillés rappelle que l'amélioration peut renvoyer à une nouvelle rédaction.}
\end{popfigure}

\begin{aretenir}[L'idée directrice]
Un CV efficace n'est pas écrit, il est \cle{construit}. Chaque étape a sa fonction : sans analyse préalable, le CV est vide ; sans tri, il est confus ; sans relecture, il est fautif ; sans amélioration, il reste un brouillon.
\end{aretenir}

\courssub{Étape 1 : analyser sa propre situation}

Avant d'écrire la moindre ligne, pose-toi les bonnes questions. L'employeur qui lira ton CV ne te connaît pas : que doit-il savoir de toi en trente secondes de lecture ?

\begin{methode}[Rédiger son premier CV : la fiche d'identité professionnelle]
Avant toute rédaction, remplis sur une feuille de brouillon une \cle{fiche d'identité professionnelle} en répondant par écrit à ces questions :
\begin{enumerate}
\item \textbf{Identité} : nom, prénom, date et lieu de naissance, adresse, téléphone, adresse électronique (sérieuse !).
\item \textbf{Formation} : quelle filière technique ? quel diplôme préparé ou obtenu (Probatoire, Baccalauréat technique) ? dans quel établissement ?
\item \textbf{Expériences} : quels stages en entreprise ai-je effectués ? où, quand, combien de temps ?
\item \textbf{Compétences techniques} : que sais-je faire concrètement grâce à ma filière (câbler, usiner, saisir, dessiner un plan, tenir une comptabilité...) ?
\item \textbf{Langues et informatique} : quelles langues parlées et écrites ? quels logiciels maîtrisés ?
\item \textbf{Centres d'intérêt} : quelles activités montrent ma personnalité (sport, club, association) ?
\end{enumerate}
Cette fiche remplie est ta \emph{matière première} : tu n'écriras le CV qu'ensuite, en puisant dedans.
\end{methode}

\begin{exemplebox}[Exemple de fiche d'identité professionnelle remplie]
\textbf{Identité} : Ngo Bell Sandrine, née le 12 mars 2007 à Bafoussam ; Yaoundé, quartier Nkolbisson ; tél. 6 90 00 00 00 ; sandrine.ngobell@mail.fr.\par
\textbf{Formation} : Terminale F3 (Électrotechnique), Lycée technique de Yaoundé ; Probatoire F3 obtenu en 2025.\par
\textbf{Stages} : un mois à ENEO (agence de Mvog-Ada), juillet 2025 ; deux semaines dans un atelier d'installation électrique à Mfoundi, avril 2024.\par
\textbf{Compétences} : réaliser des installations domestiques ; lire un schéma électrique ; utiliser un multimètre ; respecter les règles de sécurité.\par
\textbf{Langues} : français (courant), anglais (niveau scolaire), ewondo (parlé).\par
\textbf{Informatique} : Word, Excel, initiation à AutoCAD.\par
\textbf{Centres d'intérêt} : handball, club sciences du lycée.
\end{exemplebox}

\courssub{Étape 2 : collecter et trier les informations}

Toutes les informations de ta fiche ne se valent pas. Le tri consiste à distinguer ce qui est \cle{pertinent} pour un employeur de ce qui est inutile, voire gênant.

\begin{center}
\begin{tabularx}{\linewidth}{|l|X|X|}
\hline
\rowcolor{popblueL} \textbf{Critère} & \textbf{Pertinent (à garder)} & \textbf{Non pertinent (à écarter)} \\
\hline
Utilité pour le poste & stages, compétences techniques, diplômes & notes détaillées de chaque trimestre \\
\hline
Actualité & stage de l'année en cours & prix de récitation du CM2 \\
\hline
Image donnée & adresse électronique sérieuse & surnom, adresse fantaisiste \\
\hline
Vie privée & état civil simple si demandé & opinions, religion, détails familiaux \\
\hline
\end{tabularx}
\end{center}

\begin{attention}[Erreur fréquente : recopier le CV d'un camarade]
Un CV est \cle{personnel} : il doit refléter un parcours \emph{authentique}. Recopier le CV d'un camarade, ou inventer un stage jamais effectué, est une faute grave : l'employeur vérifie facilement les informations (un appel téléphonique suffit) et un mensonge découvert disqualifie immédiatement le candidat. Inspire-toi de la \emph{forme} des modèles, jamais du \emph{contenu} des autres.
\end{attention}

\courssub{Étape 3 : rédiger le premier jet}

Avec ta fiche triée, rédige maintenant un premier jet en respectant la structure externe (titre, rubriques clairement séparées, présentation aérée) et la structure interne (informations exactes, ordre antéchronologique pour les formations et les expériences, tâches exprimées à l'\cle{infinitif}).

\begin{exoresolu}[Rédiger la rubrique « Expérience professionnelle »]
\textbf{Énoncé.} À partir de la fiche de Sandrine Ngo Bell (exemple ci-dessus), rédigez la rubrique « Expérience professionnelle » de son CV, en respectant l'ordre antéchronologique et la formulation à l'infinitif.
\tcblower
\textbf{Solution détaillée.}\par
\textbf{EXPÉRIENCE PROFESSIONNELLE}\par
\textbf{Juillet 2025} — Stage d'un mois, ENEO, agence de Mvog-Ada (Yaoundé) :\par
— Assister les techniciens dans le raccordement des abonnés ;\par
— Relever les compteurs et vérifier les installations ;\par
— Appliquer les consignes de sécurité sur le terrain.\par
\textbf{Avril 2024} — Stage de deux semaines, atelier d'installation électrique, Mfoundi :\par
— Assembler des circuits électriques domestiques ;\par
— Lire et reproduire des schémas de câblage.\par
\emph{Remarques méthodologiques} : le stage le plus récent (2025) figure en premier (ordre antéchronologique) ; chaque tâche commence par un verbe à l'infinitif ; les lieux et les durées sont précisés, ce qui rend le parcours vérifiable et crédible.
\end{exoresolu}

\courssub{Étape 4 : relire avec une grille de critères}

Le premier jet n'est jamais la version finale. La relecture se fait à l'aide d'une \cle{grille de critères} qui transforme une impression vague (« c'est pas mal ») en diagnostic précis.

\begin{exemplebox}[Grille d'évaluation du CV sur 20 points]
\begin{center}
\begin{tabularx}{\linewidth}{|l|X|c|}
\hline
\rowcolor{popblueL} \textbf{Critère} & \textbf{Ce que l'on vérifie} & \textbf{Points} \\
\hline
Rubriques complètes & identité, formation, expérience, compétences, langues, centres d'intérêt & 5 \\
\hline
Ordre et cohérence & ordre antéchronologique, logique interne des rubriques & 4 \\
\hline
Langue & orthographe, grammaire, tâches à l'infinitif, vocabulaire technique exact & 5 \\
\hline
Mise en page & aération, alignement, titres de rubriques visibles, une page maximum & 3 \\
\hline
Adaptation & adéquation du contenu au poste visé & 3 \\
\hline
\textbf{Total} & & \textbf{20} \\
\hline
\end{tabularx}
\end{center}
\end{exemplebox}

\courssub{Étape 5 : améliorer la version de travail}

L'amélioration porte sur trois chantiers précis :
\begin{itemize}
\item \textbf{Corriger les fautes} : orthographe d'usage (« expérience » et non « expèrience »), accords, majuscules aux titres de rubriques ;
\item \textbf{Reformuler les tâches à l'infinitif} : remplacer « j'ai fait la saisie des factures » par « Saisir les factures » ; la phrase personnelle à la première personne n'a pas sa place dans les rubriques d'un CV ;
\item \textbf{Aérer la présentation} : espacer les rubriques, aligner les dates, utiliser des tirets identiques, limiter le document à une page.
\end{itemize}

\courssub{La confrontation en binômes : la lecture croisée}

Après la relecture personnelle vient la \cle{lecture croisée} : chaque élève confie son CV à un binôme qui l'évalue avec la grille. Cette étape est décisive, car le regard d'autrui repère ce que l'auteur ne voit plus.

\begin{methode}[Conduire une lecture croisée efficace]
\begin{enumerate}
\item Lis le CV de ton camarade en entier, sans rien annoter, pour avoir une impression globale.
\item Reprends la grille critère par critère et \textbf{coche} chaque critère validé.
\item Pour chaque critère \emph{non validé}, propose une \textbf{amélioration concrète} : pas « la mise en page est mauvaise », mais « aligne les dates sur la même verticale et laisse une ligne vide entre les rubriques ».
\item Formule tes remarques avec \cle{bienveillance} et \textbf{justifie}-les toujours : « je n'ai pas validé le critère langue parce que trois tâches sont à l'indicatif au lieu de l'infinitif ».
\item Remets la grille annotée à ton camarade et échangez oralement sur les points de désaccord.
\end{enumerate}
\end{methode}

\begin{attention}[Pièges de la lecture croisée]
Deux écueils à éviter : la complaisance (tout valider pour faire plaisir, ce qui prive ton camarade d'une vraie amélioration) et la critique gratuite (rejeter sans proposer de solution). Une remarque utile est toujours \emph{justifiée} par la grille et \emph{accompagnée} d'une proposition concrète.
\end{attention}

\courssub{Intégration des remarques et version définitive}

De retour en possession de ton CV annoté :
\begin{enumerate}
\item examine chaque remarque et décide, en te justifiant, de l'accepter ou non ;
\item intègre les corrections retenues dans une \cle{version définitive} propre, dactylographiée si possible ;
\item effectue une dernière relecture complète, à voix haute, avant de remettre le document au professeur.
\end{enumerate}

Le \textbf{rôle du professeur} pendant toute cette séquence est celui d'un médiateur : il circule entre les binômes, apporte des éclairages ciblés (une règle de grammaire, un exemple de formulation à l'infinitif), arbitre les désaccords de la lecture croisée et \cle{valide} les critères de la grille afin que toute la classe évalue de la même manière.

\begin{exercice}[À toi de jouer]
Remplis ta fiche d'identité professionnelle, rédige le premier jet complet de ton CV (une page maximum), puis procède à la lecture croisée avec ton binôme à l'aide de la grille sur 20 points. Remets au professeur : le premier jet, la grille annotée et la version définitive.
\end{exercice}

\begin{corrige}[Exercice : indications de correction]
La correction est individuelle, mais le professeur vérifie : la présence des six rubriques ; l'ordre antéchronologique ; la formulation des tâches à l'infinitif ; l'absence de fautes d'orthographe ; une mise en page aérée tenant sur une page ; la trace écrite des remarques du binôme et leur intégration visible dans la version définitive.
\end{corrige}

\begin{savaistu}[Des ateliers gratuits pour les jeunes diplômés]
Au Cameroun, des entreprises comme la SABC et plusieurs ONG basées à Douala organisent régulièrement des \cle{ateliers gratuits} de rédaction de CV et de préparation aux entretiens d'embauche à l'intention des jeunes diplômés. Savoir présenter son parcours en une page claire est une compétence qui se travaille toute la vie : le CV que tu rédiges aujourd'hui en classe est le premier d'une longue série !
\end{savaistu}

\coursec{Adapter le CV à une situation d'embauche}

Dans la section précédente, vous avez rédigé un premier CV complet, puis vous l'avez relu, corrigé et amélioré : structure externe claire, rubriques bien ordonnées, verbes d'action à l'infinitif, mise en page soignée. Ce CV est-il pour autant « prêt à l'emploi » ? Pas encore. Un CV n'est pas un document figé que l'on photocopie et que l'on envoie partout : c'est un \cle{outil de communication} que l'on ajuste à chaque destinataire. C'est l'objet de cette section : apprendre à \cle{adapter} son CV à une \cle{situation d'embauche} précise, compétence directement exigible au Baccalauréat technique.

\courssub{Le principe fondamental : un CV par candidature}

Commençons par une observation simple. Imaginez deux élèves de Terminale technique qui répondent à deux offres différentes. Le premier envoie exactement le même CV aux deux entreprises ; le second modifie son CV pour chaque offre. Lequel a le plus de chances d'être convoqué à un entretien ? Le second, évidemment. Pourquoi ? Parce que le recruteur ne cherche pas « un bon élève » en général : il cherche \emph{la} personne qui correspond à \emph{son} besoin précis.

\begin{aretenir}[Le principe de l'adaptation]
Il n'existe pas \textbf{un} CV unique, valable partout. Il existe \textbf{un CV par candidature}. Le fond (votre identité, votre parcours réel) reste le même et doit toujours être véridique ; c'est la \textbf{forme}, l'\textbf{ordre des rubriques} et la \textbf{sélection des informations} qui changent selon l'offre.
\end{aretenir}

\begin{attention}[L'erreur du CV générique]
L'erreur la plus fréquente consiste à envoyer le même CV générique à toutes les offres. Le recruteur repère immédiatement le manque de ciblage : un titre vague (« À la recherche d'un emploi »), des compétences qui ne correspondent pas à l'annonce, aucun mot de l'offre repris. Résultat : le CV est écarté en quelques secondes. Un recruteur consacre souvent \textbf{moins d'une minute} à la première lecture d'un CV : il faut donc que la correspondance avec le poste soit visible \emph{immédiatement}.
\end{attention}

\courssub{Analyser la situation d'embauche}

Avant de toucher à votre CV, il faut d'abord comprendre \emph{à qui} vous écrivez et \emph{pour quoi}. C'est l'étape d'\cle{analyse de la situation d'embauche}.

\begin{definition}[Situation d'embauche]
La \cle{situation d'embauche} est l'ensemble formé par :
\begin{itemize}
\item l'\textbf{employeur} (qui recrute ? une entreprise, une administration, une ONG, un particulier ?) ;
\item le \textbf{poste} proposé (intitulé, tâches à accomplir, lieu de travail) ;
\item le \textbf{profil recherché} (diplôme exigé, expérience, compétences techniques, qualités personnelles, langues) ;
\item les \textbf{conditions de candidature} (documents à fournir, date limite, mode de dépôt).
\end{itemize}
\end{definition}

Analyser la situation d'embauche, c'est donc répondre à trois questions, dans cet ordre :
\begin{enumerate}
\item \textbf{Qui recrute ?} Une grande entreprise de BTP n'attend pas le même discours qu'une petite ONG humanitaire. La première valorise la rigueur technique et le rendement ; la seconde valorise l'engagement, la polyvalence et le sens du relationnel.
\item \textbf{Quel poste ?} Il faut repérer l'intitulé exact du poste et les tâches annoncées : ce sont elles que votre CV doit viser.
\item \textbf{Quel profil est recherché ?} On relève dans l'annonce le diplôme, l'expérience, les compétences (techniques, linguistiques, informatiques) et les qualités humaines demandées.
\end{enumerate}

\begin{methode}[Adapter son CV en quatre étapes]
\begin{enumerate}
\item \textbf{Surligner les mots-clés de l'offre} : soulignez dans l'annonce le poste, les compétences exigées, les langues, les logiciels, les qualités demandées (ex. : « bilingue », « maîtrise de Sage Saari », « sens de l'organisation »).
\item \textbf{Établir la correspondance} : pour chaque mot-clé, cherchez dans votre parcours l'expérience, la formation ou la compétence qui y répond.
\item \textbf{Vérifier que chaque mot-clé trouve une réponse dans le CV} : si l'offre demande « anglais courant », la mention « anglais courant » doit apparaître clairement dans votre rubrique « Langues ».
\item \textbf{Réorganiser et reformuler} : placez en tête ce qui répond le mieux à l'offre, adaptez le titre, reprenez le vocabulaire exact de l'annonce.
\end{enumerate}
\end{methode}

\courssub{Sélectionner et réorganiser les informations}

Une fois l'offre analysée, trois opérations concrètes permettent d'adapter le CV.

\textbf{1. Sélectionner les informations.} Votre CV n'a pas à tout dire : il doit dire \emph{ce qui intéresse ce recruteur-là}. Mettez en tête les compétences et les expériences qui répondent à l'offre ; réduisez ou supprimez ce qui n'a aucun rapport avec le poste. Un stage dans un atelier de menuiserie est décisif pour un poste de technicien en bois, mais inutile pour un poste de secrétaire comptable : il peut alors être résumé en une ligne ou omis.

\textbf{2. Adapter le titre.} Le titre du CV est la première ligne lue par le recruteur. Il doit \textbf{préciser le poste visé}, en reprenant l'intitulé de l'annonce. Comparez :
\begin{itemize}
\item Titre générique (à éviter) : « Curriculum Vitae » ou « À la recherche d'un emploi ».
\item Titre adapté (à privilégier) : « Candidature au poste de \textbf{technicien en froid et climatisation} ».
\end{itemize}

\textbf{3. Réorganiser les rubriques.} L'ordre des rubriques n'est pas immuable : il dépend de votre profil et de l'offre.
\begin{itemize}
\item Si votre \textbf{expérience} est plus parlante que vos diplômes (plusieurs stages, emplois saisonniers, apprentissage en atelier), placez la rubrique « Expérience professionnelle » \textbf{avant} « Formation ».
\item Si vous êtes \textbf{jeune diplômé} sans expérience, placez « Formation » en premier : c'est votre meilleur argument.
\end{itemize}

\begin{exemplebox}[Exemple chiffré : le poste de secrétaire comptable bilingue]
Une entreprise de Douala recrute un(e) « \textbf{secrétaire comptable bilingue} ». L'annonce exige : anglais courant, maîtrise du logiciel Sage Saari, sens de l'organisation. Votre rubrique « Compétences » doit donc être réorganisée ainsi :
\begin{enumerate}
\item Anglais courant (lu, écrit, parlé) — niveau B2 ;
\item Maîtrise de Sage Saari (comptabilité, facturation) ;
\item Word, Excel, PowerPoint ;
\item Sens de l'organisation et de la confidentialité.
\end{enumerate}
Avant l'adaptation, « anglais courant » figurait peut-être en quatrième position, noyé entre « cuisine » et « football ». Après adaptation, les deux compétences décisives passent \textbf{en tête} : en moins de dix secondes de lecture, le recruteur voit que le profil correspond à la totalité de ses exigences principales.
\end{exemplebox}

\courssub{Reprendre le vocabulaire de l'offre}

Un point fin, souvent négligé : il faut \cle{reprendre les termes exacts de l'annonce}. Si l'offre demande un « technicien de maintenance », n'écrivez pas « bricoleur polyvalent » mais « technicien de maintenance ». Si l'offre parle de « gestion de stock », employez « gestion de stock » et non « rangement des marchandises ».

Pourquoi ? D'abord parce que le recruteur \emph{reconnaît} ses propres mots : la correspondance du profil saute aux yeux. Ensuite parce que, dans les grandes structures, les CV sont parfois triés à partir de \cle{mots-clés} : un CV qui ne contient pas les termes de l'offre risque d'être éliminé avant même d'être lu.

\begin{attention}[Adapter ne signifie pas mentir]
Reprendre le vocabulaire de l'offre et mettre en valeur ses atouts ne doit jamais conduire à \textbf{inventer} un diplôme, une expérience ou un niveau de langue. Toute fausse information peut être vérifiée (entretien, test pratique, demande d'attestation) et entraîne le rejet immédiat de la candidature, voire un licenciement. Adapter, c'est \emph{choisir et ordonner} la vérité, pas la travestir.
\end{attention}

\courssub{Schéma de correspondance entre l'offre et le CV}

Le schéma suivant visualise la démarche d'adaptation : à chaque exigence de l'offre (colonne de gauche) doit correspondre une réponse précise dans une rubrique du CV (colonne de droite).

\begin{popfigure}
\begin{center}
\begin{tikzpicture}[
  exig/.style={rectangle, draw=popblue, fill=popblueL, rounded corners=2pt,
    text width=4.6cm, align=center, minimum height=1.05cm, font=\small},
  rubr/.style={rectangle, draw=poporange, fill=poporangeL, rounded corners=2pt,
    text width=5.2cm, align=center, minimum height=1.05cm, font=\small},
  lien/.style={-{Stealth[length=2.5mm]}, thick, popdark}]
% Titres des colonnes
\node[font=\bfseries\small, popblue] at (0,0.9) {Exigences de l'offre};
\node[font=\bfseries\small, poporange] at (8,0.9) {Rubriques du CV adapté};
% Colonne gauche : exigences
\node[exig] (e1) at (0,0) {Poste : technicien en froid et climatisation};
\node[exig] (e2) at (0,-1.6) {Diplôme : BT froid et climatisation};
\node[exig] (e3) at (0,-3.2) {Expérience : 1 an minimum};
\node[exig] (e4) at (0,-4.8) {Qualités : ponctualité, travail en équipe};
% Colonne droite : rubriques
\node[rubr] (r1) at (8,0) {Titre : « Candidature au poste de technicien en froid et climatisation »};
\node[rubr] (r2) at (8,-1.6) {Formation : BT froid et climatisation, LTP Bafoussam, 2025};
\node[rubr] (r3) at (8,-3.2) {Expérience : stage de 6 mois, froid industriel};
\node[rubr] (r4) at (8,-4.8) {Compétences / qualités : ponctualité, esprit d'équipe};
% Flèches de correspondance
\draw[lien] (e1.east) -- (r1.west);
\draw[lien] (e2.east) -- (r2.west);
\draw[lien] (e3.east) -- (r3.west);
\draw[lien] (e4.east) -- (r4.west);
\end{tikzpicture}
\end{center}
\legende{Diagramme à double entrée : chaque exigence de l'offre d'emploi (colonne bleue) doit trouver une réponse visible dans une rubrique du CV adapté (colonne orange). Les flèches matérialisent la correspondance que le recruteur doit percevoir en moins d'une minute.}
\end{popfigure}

\courssub{Cas pratique : un même parcours, deux offres différentes}

Voyons maintenant, sur un cas complet, comment un \emph{même} parcours donne lieu à deux CV \emph{différents}.

\textbf{Le parcours de départ.} Aminatou, 19 ans, vient d'obtenir son Baccalauréat technique (spécialité Génie civil) au Lycée technique de Bafoussam. Son parcours comprend : un stage de deux mois sur un chantier de construction (mesures, coffrage), six mois d'aide bénévole dans une association de quartier (organisation de distributions, accueil du public), un bon niveau d'anglais, la maîtrise d'AutoCAD et d'Excel, le permis B.

\textbf{Offre 1 : entreprise de BTP à Bafoussam.} Poste : « technicien(ne) de chantier ». Exigences : Bac technique Génie civil, lecture de plans, AutoCAD, expérience de chantier, permis B apprécié.

\textbf{Offre 2 : ONG humanitaire à Maroua.} Poste : « assistant(e) logistique ». Exigences : sens de l'organisation, contact avec les bénéficiaires, Excel, anglais, mobilité.

\begin{center}
\begin{tabularx}{\linewidth}{|l|X|X|}
\hline
\rowcolor{popblueL} \textbf{Élément du CV} & \textbf{CV 1 : BTP (Bafoussam)} & \textbf{CV 2 : ONG (Maroua)} \\
\hline
Titre & Candidature au poste de technicienne de chantier & Candidature au poste d'assistante logistique \\
\hline
Rubrique en tête & Expérience : stage chantier (mesures, coffrage) & Expérience : bénévolat associatif (organisation, accueil) \\
\hline
Compétences mises en avant & Lecture de plans, AutoCAD, topographie de base & Organisation de distributions, Excel, relationnel \\
\hline
Langues & Français, anglais (bon niveau) & Anglais courant, français (langues de travail) \\
\hline
Informations réduites & Bénévolat résumé en une ligne & Stage chantier résumé en une ligne \\
\hline
Atout supplémentaire & Permis B (déplacements chantier) & Permis B + disponibilité et mobilité nationale \\
\hline
\end{tabularx}
\end{center}

On le constate : \textbf{aucune information n'est inventée}. Seuls l'ordre, la sélection et le vocabulaire changent. Aminatou présente à chaque recruteur le visage de son parcours qui correspond à son besoin.

\begin{exoresolu}[Adapter un CV à une offre]
\textbf{Énoncé.} Voici une offre : « Une grande entreprise agro-industrielle de Garoua recrute un(e) secrétaire de direction. Profil : Bac technique, excellente maîtrise de Word et Excel, français impeccable, anglais souhaité, discrétion exigée. » Voici les compétences brutes d'un candidat : « football ; Word et Excel ; chant choral ; anglais scolaire ; français courant ; discrétion ; conduite de moto ». 1) Sélectionnez les compétences à conserver. 2) Ordonnez-les. 3) Rédigez le titre du CV.
\tcblower
\textbf{Solution détaillée.}
1) \textbf{Sélection} : on conserve Word et Excel, l'anglais, le français, la discrétion (ce sont les mots-clés de l'offre). On écarte le football, le chant choral et la conduite de moto, sans rapport avec le poste (ils pourraient éventuellement figurer dans « Centres d'intérêt », en une ligne).
2) \textbf{Ordre} (du plus décisif au moins décisif, selon l'offre) : maîtrise de Word et Excel ; français impeccable (lu, écrit, parlé) ; anglais (bon niveau scolaire) ; discrétion et sens de la confidentialité.
3) \textbf{Titre du CV} : « Candidature au poste de secrétaire de direction ».
Le candidat a ainsi répondu, point par point, aux quatre exigences de l'annonce, en reprenant son vocabulaire exact.
\end{exoresolu}

\courssub{Articulation avec la lettre de motivation}

Le CV ne voyage jamais seul : il est accompagné d'une \cle{lettre de motivation}. Il faut bien distinguer leurs rôles, qui sont complémentaires :
\begin{itemize}
\item le \textbf{CV informe} : il présente, de manière brève et structurée, les faits (identité, formation, expériences, compétences) ;
\item la \textbf{lettre argumente} : elle explique \emph{pourquoi} vous postulez, \emph{pourquoi} ce poste vous attire et \emph{en quoi} votre parcours répond au besoin de l'employeur.
\end{itemize}

\begin{aretenir}[CV et lettre : deux documents complémentaires]
Le CV dit \emph{ce que} vous avez fait ; la lettre dit \emph{pourquoi} cela vous rend apte au poste. Les deux doivent être cohérents et adaptés à la même offre : le titre du CV, les mots-clés repris et les arguments de la lettre forment un ensemble uni. Au Baccalauréat technique, savoir distinguer ces deux fonctions est un atout pour l'épreuve de production d'écrits.
\end{aretenir}

\begin{savaistu}[Le CV au Baccalauréat technique]
Au Baccalauréat technique, l'épreuve de production d'écrits peut vous demander de \textbf{rédiger un CV à partir d'une offre d'emploi fournie dans le sujet}. Le correcteur évalue alors précisément les compétences travaillées dans cette section : avez-vous repris l'intitulé exact du poste dans le titre ? Avez-vous sélectionné les informations qui répondent à l'offre ? Avez-vous ordonné les rubriques de manière pertinente ? Un CV générique, sans lien visible avec l'annonce du sujet, est lourdement pénalisé : l'adaptation est donc un critère de notation, pas un simple conseil pratique.
\end{savaistu}

\begin{exercice}[À vous de jouer]
On vous donne l'offre suivante : « Hôtel La Falaise, Yaoundé, recrute un(e) réceptionniste. Profil : accueil de la clientèle, français et anglais courants, maîtrise de l'outil informatique, présentation soignée, disponibilité (travail le week-end). »
1) Relevez les cinq mots-clés de cette offre.
2) À partir de votre propre parcours, indiquez pour chaque mot-clé l'information que vous feriez figurer dans votre CV et dans quelle rubrique.
3) Rédigez le titre adapté de votre CV.
4) Expliquez dans quel ordre vous placerez les rubriques « Formation » et « Expérience », en justifiant votre choix.
\end{exercice}

\begin{corrige}[Exercice : adapter son CV à l'offre de réceptionniste]
1) \textbf{Mots-clés} : accueil de la clientèle ; français et anglais courants ; maîtrise de l'outil informatique ; présentation soignée ; disponibilité (week-end).
2) \textbf{Correspondances attendues} (exemple de réponse) : accueil de la clientèle $\rightarrow$ rubrique « Expérience » (stage à l'accueil d'une administration ou d'un commerce) ; langues $\rightarrow$ rubrique « Langues » : « français courant, anglais courant » ; informatique $\rightarrow$ rubrique « Compétences » : « Word, Excel, messagerie » ; présentation $\rightarrow$ rubrique « Qualités » : « présentation soignée, sens du contact » ; disponibilité $\rightarrow$ ligne « Informations complémentaires » : « disponible immédiatement, week-ends compris ».
3) \textbf{Titre} : « Candidature au poste de réceptionniste — Hôtel La Falaise, Yaoundé ».
4) \textbf{Ordre des rubriques} : si le candidat a déjà une expérience d'accueil (stage, emploi saisonnier), « Expérience » passe avant « Formation », car c'est elle qui répond le mieux à l'offre ; s'il est jeune diplômé sans expérience, « Formation » passe en premier et les qualités personnelles (contact, présentation) sont mises en valeur.
\end{corrige}

En résumé : adapter son CV, c'est d'abord \emph{lire} l'offre (qui recrute ? quel poste ? quel profil ?), puis \emph{sélectionner} et \emph{ordonner} ses informations, \emph{reprendre} le vocabulaire exact de l'annonce, et enfin vérifier que chaque exigence trouve une réponse visible en moins d'une minute. La section suivante vous permettra de confronter vos productions à celles de vos camarades, de les améliorer encore, puis de faire le bilan de l'unité à travers le compte rendu et la remédiation.

\coursec{Compte rendu, remédiation et évaluation}

Après avoir appris à \emph{adapter son CV à une offre d'emploi précise} (section précédente), nous arrivons à l'étape finale et décisive : \textbf{valider les acquis, corriger les lacunes persistantes et évaluer la compétence en situation réelle}. Cette séance n'est pas une simple formalité administrative ; c'est le moment où l'élève prend conscience de sa progression, où l'enseignant identifie les points de blocage pour y remédier efficacement, et où l'on vérifie que l'élève est capable de produire un CV opérationnel, prêt à être déposé.

\begin{methode}[Réaliser le compte rendu collectif des acquis]
\textbf{Objectif :} Synthétiser oralement puis par écrit l'ensemble des savoirs et savoir-faire travaillés durant la leçon, en formulant chaque acquis sous forme de capacité personnelle.
\begin{enumerate}
    \item \textbf{Rappel chronologique :} L'enseignant récapitule les grandes étapes : découverte des textes fonctionnels/iconiques $\rightarrow$ structure externe $\rightarrow$ structure interne $\rightarrow$ langue (impératif/infinitif) $\rightarrow$ production guidée $\rightarrow$ adaptation à l'offre.
    \item \textbf{Formulation individuelle :} Chaque élève note sur une fiche 5 à 7 acquis en commençant par \guillemotleft Je sais... \guillemotright ou \guillemotleft Je sais faire... \guillemotright.
    \item \textbf{Mise en commun :} Au tableau, construction collective de la liste de référence. L'enseignant reformule pour garantir la terminologie exacte (ex. : \guillemotleft Je sais identifier les rubriques obligatoires et optionnelles \guillemotright plutôt que \guillemotleft Je sais les parties du CV \guillemotright).
    \item \textbf{Validation :} La liste finale est recopiée par tous dans le cahier de leçon sous le titre \guillemotleft Bilan de la leçon : Le CV \guillemotright.
\end{enumerate}
\end{methode}

\courssub{Compte rendu collectif et construction de la fiche-bilan}

Le compte rendu permet de \textbf{nommer} ce qui a été appris. Sans nomination explicite, les savoirs restent implicites et difficilement transférables.

\begin{exemplebox}[Exemples de formulations « Je sais / Je sais faire » validées collectivement]
\begin{itemize}
    \item Je sais \textbf{distinguer la structure externe} (mise en page, police, longueur, photo) \textbf{de la structure interne} (rubriques, ordre, contenu) d'un CV.
    \item Je sais \textbf{choisir l'ordre des rubriques} (antéchronologique, fonctionnel, mixte) en fonction de mon profil (jeune diplômé vs expérimenté) et de l'offre.
    \item Je sais \textbf{rédaction des items d'expérience et de formation} en utilisant des \textbf{verbes d'action à l'infinitif} (ex. : \guillemotleft Gérer une équipe de 5 personnes \guillemotright) ou à l'\textbf{impératif} dans la rubrique \guillemotleft Compétences \guillemotright (ex. : \guillemotleft Maîtriser la suite Office \guillemotright).
    \item Je sais \textbf{analyser une offre d'emploi} (mots-clés, compétences requises, valeurs de l'entreprise) pour \textbf{personnaliser mon CV} (titre, profil, sélection/mise en avant des expériences).
    \item Je sais \textbf{relire mon CV} pour éliminer les fautes d'orthographe, les abréviations obscures, les informations inutiles (état civil détaillé, références) et vérifier la cohérence des dates.
\end{itemize}
\end{exemplebox}

À l'issue de ce tour de table, la classe construit \textbf{la fiche-bilan « Les 10 règles d'or du CV réussi »}. Cette fiche, affichée en classe et photocopiée, sert de \emph{mémo} permanent pour les rédactions futures (stages, job d'été, post-Bac).

\begin{aretenir}[Les 10 règles d'or du CV réussi (Fiche-bilan collective)]
\begin{enumerate}
    \item \textbf{Un CV = Une page maximum} (sauf expérience > 10 ans justifiée).
    \item \textbf{Mise en page aérée, police lisible} (Arial, Calibri, 11-12 pts), \textbf{pas de couleur excessive}.
    \item \textbf{Photo professionnelle} (si usage local) : fond neutre, tenue correcte, sourire naturel.
    \item \textbf{Titre ciblé} : Intitulé du poste visé + spécialité (ex. : \guillemotleft Technicien en Maintenance Industrielle \guillemotright).
    \item \textbf{Profil / Accroche} : 3 lignes max, \emph{sur mesure} pour l'offre (compétences clés + valeur ajoutée).
    \item \textbf{Ordre antéchronologique} (du plus récent au plus ancien) dans \textbf{Expériences} et \textbf{Formation}.
    \item \textbf{Verbes d'action à l'infinitif} pour décrire les missions/réalisations (pas de \guillemotleft J'ai fait \guillemotright).
    \item \textbf{Zéro faute} : relecture à voix haute + correcteur + relecture par un tiers.
    \item \textbf{Adaptation systématique} : Mots-clés de l'offre repris, expériences pertinentes mises en gras.
    \item \textbf{Fichier PDF} nommé \guillemotleft Nom\_Prenom\_CV\_Poste.pdf \guillemotright + mail d'accompagnement poli.
\end{enumerate}
\end{aretenir}

\begin{popfigure}
\begin{tikzpicture}[
    mindmap,
    concept color=popblue,
    level 1/.style={sibling angle=90, level distance=4.5cm},
    level 2/.style={sibling angle=45, level distance=3.5cm},
    level 3/.style={sibling angle=30, level distance=3cm},
    every node/.style={concept, execute at begin node=\hskip0pt, font=\small},
    popblue/.style={concept color=popblue},
    poporange/.style={concept color=poporange},
    popgreen/.style={concept color=popgreen},
    poppurple/.style={concept color=poppurple},
    poppink/.style={concept color=poppink},
    popgold/.style={concept color=popgold},
]
\node[popblue, align=center]{Le CV\\Terminale Tech}
    child[poporange] {node {Structure Externe}
        child {node {Mise en page (1 page)}
        child {node {Police, marges, interlignes}}}
        child {node {Photo pro (optionnelle)}
        child {node {Format PDF}}}
        child {node {En-tête : Identité + Titre}}}
    child[popgreen] {node {Structure Interne}
        child {node {Rubriques obligatoires}
            child {node {Expériences (antéchronol.)}
            child {node {Verbes infinitif}}}
            child {node {Formation}}
            child {node {Compétences (tech/transv.)}}}
        child {node {Rubriques optionnelles}
            child {node {Langues, Centres d'intérêt}}
            child {node {Projets, Certificats}}}
        child {node {Profil / Accroche}}}
    child[poppurple] {node {Langue : Modes}
        child {node {Infinitif : Missions, Tâches}
            child {node {Gérer, Concevoir, Analyser}}}
        child {node {Impératif : Compétences, Objectifs}
            child {node {Maîtriser, Assurer, Garantir}}}
        child {node {Concision, Impersonnalité}}}
    child[poppink] {node {Adaptation à l'Offre}
        child {node {Analyse mots-clés}
            child {node {Compétences requises}}}
        child {node {Personnalisation}
            child {node {Titre, Profil, Ordre rubriques}}}
        child {node {Lettre de motivation associée}}};
\end{tikzpicture}
\legende{Carte mentale récapitulative de la leçon « Production d'écrits utilitaires : rédaction d'un CV » (Terminale Technique).}
\end{popfigure}

\courssub{Analyse des difficultés récurrentes et remédiation ciblée}

L'analyse des productions de la classe (brouillons, versions intermédiaires, CV adaptés) révèle systématiquement des \textbf{points de blocage partagés}. Les identifier collectivement désamorce la honte de l'erreur et permet une remédiation efficace.

\begin{attention}[Remédiation $\neq$ Punition]
L'erreur fréquente, chez l'élève comme parfois chez l'enseignant, est de considérer la remédiation comme une \textbf{sanction} ou une \guillemotleft séance de rattrapage pour les nuls \guillemotright.
\textbf{Realité pédagogique :} La remédiation est un \textbf{second apprentissage}, ciblé, explicite et bienveillant, sur les points \emph{non encore maîtrisés}. Elle concerne \emph{tous} les élèves sur des aspects différents. L'élève qui a réussi la structure peut avoir des lacunes en langue ; celui qui maîtrise la langue peut peiner sur l'adaptation.
\end{attention}

\textbf{Difficultés types observées en Terminale Technique :}
\begin{itemize}
    \item \textbf{Langue :} Confusion Infinitif / Impératif / Participe passé. Ex. : \guillemotleft *Géré une équipe \guillemotright (participe passé au lieu de l'infinitif \guillemotleft Gérer \guillemotright) ; \guillemotleft *Maîtrise l'anglais \guillemotright (indicatif/impératif 2e pers. au lieu de l'infinitif \guillemotleft Maîtriser \guillemotright).
    \item \textbf{Structure :} Mélange des rubriques (stages dans Formation, loisirs dans Compétences), ordre chronologique au lieu d'antéchronologique, dates approximatives (année seule sans mois).
    \item \textbf{Contenu :} Phrases complètes au lieu de items nominaux/verbaux (\guillemotleft J'ai été chargé de la gestion des stocks \guillemotright $\rightarrow$ \guillemotleft Gestion des stocks \guillemotright) ; tâches génériques sans résultat chiffré (\guillemotleft Vente \guillemotright vs \guillemotleft Augmentation du CA de 15\% sur 6 mois \guillemotright).
    \item \textbf{Adaptation :} CV \guillemotleft passe-partout \guillemotright non modifié pour l'offre test ; titre du CV inchangé ; compétences clés de l'offre absentes du CV.
\end{itemize}

\textbf{Remédiation 1 : Réemploi sur l'Impératif et l'Infinitif (Atelier de 20 min)}

\begin{exoresolu}[Exercice de réemploi : Choisir le mode verbal adapté]
\textbf{Consigne :} Complétez les items de CV suivants en conjuguant le verbe entre parenthèses à l'\textbf{infinitif} (missions/réalisations) ou à l'\textbf{impératif présent} (compétences/qualités), selon la rubrique indiquée.
\begin{enumerate}
    \item \textbf{[Rubrique Expériences]} Stage chez CAMWATER -- \textit{(surveiller)} la qualité de l'eau ; \textit{(rédiger)} des rapports quotidiens ; \textit{(participer)} à la maintenance des pompes.
    \item \textbf{[Rubrique Compétences Techniques]} \textit{(maîtriser)} la soudure TIG/MIG ; \textit{(utiliser)} AutoCAD 2D/3D ; \textit{(programmer)} des automates Siemens S7.
    \item \textbf{[Rubrique Compétences Transverses]} \textit{(communiquer)} en anglais technique (niveau B2) ; \textit{(gérer)} les priorités sous contrainte de délai ; \textit{(respecter)} les normes QHSE.
    \item \textbf{[Rubrique Projet / Objectif]} \textit{(intégrer)} une équipe de maintenance préventive ; \textit{(développer)} mes compétences en diagnostic de pannes hydrauliques.
\end{enumerate}
\tcblower
\textbf{Solution et Justification :}
\begin{enumerate}
    \item \textbf{Surveiller / Rédiger / Participer} (Infinitif). \emph{Justification :} Rubrique \guillemotleft Expériences \guillemotright $\rightarrow$ description de missions/réalisations passées $\rightarrow$ \textbf{Infinitif} (valeur de nom d'action, impersonnel, standard du CV).
    \item \textbf{Maîtriser / Utiliser / Programmer} (Infinitif). \emph{Justification :} Rubrique \guillemotleft Compétences \guillemotright $\rightarrow$ énumération de savoir-faire $\rightarrow$ \textbf{Infinitif} (usage dominant moderne, neutre, liste). \textit{Note : L'impératif « Maîtrisez » est possible mais moins standard dans une liste à puces.}
    \item \textbf{Communiquer / Gérer / Respecter} (Infinitif). \emph{Même raisonnement : liste de compétences.}
    \item \textbf{Intégrer / Développer} (Infinitif). \emph{Justification :} Rubrique \guillemotleft Projet / Objectif \guillemotright $\rightarrow$ projection future $\rightarrow$ \textbf{Infinitif} (valeur de finalité, d'objectif). L'impératif (\guillemotleft Intégrez \guillemotright) s'adresserait au recruteur (inapproprié).
\end{enumerate}
\emph{Règle mnémotechnique :} \textbf{Dans 95\% des rubriques d'un CV moderne, c'est l'INFINITIF qui s'impose.} L'impératif est réservé aux titres de rubriques (\guillemotleft \textbf{Maîtriser} l'anglais \guillemotright comme intitulé) ou aux consignes.
\end{exoresolu}

\textbf{Remédiation 2 : Reconstitution de structure (Jeu de cartes / Tri)}

\begin{experience}[Atelier « Remets de l'ordre dans ce CV »]
\textbf{Matériel :} Un CV \guillemotleft brouillon \guillemotright découpé en bandes (une bande = une rubrique ou un item) : Identité, Titre, Profil, Expérience 1 (stage), Expérience 2 (job d'été), Formation (Bac, BTS), Compétences Techniques, Langues, Centres d'intérêt, Références. Les bandes sont mélangées dans une enveloppe par binôme.
\textbf{Protocole :}
\begin{enumerate}
    \item Trier les bandes en \guillemotleft Rubriques principales \guillemotright vs \guillemotleft Contenu des rubriques \guillemotright.
    \item Ordonner les rubriques principales selon la \textbf{structure standard antéchronologique} adaptée à un jeune diplômé (Formation avant Expérience ? ou Expérience avant Formation ? $\rightarrow$ Débat : si stages significatifs $\rightarrow$ Expérience d'abord).
    \item À l'intérieur de \guillemotleft Expériences \guillemotright et \guillemotleft Formation \guillemotright, ordonner \textbf{antéchronologiquement} (plus récent en haut).
    \item Coller le CV reconstitué sur une affiche A3 et justifier l'ordre choisi à l'oral.
\end{enumerate}
\textbf{Observations attendues :} Débat sur la place de \guillemotleft Centres d'intérêt \guillemotright (fin), pertinence de \guillemotleft Références \guillemotright (souvent \guillemotleft Sur demande \guillemotright), cohérence des dates.
\textbf{Interprétation :} L'exercice ancre la \textbf{macro-structure} (architecture du document) et la \textbf{micro-structure} (ordonnancement temporel interne).
\end{experience}

\courssub{Correction individualisée des versions définitives}

Pendant que les groupes travaillent sur la remédiation, l'enseignant procède à la \textbf{correction individualisée} des versions définitives des CV (remises en début de séance ou séance précédente). Cette correction suit une \textbf{grille critériée} connue des élèves (transparence de l'évaluation).

\begin{exemplebox}[Barème type Baccalauréat Technique — Production d'un CV (20 points)]
\begin{center}
\begin{tabularx}{\linewidth}{|l|X|c|}\hline
\textbf{Critère} & \textbf{Indicateurs de réussite} & \textbf{Points} \\ \hline
\textbf{1. Respect de la situation} & Titre ciblé, Profil adapté, Mots-clés de l'offre intégrés, Cohérence profil/poste & 4 \\ \hline
\textbf{2. Structure \& Organisation} & Rubriques complètes/pertinentes, Ordre logique (antéchronol.), Hiérarchie visuelle claire (titres, puces), Longueur 1 page & 6 \\ \hline
\textbf{3. Langue \& Formulation} & Verbes d'action à l'infinitif (missions), Vocabulaire technique précis, Zéro faute (orthographe, grammaire, syntaxe), Concision (style télégraphique) & 6 \\ \hline
\textbf{4. Présentation \& Soin} & Mise en page aérée/équilibrée, Police pro, Photo adaptée (si requise), Fichier PDF propre, Nom de fichier normé & 4 \\ \hline
\end{tabularx}
\end{center}
\par\small\textit{Grille d'évaluation sommative alignée sur les attendus du Probatoire/Bac Technique.}\par
\end{exemplebox}

\begin{methode}[Correction individualisée efficace (en 5 min par élève)]
\begin{enumerate}
    \item \textbf{Lecture rapide (30s) :} Impression globale (mise en page, longueur, photo).
    \item \textbf{Contrôle structure (1 min) :} Rubriques présentes ? Ordre ? Dates cohérentes ?
    \item \textbf{Analyse langue/contenu (2 min) :} Survol des items $\rightarrow$ repérage verbes (infinitif ?), vocabulaire précis ? Fautes ? Adaptation à l'offre (mots-clés) ?
    \item \textbf{Feedback ciblé (1,5 min) :} \textbf{1 point fort} + \textbf{2 axes de progression prioritaires} notés en marge ou sur fiche annexe. \emph{Ne pas tout corriger à la place de l'élève.}
    \item \textbf{Note sur 20} reportée sur grille.
\end{enumerate}
\end{methode}

\courssub{Évaluation sommative : rédaction en temps limité}

L'évaluation sommative valide la \textbf{compétence transférable} : \guillemotleft Produire un CV adapté à une offre inédite, en autonomie, en temps contraint \guillemotright. C'est la simulation la plus proche de la réalité (concours, réponse à une annonce).

\begin{exercice}[Sujet d'évaluation sommative type (Durée : 1h30)]
\textbf{Situation :} Vous êtes élève en Terminale Technique, spécialité \textbf{Électrotechnique} (ou votre spécialité). Vous répondez à l'offre ci-dessous pour un \textbf{stage de fin d'études (2 mois)} ou un \textbf{premier emploi} (selon contexte).
\vspace{0.5em}

\noindent\textbf{OFFRE D'EMPLOI / STAGE — ENTREPRISE \guillemotleft ENERGICAM \guillemotright (Douala)}
\begin{quote}
\textbf{Poste :} Technicien(ne) Maintenance Électrique / Automatisme (Stage / Junior)
\textbf{Missions :} Maintenance préventive/curative lignes de production ; Diagnostic pannes variateurs/automates ; Lecture plans électriques/schéma câblage ; Rédaction rapports d'intervention ; Respect strict consignes QHSE.
\textbf{Profil :} Bac+2/3 Électrotech/Automatisme (ou Terminale Tech option Électrotech pour stage). Maîtrise : Schémas unifilaires, Automates (Siemens/Schneider), Variateurs de vitesse, Outils de mesure (multimètre, oscilloscope). Anglais technique lu. Rigueur, réactivité, esprit d'équipe, permis B.
\textbf{Contact :} RH@energicam.cm — Objet : \guillemotleft Candidature Maintenance - [Votre Nom] \guillemotright
\end{quote}

\textbf{Consignes :}
\begin{enumerate}
    \item Analysez l'offre (surlignez 5 mots-clés techniques, 3 compétences transverses, 1 valeur entreprise).
    \item Rédigez votre \textbf{CV complet, adapté}, sur \textbf{une page maximum}.
    \item Respectez la structure externe (mise en page, police, en-tête, titre, photo optionnelle) et interne (rubriques ordonnées, items à l'infinitif, dates précises).
    \item Inventez des détails plausibles pour vos stages/projets scolaires (noms d'entreprises locales, réalisations chiffrées).
    \item Enregistrez/Remettez en format \textbf{PDF} nommé \guillemotleft NOM\_Prenom\_CV\_Energicam.pdf \guillemotright.
\end{enumerate}
\end{exercice}

\courssub{Ouverture : Le CV dans le portfolio de l'élève}

Le CV rédigé, corrigé, évalué et amélioré ne doit pas finir au fond d'un tiroir. C'est le \textbf{premier document professionnel} de l'élève technicien.

\begin{savaistu}[Le CV, outil vivant du portfolio professionnel]
Un CV bien conservé et \textbf{mis à jour à chaque nouvelle expérience} (stage, job d'été, projet bénévole, certification, permis, formation continue) fait gagner un \textbf{temps précieux} lors des dépôts de candidature futurs (BTS, Licence Pro, Emploi, Concours fonction publique).
\begin{itemize}
    \item \textbf{Format maître :} Conserver le fichier source (Word / LibreOffice / LaTeX) modifiable + version PDF à jour.
    \item \textbf{Stockage :} Clé USB \guillemotleft Portfolio \guillemotright + Cloud (Drive, OneDrive) + Mail personnel (envoi à soi-même).
    \item \textbf{Rituel de mise à jour :} Tous les 6 mois, ou après chaque stage/projet majeur : ajouter l'expérience, actualiser les compétences, vérifier les coordonnées, rafraîchir la photo si besoin.
    \item \textbf{Versions dérivées :} À partir du maître, créer versions \guillemotleft Stage \guillemotright, \guillemotleft Emploi \guillemotright, \guillemotleft Concours \guillemotright (ajout notes, mentions).
\end{itemize}
\end{savaistu}

\begin{propriete}[Efficacité de l'écrit utilitaire]
La qualité d'un écrit utilitaire se juge à son \textbf{efficacité} : \textbf{un bon CV est celui qui décroche un entretien}. La beauté formelle, la perfection grammaticale, l'originalité du design sont \emph{secondaires} si le CV ne \guillemotleft passe \guillemotright pas le filtre du recruteur (ATS ou lecture humaine 6-10 secondes). L'efficacité = \textbf{Pertinence (mots-clés)} + \textbf{Lisibilité (structure)} + \textbf{Crédibilité (concret, chiffré, sans faute)}.
\end{propriete}

\courssub{Bilan final de la leçon}

Cette leçon a conduit l'élève de la \textbf{découverte} (textes fonctionnels/iconiques) à la \textbf{maîtrise d'un genre professionnel codifié} (le CV), en passant par l'\textbf{analyse structurelle}, l'\textbf{appropriation linguistique} (modes verbaux), la \textbf{production itérative} (brouillon $\rightarrow$ guidé $\rightarrow$ adapté) et l'\textbf{évaluation certificative}.

L'élève de Terminale Technique dispose désormais d'un \textbf{outil concret, certifié par une note, stocké dans son portfolio}, qu'il pourra faire évoluer tout au long de son parcours de technicien supérieur. La compétence \guillemotleft Rédiger un CV adapté \guillemotright est \textbf{transférable immédiatement} : recherche de stage de fin d'études, inscription en BTS/Université, premier emploi, concours administratifs.

\begin{aretenir}[Compétences validées en fin de leçon (Extrait grille MINESEC)]
\begin{itemize}
    \item \textbf{Savoir} : Structure type du CV (externe/interne), valeurs de l'infinitif/impératif dans l'écrit professionnel, critères d'adaptation à une offre.
    \item \textbf{Savoir-faire} : Analyser une offre $\rightarrow$ Produire un CV ciblé (1 page, PDF, sans faute, verbes d'action) $\rightarrow$ S'auto-évaluer sur grille.
    \item \textbf{Savoir-être} : Rigueur, sens de la communication professionnelle, autonomie, capacité d'adaptation, valorisation de son parcours technique.
\end{itemize}
\end{aretenir}

\newpage

\coursec{Activité d'intégration}

\coursec{Leçon 12 : Production d'écrits utilitaires — Rédaction d'un CV}

\courssub{Objectifs de l'activité}
\noindent
À l'issue de cet atelier d'intégration, l'élève doit être capable de :
\begin{itemize}
    \item \textbf{Analyser} une offre d'emploi réelle en identifiant le poste visé, le profil de compétences exigé (savoirs, savoir-faire, savoir-être) et les consignes de candidature formulées à l'impératif ou à l'infinitif.
    \item \textbf{Mobiliser} les connaissances sur la structure externe (mise en page, rubriques obligatoires/facultatives) et la structure interne (rédaction des entrées, chronologie inversée, verbes d'action, vocabulaire technique) du CV.
    \item \textbf{Rédiger} un CV complet, cohérent et adapté à une offre précise, en exploitant les valeurs de l'impératif (conseil, instruction) et de l'infinitif (fonction nominale, titre de rubrique, description de tâches).
    \item \textbf{Évaluer} la production d'un pair à l'aide d'une grille critériée (pertinence, lisibilité, correction linguistique, adaptation) et formuler des propositions d'amélioration argumentées.
    \item \textbf{Finaliser} une version définitive prête à être déposée, intégrant les corrections issues de la confrontation.
\end{itemize}

\courssub{Situation}
\noindent
\textbf{Contexte :} Vous êtes élèves en Terminale Technique (spécialités : Maintenance Industrielle, Secrétariat Bilingue, Génie Civil) au Lycée Technique de Douala-Bonabéri. Dans le cadre du module « Insertion Professionnelle », votre professeur de Français organise une simulation de recrutement en partenariat avec trois entreprises locales. L'objectif est de vous préparer à la recherche effective de votre premier emploi ou stage de fin d'études (Probatoire/Baccalauréat technique en vue).

\medskip
\noindent
Le professeur distribue trois offres d'emploi authentiques, reconstituées à partir d'annonces récentes parues sur le site \emph{Emploi.cm} et dans le journal \emph{Mutations}. Chaque groupe de quatre élèves tire au sort l'une des offres ci-dessous.

\medskip
\noindent
\textbf{Offre n° 1 -- Brasseries du Cameroun (Douala-Bassa)}
\begin{center}
\begin{tabularx}{\linewidth}{|l|X|}\hline
\textbf{Poste} & Technicien de Maintenance Industrielle (H/F) -- Stage professionnalisant de 12 mois, débouchant sur CDI \\ \hline
\textbf{Profil exigé} & \begin{itemize}\item Baccalauréat Technique F3 (Maintenance des Systèmes Mécaniques Automatisés) ou F2 (Électrotechnique) ;\item Maîtrise de la maintenance préventive et curative sur lignes de conditionnement (Krones, Sidel) ;\item Lecture de plans hydrauliques, pneumatiques, schémas électriques (Autocad Electrical) ;\item Notions de GMAO (SAP PM / Maximo) ;\item Anglais technique lu et écrit (niveau B1 minimum) ;\item Permis B, disponibilité immédiate pour travail posté (3x8).\end{itemize} \\ \hline
\textbf{Compétences transverses} & Rigueur, réactivité, esprit d'équipe, respect strict des consignes HSE (Hygiène Sécurité Environnement). \\ \hline
\textbf{Constitution du dossier} & \textbf{Envoyer} (infinitif d'injonction) CV actualisé + lettre de motivation + copies certifiées conformes des diplômes + extrait d'acte de naissance de moins de 3 mois \textbf{à l'adresse} : \texttt{recrutement@brasseries-cm.com} \textbf{avant le 15/01/2026}. Objet du mail : \texttt{STAGE\_MAINT\_2026\_NOM\_PRENOM}. \\ \hline
\end{tabularx}
\end{center}

\medskip
\noindent
\textbf{Offre n° 2 -- Afriland First Bank (Yaoundé, Agence Centre Administratif)}
\begin{center}
\begin{tabularx}{\linewidth}{|l|X|}\hline
\textbf{Poste} & Secrétaire Bilingue Accueil-Clientèle (H/F) -- CDD 6 mois renouvelable \\ \hline
\textbf{Profil exigé} & \begin{itemize}\item Baccalauréat Technique G2 (Secrétariat de Direction) ou G3 (Action Commerciale) ;\item Maîtrise parfaite du français et de l'anglais (écrit/oral) -- test de niveau obligatoire ;\item Maîtrise de la suite Office 365 (Word, Excel avancé : TCD, RECHERCHEV ; PowerPoint, Outlook) ;\item Techniques de secrétariat : gestion d'agenda, filtrage d'appels, rédaction de comptes rendus, archivage numérique (SharePoint) ;\item Expérience de 6 mois minimum en milieu bancaire ou administration (stage validé accepté).\end{itemize} \\ \hline
\textbf{Compétences transverses} & Discrétion professionnelle (secret bancaire), sens du service client, gestion du stress, présentation soignée. \\ \hline
\textbf{Constitution du dossier} & \textbf{Déposer} (impératif) le dossier physique à l'accueil RH ou \textbf{transmettre} (infinitif comme consigne) par voie numérique : CV + LM + photocopie CNI + attestation de stage + casier judiciaire bulletin n°3 \textbf{au plus tard le 20/01/2026}. \textbf{Préciser} (impératif) la référence : \texttt{REF/SEC/BIL/0126}. \\ \hline
\end{tabularx}
\end{center}

\medskip
\noindent
\textbf{Offre n° 3 -- Entreprise MAGNA BTP (Bafoussam, Chantier Autoroute Dschang-Bafoussam)}
\begin{center}
\begin{tabularx}{\linewidth}{|l|X|}\hline
\textbf{Poste} & Conducteur de Travaux Junior (H/F) -- CDI \\ \hline
\textbf{Profil exigé} & \begin{itemize}\item Baccalauréat Technique F4 (Génie Civil - Travaux Publics) ou BT Bâtiment ;\item Lecture de plans d'exécution (DAO : AutoCAD, Covadis) ;\item Suivi de chantier : métrés, quantification, planning (MS Project), gestion des approvisionnements ;\item Connaissance des normes CCTG (Cahier des Clauses Techniques Générales) routières ;\item Encadrement d'une équipe de 10 à 15 ouvriers (maçons, ferrailleurs, manutentionnaires) ;\item Mobilité nationale (chantiers itinérants), permis C souhaité.\end{itemize} \\ \hline
\textbf{Compétences transverses} & Leadership, sens de l'organisation, reporting quotidien au chef de projet, respect des délais et budget. \\ \hline
\textbf{Constitution du dossier} & \textbf{Adresser} (impératif) candidature (CV détaillé + LM + copies diplômes + permis) \textbf{par courrier recommandé avec AR} à : \texttt{MAGNA BTP, BP 145 Bafoussam} \textbf{ou par mail} : \texttt{rh@magna-btp.cm}. \textbf{Date limite : 25/01/2026}. \textbf{Joindre} (impératif) obligatoirement une attestation de visite médicale de moins de 6 mois. \\ \hline
\end{tabularx}
\end{center}

\medskip
\noindent
\textbf{Contrainte commune :} Le candidat fictif dont vous allez rédiger le CV est un jeune diplômé de 2024 (âgé de 19-20 ans), sans expérience professionnelle longue, mais ayant effectué des stages significatifs et des projets scolaires techniques pertinents. Vous devez inventer son identité (nom, prénom, contact réaliste au Cameroun) et son parcours pour qu'il \textbf{colle exactement} au profil exigé.

\courssub{Consignes}
\noindent
Travail en groupes de 4 élèves. Durée : 2 séances (3h). Le professeur circule pour valider les étapes.

\begin{enumerate}
    \item \textbf{Analyse de l'offre (30 min) :} Lisez l'offre attribuée. Remplissez la fiche d'analyse ci-dessous :
    \begin{itemize}
        \item Poste visé / Entreprise / Localisation.
        \item 5 compétences techniques « dures » (hard skills) explicites.
        \item 3 compétences transverses (soft skills) explicites ou implicites.
        \item Relevez \textbf{tous les verbes de consigne} (impératif / infinitif) utilisés dans la rubrique « Constitution du dossier ». Classez-les dans un tableau à deux colonnes et expliquez la nuance de valeur (ordre, conseil, procédure administrative).
        \item Identifiez les mots-clés à faire apparaître dans le CV pour passer un tri automatisé (ATS).
    \end{itemize}
    \item \textbf{Rédaction du CV candidat (60 min) :} Construisez le CV du candidat idéal.
    \begin{itemize}
        \item \textbf{Structure externe :} Choisissez un modèle (chronologique inversé, fonctionnel ou mixte) justifié. Respectez les normes : photo professionnelle (optionnelle mais recommandée au Cameroun), en-tête complet, rubriques ordonnées, mise en page aérée (police 11-12, interligne 1,15, marges 2 cm), 1 page max.
        \item \textbf{Structure interne :} Rédigez chaque entrée (Formation, Stages/Expériences, Projets techniques, Compétences techniques, Langues, Centres d'intérêt) en utilisant des \textbf{verbes d'action à l'infinitif} (ex. : \emph{Réaliser la maintenance préventive sur...}, \emph{Gérer l'agenda de la direction...}, \emph{Superviser le ferraillage de...}). Quantifiez les résultats (durées, effectifs, volumes, taux de réussite).
        \item Adaptez le vocabulaire technique à la spécialité (maintenance, secrétariat, BTP).
    \end{itemize}
    \item \textbf{Évaluation par les pairs et amélioration (45 min) :} Échangez votre production avec un groupe voisin (offre différente). Utilisez la \textbf{Grille d'évaluation CV} fournie (critères : Pertinence/Adaptation -- 4 pts ; Structure externe -- 3 pts ; Structure interne / Contenu -- 5 pts ; Langue / Modes verbaux -- 3 pts ; Orthographe/Typographie -- 3 pts ; Total 18 pts). Annotez le CV : surlignez les forces (vert), les faiblesses (rouge), écrivez 3 conseils précis d'amélioration au dos.
    \item \textbf{Version définitive (15 min) :} Récupérez votre CV annoté. Intégrez les améliorations pertinentes. Remettez la version finale (papier + fichier .pdf sur clé USB du professeur) pour affichage.
    \item \textbf{Compte rendu oral (10 min) :} Un porte-parole par groupe présente : l'offre choisie, la stratégie d'adaptation (pourquoi ce modèle ? quels mots-clés ?), la principale difficulté rencontrée, la correction apportée suite à l'évaluation.
\end{enumerate}

\courssub{Ressources à mobiliser}
\begin{aretenir}[Rappel : Structure Externe du CV]
\textbf{Rubriques obligatoires :} En-tête (Identité, Coordonnées, Titre/Objectif), Formation (diplômes, mentions, établissement, dates), Expériences professionnelles / Stages (entreprise, dates, poste, missions, résultats), Compétences techniques (logiciels, machines, normes, langues), Centres d'intérêt (liens avec le poste si possible).
\textbf{Rubriques facultatives :} Photo, Projets scolaires / PFE, Certifications / Permis, Références.
\textbf{Mise en page :} Police sans empattement (Calibri, Arial, Roboto), alignement gauche, puces cohérentes (● ou --), gras sur les intitulés de poste/école, italique sur les lieux/dates. Fichier final : \textbf{PDF} (nom : \texttt{NOM\_Prenom\_CV\_Poste.pdf}).
\end{aretenir}

\begin{aretenir}[Rappel : Structure Interne \& Verbes d'Action]
\begin{itemize}
    \item \textbf{Chronologie inversée :} Du plus récent au plus ancien.
    \item \textbf{Formule standard d'une entrée Expérience :} \emph{Intitulé du poste} -- \emph{Entreprise, Ville} -- \emph{Dates} \newline $\bullet$ \textbf{Verbe infinitif} + Objet + Contexte + \textbf{Résultat chiffré}.
    \item \textbf{Exemples :} \emph{Diagnostiquer} les pannes hydrauliques sur presse à injecter (réduction temps d'arrêt 15\,\%) ; \emph{Rédiger} les comptes rendus de réunions du comité de direction (15 PV/an) ; \emph{Contrôler} la conformité des armatures selon plans d'exécution (500 ml linéaires/semaine).
    \item \textbf{Mots-clés (ATS) :} Reprendre les termes exacts de l'offre (ex. : \emph{GMAO SAP PM}, \emph{RECHERCHEV}, \emph{CCTG routières}).
\end{itemize}
\end{aretenir}

\begin{propriete}[Valeurs de l'Impératif et de l'Infinitif dans les offres d'emploi]
\begin{itemize}
    \item \textbf{Impératif (2\textsuperscript{e} pers. sg/pl ou 1\textsuperscript{re} pl) :} Valeur \textbf{injonctive} (ordre, consigne stricte, procédure obligatoire). Ex. : \emph{Envoyez}, \emph{Déposez}, \emph{Adressez}, \emph{Joignez}, \emph{Précisez}. Marque l'autorité de l'employeur / l'obligation administrative.
    \item \textbf{Infinitif (souvent prépositionnel : \emph{à} + inf. ou seul en titre) :} Valeur \textbf{nominale / procédurale}. Ex. : \emph{Transmettre par voie numérique}, \emph{Préciser la référence}. Désigne l'action comme une étape de la procédure, sans s'adresser directement au sujet (impersonnel). Utilisé aussi pour les \textbf{titres de rubriques} (Formation, Expériences) et la \textbf{description des missions} dans le CV (style télégraphique, impersonnel, concis).
    \item \textbf{Infinitif d'injonction (forme infinitive, valeur impérative) :} Fréquent dans les consignes administratives écrites. Ex. : \emph{Envoyer} (au lieu d'\emph{Envoyez}), \emph{Joindre} (au lieu de \emph{Joignez}). Équivalent fonctionnel à l'impératif, mais plus neutre / impersonnel.
\end{itemize}
\end{propriete}

\begin{methode}[Adapter son CV à une offre (3 étapes)]
\begin{enumerate}
    \item \textbf{Décortiquer l'offre :} Surligner \texttt{\textcolor{popblue}{compétences techniques}}, \texttt{\textcolor{poporange}{soft skills}}, \texttt{\textcolor{popgreen}{mots-clés}}, \texttt{\textcolor{poppink}{consignes}}.
    \item \textbf{Faire le matching :} Pour chaque exigence, cocher : \emph{Acquis (stage/projet)} / \emph{Partiel (à valoriser)} / \emph{Manquant (à former / ne pas mentir)}.
    \item \textbf{Réécrire le CV :} Ordonner les rubriques selon l'offre (Compétences techniques en 1\textsuperscript{er} si technique pur), injecter les mots-clés, reformuler les missions avec les verbes de l'offre, supprimer le superflu.
\end{enumerate}
\end{methode}

\begin{attention}[Pièges fréquents à éviter]
\begin{itemize}
    \item \textbf{Mentir sur les compétences :} Un recruteur technique teste (QCM, mise en situation). L'honnêteté « Partiel (à valoriser) » est préférable.
    \item \textbf{CV « passe-partout » :} Un même CV pour toutes les offres = rejet quasi certain. L'adaptation (mots-clés, ordre rubriques) est obligatoire.
    \item \textbf{Négliger la typographie :} Espace insécable avant \%, ; : ! ? ; guillemets français « » ; tirets demi-cadratin -- pour les puces. Fichier \texttt{.pdf} non modifiable, nommé \texttt{NOM\_Prenom\_CV\_Poste.pdf}.
    \item \textbf{Oublier les pièces jointes :} L'offre liste des pièces (acte de naissance, casier, visite médicale). Le mail/lettre doit lister explicitement les PJ.
\end{itemize}
\end{attention}

\courssub{Démarche et corrigé commenté}
\begin{exoresolu}[Atelier « Décroche ton premier emploi » -- Exemple traité : Offre n°1 (Brasseries du Cameroun)]
\textbf{Énoncé :} Analyse de l'offre 1 (Brasseries du Cameroun) et rédaction du CV adapté pour un profil F3 Maintenance Industrielle.
\tcblower
\textbf{Étape 1 -- Analyse de l'offre (Groupe A)}
\begin{center}
\begin{tabularx}{\linewidth}{|l|X|}\hline
\textbf{Poste / Entreprise} & Technicien Maintenance Industrielle -- Brasseries du Cameroun (Douala-Bassa) \\ \hline
\textbf{Hard Skills (5)} & 1. Maintenance préventive/curative lignes conditionnement (Krones, Sidel) ; 2. Lecture plans hydrauliques/pneumatiques/électriques (Autocad Elec) ; 3. GMAO (SAP PM/Maximo) ; 4. Anglais technique B1 ; 5. Travail posté 3x8 / Permis B. \\ \hline
\textbf{Soft Skills (3)} & 1. Rigueur (HSE) ; 2. Réactivité (pannes) ; 3. Esprit d'équipe (équipe maintenance/production). \\ \hline
\textbf{Verbes de consigne (Offre 1)} & \begin{tabular}{ll}
\textbf{Forme / Mode} & \textbf{Verbe / Valeur} \\
Infinitif d'injonction & \emph{Envoyer} (valeur impérative, consigne stricte) \\
-- & (Aucun infinitif procédural pur dans cette offre) \\
-- & (Aucun impératif morphologique en \emph{-ez}) \\
\end{tabular} \\ \hline
\textbf{Analyse modes} & L'offre 1 utilise la forme \textbf{infinitif d'injonction} (\emph{Envoyer}) pour la consigne principale. Cette forme, bien que morphologiquement infinitif, a une \textbf{valeur injonctive} équivalente à l'impératif. Elle marque l'obligation stricte de la procédure. L'infinitif procédural (\emph{à transmettre}, \emph{à joindre}) n'apparaît pas ici mais sera utilisé \textbf{dans le CV} pour les missions. \\ \hline
\textbf{Mots-clés ATS} & Maintenance préventive, Krones, Sidel, Hydraulique, Pneumatique, Autocad Electrical, GMAO, SAP PM, Maximo, Anglais technique, HSE, 3x8. \\ \hline
\end{tabularx}
\end{center}

\textbf{Étape 2 -- Construction du candidat fictif \& Rédaction du CV}
\noindent
\textbf{Identité inventée :} \textsc{Ngono} Jean-Pierre, né le 12/03/2005 à Douala (Cameroun), Célibataire, Permis B. \\
\textbf{Contact :} +237 6 90 12 34 56 -- \texttt{jp.ngono@email.com} -- Quartier Bonabéri, Douala. \\
\textbf{Titre CV :} \textbf{Technicien Maintenance Industrielle (F3) -- Spécialisation Lignes de Conditionnement}

\medskip
\noindent
\textbf{Choix du modèle :} \textbf{Chronologique inversé mixte} (Compétences techniques mises en avant avant Expériences car jeune diplômé, mais stages significatifs).

\medskip
\noindent
\textbf{Contenu du CV (Structure Interne -- Extraits commentés) :}

\begin{itemize}
    \item \textbf{FORMATION}
    \begin{itemize}
        \item 2024 : \textbf{Baccalauréat Technique F3 -- Maintenance des Systèmes Mécaniques Automatisés} -- Mention Assez Bien (13,2/20). Lycée Technique de Douala-Bonabéri. \emph{Projet de fin d'études :} \emph{Concevoir} et \emph{programmer} une station de graissage automatique sur automate Siemens S7-1200 (Note : 16/20).
        \item 2022 : \textbf{Probatoire Technique F3} -- Mention Bien. Lycée Technique de Douala-Bonabéri.
    \end{itemize}
    
    \item \textbf{COMPÉTENCES TECHNIQUES (Mots-clés injectés)}
    \begin{itemize}
        \item \textbf{Maintenance :} Préventive (planification, gammes), Curative (diagnostic pannes mécaniques/hydrauliques/électriques). Machines : \textbf{Krones (remplisseuse/boucheuse)}, \textbf{Sidel (souffleuse)} -- \emph{Connaissance théorique approfondie + pratique stage}.
        \item \textbf{Outils CAO/DAO :} \textbf{Autocad Electrical} (schémas), SolidWorks (pièces), Lecture plans ISO (hydraulique/pneumatique).
        \item \textbf{GMAO :} \textbf{SAP PM} (création ordres, clôture, historique) -- \emph{Utilisation quotidienne en stage}.
        \item \textbf{Automatismes :} Automates Siemens (S7-300/1200/1500), TIA Portal (bases), Capteurs/Actionneurs.
        \item \textbf{HSE :} Port EPI, Consignation (LOTO), Analyse de risques (AMDEC basique), Permis de feu.
    \end{itemize}
    
    \item \textbf{EXPÉRIENCES PROFESSIONNELLES / STAGES (Chronologie inversée, Verbes infinitif + Résultats)}
    \begin{itemize}
        \item \textbf{Juin -- Août 2024 (3 mois) : Stagiaire Technicien Maintenance -- Brasseries du Cameroun, Douala-Bassa.}
        \begin{itemize}
            \item \emph{Effectuer} la maintenance préventive niveau 1 sur ligne Krones (remplisseuse 50 000 b/h) : graissage, contrôle jeux, changement filtres -- \textbf{Résultat :} 0 panne majeure sur la période.
            \item \emph{Assister} le technicien senior dans le diagnostic d'une panne récurrente sur boucheuse (casse bouteilles) : \emph{analyser} vibrométrie, \emph{remplacer} came usée -- \textbf{Résultat :} Taux de casse réduit de 1,2\,\% à 0,3\,\%.
            \item \emph{Renseigner} la GMAO \textbf{SAP PM} : création notifications, ordres de travail, feedbacks (15 interventions/semaine).
            \item \emph{Appliquer} strictement les consignes \textbf{HSE} : consignation LOTO avant toute intervention, port EPI complet.
        \end{itemize}
        \item \textbf{Janv -- Mars 2023 (2 mois) : Stagiaire Maintenance -- SOCAPALM (Société Camerounaise de Palmeraies), Dizangué.}
        \begin{itemize}
            \item \emph{Réaliser} la maintenance curative sur convoyeurs à bande et broyeurs (hydraulique, pneumatique, mécanique).
            \item \emph{Lire} et \emph{interpréter} des plans hydrauliques/pneumatiques pour localiser fuites et blocages.
            \item \emph{Participer} à la révision annuelle du broyeur principal (démontage/remontage roulements, alignement laser).
        \end{itemize}
    \end{itemize}
    
    \item \textbf{PROJETS TECHNIQUES (Valorisation PFE / Travaux pratiques)}
    \begin{itemize}
        \item \textbf{PFE 2024 :} \emph{Conception d'un système de graissage centralisé piloté par automate Siemens S7-1200.} \emph{Programmer} la logique (Grafcet/SCL), \emph{câbler} l'armoire, \emph{tester} en simulation. \textbf{Compétences :} Automatisme, Électrotechnique, Gestion de projet.
        \item \textbf{TP Intégration 2023 :} \emph{Monter} et \emph{régler} un circuit pneumatique séquentiel (vérins, distributeurs, capteurs) selon cahier des charges.
    \end{itemize}
    
    \item \textbf{LANGUES}
    \begin{itemize}
        \item Français : Langue maternelle.
        \item Anglais : Niveau \textbf{B1 (TOEIC 650 -- 2023)} -- Lecture plans techniques, notices constructeurs (Krones/Sidel), rédaction rapports d'intervention simples.
    \end{itemize}
    
    \item \textbf{CENTRES D'INTÉRÊT}
    \begin{itemize}
        \item Football (club local) -- Esprit d'équipe, endurance.
        \item Bricolage automobile / Mécanique de précision -- Déxtérité, curiosité technique.
        \item Bénévolat : Maintenance matériel informatique école (club robotique).
    \end{itemize}
\end{itemize}

\textbf{Étape 3 -- Évaluation par les pairs (Groupe B évalue Groupe A) -- Extrait de la grille remplie}
\begin{center}
\begin{tabularx}{\linewidth}{|l|X|c|}\hline
\textbf{Critère} & \textbf{Observations / Commentaires} & \textbf{Note / 18} \\ \hline
Pertinence / Adaptation (4) & Profil F3 parfait. Mots-clés Krones, Sidel, SAP PM, Autocad Elec, HSE, 3x8, Anglais B1 \textbf{tous présents}. Stage Brasseries mis en 1er (pertinent). Titre CV ciblé. & 4/4 \\ \hline
Structure Externe (3) & 1 page, aéré, police Calibri 11, photo pro (simulée), en-tête complet, rubriques claires, puces cohérentes. PDF nommé correctement. & 3/3 \\ \hline
Structure Interne / Contenu (5) & Verbes infinitif systématiques (\emph{Effectuer, Assister, Analyser, Remplacer, Renseigner, Appliquer, Réaliser, Lire, Participer, Concevoir, Programmer, Câbler}). Résultats \textbf{chiffrés} (taux casse, nb interventions). Chronologie inversée respectée. & 5/5 \\ \hline
Langue / Modes (3) & Zéro faute d'accord. Infinitif d'injonction respecté dans la lettre d'accompagnement (objet mail). Style télégraphique (infinitif) parfait dans le CV. & 3/3 \\ \hline
Orthographe / Typographie (3) & 1 coquille : « \emph{remplacér} » au lieu de « \emph{remplacer} » (ligne stage Brasseries). 1 espace insécable manquante avant \% (0,3\,\%). & 2/3 \\ \hline
\textbf{TOTAL} & \textbf{Excellent CV, prêt à envoyer après correction mineure.} & \textbf{17/18} \\ \hline
\end{tabularx}
\end{center}

\textbf{Conseils d'amélioration formulés par Groupe B :}
\begin{enumerate}
    \item Corriger la coquille « remplacér $\to$ remplacer ».
    \item Ajouter espace insécable avant symbole \% (0,3\,\%).
    \item Préciser la durée du stage Brasseries dans l'en-tête de l'entrée : « Juin -- Août 2024 (3 mois) » $\to$ déjà fait, valider.
    \item \textbf{Force majeure :} Mentionner explicitement « Disponibilité 3x8 » dans la rubrique « Disponibilités / Mobilité » ou en bas de CV.
\end{enumerate}

\textbf{Étape 4 -- Version Définitive (Groupe A intègre)}
\noindent
Corrections apportées. Ajout d'une ligne « \textbf{Disponibilité :} Immédiate -- Travail posté 3x8 accepté -- Mobilité nationale : Oui » sous l'en-tête. Fichier final généré : \texttt{NGONO\_Jean-Pierre\_CV\_TechnicienMaintenance\_Brasseries.pdf}.

\textbf{Étape 5 -- Compte rendu oral (Porte-parole Groupe A)}
\noindent
\emph{« Nous avons choisi l'offre Brasseries. Stratégie : mettre le stage Brasseries en tout premier (expérience la plus parlante), créer une rubrique "Compétences Techniques" dense avant "Expériences" pour capter les mots-clés ATS (Krones, Sidel, SAP PM). Difficulté : quantifier les résultats du stage SOCAPALM (moins encadré). Solution : estimation réaliste "5 interventions/jour". Correction apportée : coquille "remplacér" et ajout dispo 3x8. Le CV est maintenant "taillé sur mesure". »}
\end{exoresolu}

\courssub{Bilan de l'activité}
\noindent
Cet atelier d'intégration a permis de vérifier l'acquisition effective des compétences visées par la leçon 12 dans une situation professionnalisante authentique :

\begin{enumerate}
    \item \textbf{Transfert des savoirs :} Les élèves ont démontré leur capacité à passer de l'analyse théorique (structure externe/interne, modes verbaux) à la production concrète. Le choix du modèle de CV (chronologique inversé mixte) justifié par le profil « jeune diplômé avec stages forts » montre une compréhension fine de la \textbf{structure externe} comme outil de stratégie de communication.
    \item \textbf{Maîtrise de la structure interne :} L'usage systématique de l'\textbf{infinitif de narration professionnelle} (\emph{Effectuer, Diagnostiquer, Programmer, Renseigner}) pour décrire les missions, couplé à la \textbf{quantification des résultats} (taux de casse, nombre d'interventions, volumes), valide l'acquisition de la rédaction « télégraphique » efficace attendue au Bac technique.
    \item \textbf{Compétence linguistique ciblée :} L'exercice de relevé des verbes de consigne dans les offres (Infinitif d'injonction : \emph{Envoyer} ; Impératif morphologique : \emph{Déposez, Adressez, Joignez, Précisez} ; Infinitif procédural : \emph{transmettre}) a ancré la distinction entre \textbf{valeur injonctive} (relation hiérarchique candidat/recruteur) et \textbf{valeur nominale/procédurale} (description de processus). Les élèves ont réinvesti cette distinction : impératif/infinitive d'injonction dans la lettre de motivation (\emph{Je vous prie d'agréer...}), infinitif dans le CV.
    \item \textbf{Adaptation contextuelle (Cameroun) :} L'intégration de réalités locales (Brasseries du Cameroun, Afriland First Bank, MAGNA BTP, mentions Probatoire/Bac Tech F3/G2/F4, format dossier physique + numérique, photo, extrait acte de naissance, casier judiciaire, visite médicale) ancre l'apprentissage dans le vécu des élèves et les prépare aux exigences réelles du marché de l'emploi camerounais.
    \item \textbf{Évaluation formative par les pairs :} La grille critériée (18 pts) a objectivé l'évaluation. L'échange entre groupes (offres différentes) a forcé l'objectivité : on évalue la \textbf{qualité de l'adaptation} à l'offre, pas la spécialité technique. Les corrections mineures (coquilles, typographie) relevées prouvent l'attention au détail, critère éliminatoire en recrutement.
    \item \textbf{Remédiation immédiate :} La version définitive corrigée en séance, validée par le professeur et affichée, clôture le cycle « Production -- Confrontation -- Amélioration » préconisé par le programme. Les CV affichés servent de modèles de référence pour la classe entière.
\end{enumerate}

\medskip
\noindent
\cle{Compétence transversale} : Au-delà du français technique, cet atelier développe l'\textbf{employabilité} : savoir se vendre par écrit, décoder une offre, respecter une procédure administrative, travailler en équipe, accepter la critique constructive. C'est la finalité ultime de l'enseignement technique au Cameroun.

\newpage

\coursec{Méthodes et exercices résolus}

\coursec{Production d'écrits utilitaires : rédaction d'un CV}

\courssub{Analyser une offre d'emploi et un texte iconique (annonce, affiche)}

\begin{methode}[Analyser un texte fonctionnel de recrutement (offre d'emploi)]
\textbf{Objectif :} Identifier les attentes du recruteur (compétences, qualités, diplômes) pour cibler son CV.
\begin{enumerate}
    \item \textbf{Lire le texte en entier} sans surligner pour saisir le contexte (entreprise, poste, secteur).
    \item \textbf{Repérer les mots-clés} : intitulé du poste, verbes d'action (\cle{concevoir}, \cle{gérer}, \cle{réaliser}, \cle{animer}), compétences techniques (\cle{hard skills}) et transversales (\cle{soft skills}), diplômes requis, expérience demandée.
    \item \textbf{Classifier les exigences} dans un tableau à deux colonnes : \og Obligatoires \fg (critères éliminatoires) et \og Souhaitées \fg (critères différenciants).
    \item \textbf{Analyser le texte iconique associé} (logo, mise en page, photo, slogan) : quelle image l'entreprise renvoie-t-elle ? (dynamique, sérieuse, innovante, sociale). Adapter le ton de son CV en conséquence.
    \item \textbf{Synthétiser} : rédiger une phrase résumant le \cle{profil idéal} recherché.
\end{enumerate}
\end{methode}

\begin{exoresolu}[Analyse d'une offre pour un poste de Technicien en Maintenance Industrielle]
\textbf{Énoncé :} Voici une offre d'emploi. Appliquez la méthode pour en extraire les éléments clés à faire apparaître sur votre CV.

\emph{Offre :} \og Entreprise agroalimentaire leader au Cameroun recherche un Technicien de Maintenance H/F. Bac+2/3 (BTS/DUT Maintenance, Électrotechnique, Mécanique). Expérience 2 ans minimum sur ligne de production automatisée. Maîtrise impérative : GMAO (SAP PM), lecture plans électriques/hydrauliques, variateurs de vitesse, automatismes Siemens/Schneider. Qualités : rigueur, réactivité, esprit d'équipe, respect normes HSE. Permis B. Poste en 3x8. \fg

\tcblower
\textbf{Solution détaillée :}

\noindent \textbf{Étape 1 : Contexte.} Secteur agroalimentaire (contraintes HSE fortes, production continue). Poste : Technicien Maintenance (opérationnel, terrain).

\noindent \textbf{Étape 2 : Mots-clés repérés.}
\begin{itemize}
    \item \textbf{Diplôme :} Bac+2/3 (BTS/DUT Maintenance, Électrotechnique, Mécanique).
    \item \textbf{Expérience :} 2 ans min, \emph{ligne de production automatisée}, \emph{3x8}.
    \item \textbf{Hard Skills (Techniques) :} GMAO (SAP PM), lecture plans (électrique/hydraulique), variateurs vitesse, automatismes (Siemens, Schneider).
    \item \textbf{Soft Skills :} Rigueur, réactivité, esprit d'équipe, respect normes HSE.
    \item \textbf{Autre :} Permis B.
\end{itemize}

\noindent \textbf{Étape 3 : Tableau de classification.}

\begin{center}
\begin{tabularx}{\linewidth}{|X|X|}
\hline
\rowcolor{popblueL}
\textbf{Critères Obligatoires (Eliminatoires)} & \textbf{Critères Souhaités (Différenciants)} \\
\hline
- Diplôme BTS/DUT Maintenance ou équivalent & - Connaissance spécifique SAP PM (autre GMAO acceptée ?) \\
- Expérience 2 ans maintenance industrielle & - Maîtrise marques Siemens/Schneider (autres marques = adaptable) \\
- Lecture plans électriques/hydrauliques & - Expérience agroalimentaire (HSE spécifique) \\
- Variateurs de vitesse, automatismes de base & - Disponibilité immédiate pour 3x8 \\
- Permis B & \\
- Respect normes HSE (réglementaire) & \\
\hline
\end{tabularx}
\end{center}

\noindent \textbf{Étape 4 : Texte iconique (imaginé).} Logo vert/bleu (agro/nature), usine en arrière-plan, casque/tenue sécurité. Image : \emph{sécurité, hygiène, technologie, continuité}. $\rightarrow$ CV doit respirer la \cle{fiabilité} et la \cle{sécurité}.

\noindent \textbf{Étape 5 : Synthèse du profil idéal.}
\og Technicien maintenance polyvalent (électro/méca), autonome sur machines automatisées, expérimenté en environnement industriel contraint (HSE, 3x8), capable de gérer la traçabilité via GMAO. \fg

\noindent \textbf{Conclusion :} Votre CV doit mettre en \textbf{gras} (ou en premier) : BTS Maintenance, 2 ans+ expé ligne auto, GMAO, lecture plans, HSE, Permis B, dispo 3x8.
\end{exoresolu}

\courssub{Maîtriser la structure externe du CV (Mise en page, Lisibilité, Normes)}

\begin{methode}[Structurer la mise en page d'un CV professionnel]
\textbf{Objectif :} Produire un document lisible en 30 secondes (règle du \cle{scan rapide}), aéré, sans faute, au format standard (PDF, 1 page idéalement).
\begin{enumerate}
    \item \textbf{Format \& Fichier :} Format \cle{A4}, marges 2 cm min. Police unique sans empattement (\emph{Calibri, Arial, Roboto, Open Sans}) taille 10--11 pour le corps, 13--14 pour les titres. Enregistrer en \textbf{PDF} (nom : \texttt{Nom\_Prenom\_CV\_Poste.pdf}).
    \item \textbf{En-tête (Identité) :} Nom \textsc{PRÉNOM} (majuscules pour le nom). Coordonnées alignées : Téléphone (+237...), Email pro (prenom.nom@email.com), Ville/Quartier (ex: Douala, Bonapriso), LinkedIn/Profil pro (optionnel). \textbf{Pas de photo obligatoire} (sauf demande), si photo : fond neutre, tenue pro, format identité.
    \item \textbf{Organisation en blocs visuels distincts} (séparés par filets fins ou espaces blancs) : 1. \cle{Profil / Accroche} (3 lignes max), 2. \cle{Expériences} (ordre antéchronologique), 3. \cle{Formation}, 4. \cle{Compétences Techniques}, 5. \cle{Langues / Certifications}, 6. \cle{Centres d'intérêt} (ciblés).
    \item \textbf{Codes visuels constants :} Dates à gauche (MM/AAAA -- MM/AAAA), Lieu à droite. Titres de poste en \textbf{gras}, Entreprise en \textit{italique}. Puces \textbullet\ ou \textendash\ pour les missions (max 4--5 par poste).
    \item \textbf{Vérification finale :} Alignements parfaits (tabulations), pas de veuves/orphelines (titre seul en bas de page), cohérence des temps (passé pour ancien, présent pour actuel), orthographe (correcteur + relecture humaine).
\end{enumerate}
\end{methode}

\begin{exoresolu}[Correction de la structure externe d'un CV défectueux]
\textbf{Énoncé :} Un élève remet le brouillon suivant. Identifiez 5 défauts majeurs de structure externe et proposez la version corrigée (description textuelle de la mise en page).

\emph{Brouillon :} Police Comic Sans MS 12. Marge 1 cm. En-tête : "jean dupont" minuscules, email "jeandupon99@yahoo.fr", pas de tel. Photo de vacances (plage). Bloc unique dense sans séparation : "EXPERIENCES : 2020-2022 Stagiaire maintenance SOCIETE X Douala j'ai réparé des machines, fait de la soudure, géré le stock. FORMATION : 2019 BAC F2. 2021 BTS Maintenance. COMPETENCES : Word, Excel, Soudure, Mécanique, Anglais. LOISIRS : Foot, danse, cinéma."

\tcblower
\textbf{Solution détaillée :}

\noindent \textbf{Défaut 1 : Police inadaptée.} Comic Sans MS $\rightarrow$ non professionnel. \textbf{Correction :} \emph{Calibri} ou \emph{Arial} taille 11.

\noindent \textbf{Défaut 2 : Marges trop étroites (1 cm).} $\rightarrow$ effet "brouillon", imprimerie risquée. \textbf{Correction :} Marges 2 cm (haut/bas/gauche/droite).

\noindent \textbf{Défaut 3 : En-tête non standard.} Nom en minuscules, email non pro, pas de téléphone. \textbf{Correction :} \textsc{DUPO} Jean | +237 6XX XX XX XX | jean.dupont@email.pro | Douala, Akwa | linkedin.com/in/jeandupont.

\noindent \textbf{Défaut 4 : Photo inappropriée.} Photo plage $\rightarrow$ élimination directe. \textbf{Correction :} Supprimer la photo (non obligatoire au Cameroun sauf demande) ou photo fond neutre, chemise, sourire léger.

\noindent \textbf{Défaut 5 : Structure "bloc unique" illisible.} Pas de hiérarchie visuelle, phrases complètes au lieu de puces, pas d'ordre antéchronologique clair, mélange missions/résultats. \textbf{Correction :} Structure en 5 blocs distincts avec titres colorés (bleu marine) ou filets.

\noindent \textbf{Proposition de structure corrigée (squelette) :}

\begin{center}
\begin{tabularx}{\linewidth}{|X|}
\hline
\rowcolor{popblueL} \textsc{DUPO} \textbf{Jean} \hfill +237 6XX XX XX XX \hfill jean.dupont@email.pro \\
\rowcolor{popblueL} Douala, Akwa \hfill linkedin.com/in/jeandupont \hfill 24 ans, Permis B \\
\hline
\textbf{PROFIL} \hfill \textit{Technicien Maintenance Junior} \\
Technicien BTS Maintenance (Option Systèmes Automatisés), 2 ans d'alternance sur ligne agroalimentaire. Autonome en maintenance préventive/curative, lecture plans, GMAO. Rigueur HSE. Recherche poste 3x8. \\
\hline
\textbf{EXPÉRIENCES PROFESSIONNELLES} \\
\textbf{09/2021 -- 06/2023} \hfill \textbf{Technicien Maintenance (Alternance)} \\
\textit{SOCIÉTÉ X, Douala} \\
\textbullet\ Maintenance préventive niveau 1/2 sur ligne conditionneuse (Krones/Sidel). \\
\textbullet\ Dépannage curatif : variateurs, capteurs, automates Siemens S7-1200. \\
\textbullet\ Gestion pièces de rechange sous SAP PM (création notifications, bons de travail). \\
\textbullet\ Rédaction modes opératoires sécurité (Cadenassage/Consignation). \\
\hline
\textbf{FORMATION} \\
\textbf{2021 -- 2023} \hfill \textbf{BTS Maintenance des Systèmes (Option MSA)} \hfill Lycée Technique de Douala \\
\textbf{2019} \hfill \textbf{Baccalauréat F2 (Mécanique)} \hfill Lycée de New-Bell \\
\hline
\textbf{COMPÉTENCES TECHNIQUES} \\
\textbullet\ \textbf{Automatismes :} Siemens (S7-300/1200, TIA Portal), Schneider (Twido, M241). \\
\textbullet\ \textbf{Électrotechnique :} Variateurs (Altivar, Sinamics), Lecture plans IEC, Câblage armoires. \\
\textbullet\ \textbf{GMAO / IT :} SAP PM (Utilisateur avancé), Pack Office, AutoCAD Electrical (notions). \\
\textbullet\ \textbf{Méthodes :} Maintenance Préventive (Planification, Gammes), 5S, AMDEC (initiation). \\
\hline
\textbf{LANGUES \& CERTIFICATIONS} \\
\textbullet\ Français : Langue maternelle \quad Anglais : Niveau B1 (TOEIC 650) \\
\textbullet\ Habilitation Électrique BR/BC (Valide) \quad CACES R489 Cat 1B (Chariot) \\
\hline
\textbf{CENTRES D'INTÉRÊT} \\
\textbullet\ Football (Club entreprise, Capitaine) : Esprit d'équipe, gestion stress. \\
\textbullet\ Bricolage électronique (Arduino, domotique) : Veille technologique, autonomie. \\
\hline
\end{tabularx}
\end{center}

\noindent \textbf{Conclusion :} La version corrigée respecte les normes visuelles, hiérarchise l'information (Profil $\rightarrow$ Expé $\rightarrow$ Formation $\rightarrow$ Compétences) et utilise le vocabulaire technique de l'offre analysée précédemment.
\end{exoresolu}

\courssub{Rédiger la structure interne : rubriques, formulation des missions (Verbes d'action)}

\begin{methode}[Rédiger les rubriques "Expériences" et "Compétences" avec des verbes d'action (Infinitif / Impératif)]
\textbf{Objectif :} Transformer des "tâches" en "réalisations mesurables" en utilisant les \cle{valeurs de l'infinitif} (valeur nominale, injonctive, consigne) et de l'\cle{impératif} (ordre, conseil, invitation) pour dynamiser le CV.
\begin{enumerate}
    \item \textbf{Lister les activités réelles} effectuées pour chaque poste (stage, job, projet scolaire significatif).
    \item \textbf{Choisir le mode verbal adapté au CV :}
        \begin{itemize}
            \item \textbf{Infinitif (Standard CV) :} Valeur \emph{nominale} (titre de compétence) ou \emph{injonctive/consigne} (description de mission). Ex: \og \cle{Assurer} la maintenance préventive \fg ; \og \cle{Rédiger} les comptes-rendus \fg. \emph{Avantage : neutre, concis, impersonnel, valorise l'action.}
            \item \textbf{Impératif (rare en CV, plutôt lettre) :} Valeur \emph{ordre/conseil}. Ex: \og \cle{Gérez} les stocks \fg. \emph{À éviter dans le CV moderne (trop direct), réserver à la lettre de motivation.}
        \end{itemize}
    \item \textbf{Appliquer la méthode \cle{STAR} (Situation, Tâche, Action, Résultat) pour chaque puce :}
        \begin{itemize}
            \item Verbe d'action à l'infinitif (passé ou présent selon actualité).
            \item Contexte technique / Outil / Machine.
            \item Résultat chiffré ou qualitatif (gain temps, taux, budget, sécurité).
        \end{itemize}
    \item \textbf{Classer les compétences techniques (Hard Skills) par familles} (ex: Automatismes, Mécanique, Informatique, Gestion) et indiquer le \cle{niveau} (Expert / Maîtrise / Notions / En formation).
    \item \textbf{Valoriser les Soft Skills par la preuve} : ne pas écrire "Esprit d'équipe" mais \og \cle{Collaborer} au sein d'une équipe de 5 techniciens pour réduire le MTTR de 15\% \fg.
\end{enumerate}
\end{methode}

\begin{exoresolu}[Réécriture de missions "plates" en réalisations "STAR" à l'infinitif]
\textbf{Énoncé :} Transformez les phrases suivantes (style "journal intime" ou "passif") en puces de CV percutantes (Infinitif + Contexte + Résultat). Le poste visé : Technicien Maintenance.

\begin{enumerate}
    \item "J'ai réparé des machines qui tombaient souvent en panne."
    \item "On m'a demandé de faire la maintenance préventive tous les mois."
    \item "Je connais bien la soudure à l'arc et le TIG."
    \item "J'ai aidé mon chef à faire l'inventaire des pièces."
    \item "Je suis capable de lire des plans électriques."
\end{enumerate}

\tcblower
\textbf{Solution détaillée :}

\noindent \textbf{Principes linguistiques :} Le CV utilise l'\textbf{infinitif} (valeur consigne/injonctive ou nominale) pour décrire les missions. C'est une forme non conjuguée, intemporelle, qui met l'accent sur l'\textbf{action} et la \textbf{compétence} plutôt que sur l'auteur (\og Je \fg). L'impératif est exclu (trop familier/ordre direct).

\noindent \textbf{Transformation STAR (Infinitif présent pour poste actuel, passé composé ou infinitif passé pour anciens postes -- \emph{usage courant : infinitif présent intemporel pour tous}).}

\begin{center}
\begin{tabularx}{\linewidth}{|l|X|}
\hline
\rowcolor{popblueL}
\textbf{Version "Élève" (Faible)} & \textbf{Version CV "Pro" (Forte -- Infinitif + STAR)} \\
\hline
1. J'ai réparé des machines qui tombaient souvent en panne. &
\textbullet\ \textbf{Dépanner} (Infinitif) des équipements de conditionnement (Sidel/Krones) en urgence : diagnostic panne, remplacement organes (vérins, capteurs, cartes API), remise en route. \textbf{Résultat :} Réduction du temps d'arrêt moyen (MTTR) de 45 à 20 min. \\
\hline
2. On m'a demandé de faire la maintenance préventive tous les mois. &
\textbullet\ \textbf{Réaliser} la maintenance préventive planifiée (gammes hebdo/mensuelles) sur ligne automatisée : graissage, contrôle usure, réglages, relevés vibratoires. \textbf{Résultat :} Taux de disponibilité ligne maintenu à 98\%, 0 accident lié maintenance. \\
\hline
3. Je connais bien la soudure à l'arc et le TIG. &
\textbullet\ \textbf{Maîtriser} les procédés de soudage (Arc enrobé, TIG Inox/Alu, MIG/MAG) pour réparation châssis, tuyauterie presseur. \textbf{Preuve :} Qualification soudeur QMOS 9606-1 (valide), réalisation skids inox agroalimentaire. \\
\hline
4. J'ai aidé mon chef à faire l'inventaire des pièces. &
\textbullet\ \textbf{Participer} à l'inventaire tournant mensuel du magasin pièces de rechange (5000 références) sous SAP PM : comptage, étiquetage, investigation écarts. \textbf{Résultat :} Fiabilité stock 99,5\%, réduction surstock 12\%. \\
\hline
5. Je suis capable de lire des plans électriques. &
\textbullet\ \textbf{Exploiter} la documentation technique (schémas unifilaires/multifilaires IEC, nomenclatures, notices constructeurs) pour diagnostic pannes complexes et modification armoires. \textbf{Preuve :} Modif. câblage variateurs VSD (schéma + câblage + paramétrage). \\
\hline
\end{tabularx}
\end{center}

\noindent \textbf{Analyse grammaticale (Langue) :}
\begin{itemize}
    \item \og \textbf{Dépanner}, \textbf{Réaliser}, \textbf{Maîtriser}, \textbf{Participer}, \textbf{Exploiter} \fg : \textbf{Infinitifs présents}. Valeur \emph{injonctive/consigne} (équivalent à une instruction de procédure) ou \emph{nominale} (titre de la compétence). Ils remplacent avantageusement "Je dépanne / J'ai dépanné / Il faut dépanner".
    \item L'absence de sujet (\og Je \fg) généralise la compétence : c'est \textbf{ce que le candidat SAIT FAIRE}.
\end{itemize}

\noindent \textbf{Conclusion :} Chaque puce commence par un verbe d'action fort à l'infinitif, précise le contexte technique (machines, normes, logiciels) et quantifie le résultat. C'est la "signature" d'un CV technique réussi.
\end{exoresolu}

\courssub{Produire un CV complet et cohérent (Synthèse : Identité + Profil + Expé + Formation + Compétences)}

\begin{methode}[Rédiger son CV complet "Maître" (Version de base à adapter)]
\textbf{Objectif :} Produire un document unique, sans faute, qui sert de \cle{matrice} pour toutes les candidatures futures.
\begin{enumerate}
    \item \textbf{Rassembler les données brutes} : diplômes (dates, mentions, lieux), expériences (dates exactes, noms exacts entreprises, missions détaillées), compétences (logiciels, machines, langues, permis, habilitations), projets (PFEE, projets techniques), centres d'intérêt (associatifs, sportifs, techniques).
    \item \textbf{Rédiger l'Accroche (Profil) :} 3 lignes max. \og [Diplôme/Niveau] + [Spécialité] + [Expérience clé/Spécificité] + [Objectif/Disponibilité] \fg. Ex: \og Technicien Supérieur en Maintenance Industrielle (BTS MSA), 2 ans d'alternance en agroalimentaire sur lignes automatisées Siemens. Expertise maintenance préventive/curative, GMAO SAP PM, normes HSE. Disponible immédiatement pour poste en 3x8. \fg
    \item \textbf{Remplir "Expériences" (Antéchronologique) :} Pour chaque ligne : Dates | Intitulé Poste \textbf{gras} | Entreprise \textit{italique} | Lieu | 3--5 puces \textbf{Infinitif + Contexte + Résultat}.
    \item \textbf{Remplir "Formation" (Antéchronologique) :} Dates | Diplôme exact \textbf{gras} | Option/Spécialité | Établissement | Ville | Mention (si $\ge$ Assez Bien) | Mémoire/Projet pertinent (1 ligne).
    \item \textbf{Structurer "Compétences" par catégories techniques :} Informatique/Logiciels (niveau), Langues (niveau CECRL + score), Habilitations/Permis (validité), Soft Skills \emph{prouvées} par exemple.
    \item \textbf{Relire "à froid" (J+1) :} Orthographe (surtout noms propres entreprises, logiciels), cohérence dates (trous ?), alignements, longueur (1 page si $<$ 5 ans expé, 2 pages max sinon).
    \item \textbf{Exporter en PDF} : Nom\_Prenom\_CV\_Année.pdf. Tester l'ouverture sur téléphone (lisibilité).
\end{enumerate}
\end{methode}

\begin{exoresolu}[Production d'un CV "Maître" pour un élève de Terminale STI (Électrotechnique)]
\textbf{Énoncé :} À partir des éléments biographiques fictifs ci-dessous, rédigez le contenu textuel complet du CV (structure interne) en appliquant la méthode. Mettez en valeur le stage de fin d'études (PFEE).

\emph{Données brutes :}
\begin{itemize}
    \item \textbf{Identité :} Ngono Essomba Marie-Claire, 18 ans, Yaoundé (Mvan), +237 6XX XX XX XX, marie.claire.ngono@email.cm.
    \item \textbf{Diplômes :} Baccalauréat Technique F3 (Électrotechnique) session 2024, Mention Assez Bien. Brevet d'Études 2020.
    \item \textbf{Stage PFEE (Juin--Juillet 2024, 8 sem.) :} ENEO, Direction Régionale Centre, Service Maintenance Postes HTA/BT. Tuteur : M. Fouda. Missions : Suivi maintenance préventive postes HTA (disjoncteurs, relais protection), relevés thermographiques, participation remplacement transformateur 630 kVA, rédaction rapports hebdo sous Excel.
    \item \textbf{Projet Technique (Année de Terminale) :} "Automatisation d'un système de pompage solaire" : Arduino, capteurs niveau, variateur, pompe 3ph. Réalisation schéma, câblage, prog. Présenté à la Semaine de l'Excellence Technique.
    \item \textbf{Compétences Techniques :} Lecture plans unifilaires/multifilaires, Câblage armoires BT, Mesures (Multimètre, Mégaohmmètre, Caméra thermo), Notions AutoCAD Elec, Pack Office (Excel avancé : TCD, RechercheV).
    \item \textbf{Langues :} Français (Maternel), Anglais (Niveau A2/B1 scolaire - lu/écrit technique).
    \item \textbf{Habilitations :} H0B0 (théorique acquise en cours), Préparation BR/BC en cours.
    \item \textbf{Centres d'intérêt :} Volley-ball (Club scolaire, Capitaine), Bénévolat Croix-Rouge (Secourisme), Lecture revue "Énergie \& Environnement".
\end{itemize}

\tcblower
\textbf{Solution détaillée : CV Maîtresse (Contenu textuel prêt à mettre en page)}

\noindent \textsc{NGONO ESSOMBA} \textbf{Marie-Claire} \hfill +237 6XX XX XX XX \hfill marie.claire.ngono@email.cm \\
Yaoundé, Mvan \hfill 18 ans \hfill Permis B (en cours) \hfill Disponibilité : Immédiate

\medskip
\noindent \textbf{PROFIL}
\emph{Future Technicienne Supérieure en Électrotechnique (Bac F3, Mention AB), rigoureuse et autonome.} Stage validé chez ENEO (Maintenance Postes HTA/BT) et projet de fin d'études en automatisme solaire (Arduino/Variateur). Maîtrise lecture plans, câblage BT, mesures de protection. Engagée (Secourisme, Capitaine volley). Recherche alternance BTS/DUT Électrotechnique ou 1er emploi technicien maintenance à partir Septembre 2024.

\medskip
\noindent \textbf{EXPÉRIENCES PROFESSIONNELLES \& STAGES}

\noindent \textbf{06/2024 -- 07/2024 (2 mois)} \hfill \textbf{Stagiaire Technicienne Maintenance Postes HTA/BT} \\
\textit{ENEO Cameroun, Direction Régionale Centre, Yaoundé} \\
\textbullet\ \textbf{Assister} les techniciens confirmés dans la maintenance préventive systématique des postes HTA (disjoncteurs SF6, relais de protection Numériques Sepam/Ref615). \\
\textbullet\ \textbf{Effectuer} des campagnes de contrôle thermographique sur appareillages HTA/BT (Caméra FLIR) : détection 3 points chauds critiques traités en curatif. \\
\textbullet\ \textbf{Participer} au remplacement d'un transformateur 630 kVA / 30 kV - 0,4 kV : manœuvres de consignation (H0B0), contrôle isolement (Mégaohmmètre 5kV), raccordement BT, mise en service. \\
\textbullet\ \textbf{Rédiger} les comptes-rendus d'intervention hebdomadaires sous Excel (TCD pour suivi KPI : nombre interventions, taux respect planning). \\
\textbullet\ \textbf{Appliquer} strictement les consignes de sécurité (Port EPI, Balisage, Consignation, Plan de Prévention).

\medskip
\noindent \textbf{PROJETS TECHNIQUES SIGNIFICATIFS}

\noindent \textbf{2023 -- 2024} \hfill \textbf{Projet de Fin d'Études : Automatisation Pompage Solaire (Equipe 3)} \\
\textbullet\ \textbf{Concevoir} et \textbf{réaliser} l'armoire de commande (Schémas AutoCAD Elec, Câblage, Bornier, Sélecteurs). \\
\textbullet\ \textbf{Programmer} l'automate Arduino (Gestion niveau cuve, Protection pompe marche à sec, MPPT solaire simplifié). \\
\textbullet\ \textbf{Paramétrer} le variateur de fréquence (Schneider Altivar 12) pour démarrage progressif pompe 3ph 2,2 kW. \\
\textbullet\ \textbf{Présenter} le projet lors de la Semaine de l'Excellence Technique (Oral 15 min + Démonstration) : Mention "Projet Innovant".

\medskip
\noindent \textbf{FORMATION}

\noindent \textbf{2024} \hfill \textbf{Baccalauréat Technique Série F3 (Électrotechnique)} \hfill \textbf{Mention Assez Bien} \\
Lycée Technique Industriel de Yaoundé \\
\textbullet\ Spécialité : Machines tournantes, Automatismes, Électronique de puissance, Dessin technique. \\
\textbullet\ Projet PFEE : "Maintenance Postes HTA/BT" (Note : 16/20).

\noindent \textbf{2020} \hfill \textbf{Brevet d'Études du Premier Cycle} \hfill Mention Bien \\
Collège de Mvan, Yaoundé.

\medskip
\noindent \textbf{COMPÉTENCES TECHNIQUES (HARD SKILLS)}

\begin{center}
\begin{tabularx}{\linewidth}{|l|X|}
\hline
\rowcolor{popblueL}
\textbf{Domaine} & \textbf{Détails \& Niveau} \\
\hline
\textbf{Électrotechnique / Maintenance} & Lecture plans unifilaires/multifilaires (Maîtrise) ; Câblage armoires BT normes NFC 15-100 (Maîtrise) ; Mesures protection/isolement/thermographie (Maîtrise stage) ; Notions HTA (Disjoncteurs, Relais, Transfos) (Stage) \\
\hline
\textbf{Automatismes / Informatique Ind.} & Arduino (Programmation C++, Maîtrise Projet) ; Variateurs vitesse (Schneider Altivar, Notions Projet) ; AutoCAD Electrical (Schémas, Notions) ; GMAO (Notions théoriques) \\
\hline
\textbf{Bureautique / Outils} & Pack Office : Excel Avancé (TCD, RechercheV, Formules imbriquées) ; Word (Rapports techniques) ; PowerPoint (Soutenances) \\
\hline
\end{tabularx}
\end{center}

\medskip
\noindent \textbf{LANGUES \& CERTIFICATIONS}
\textbullet\ Français : Langue maternelle. \quad \textbullet\ Anglais : Niveau B1 (Scolaire + Technique lu/écrit) -- TOEIC prévu 2025. \\
\textbullet\ Habilitation Électrique : H0B0 (Validée formation) -- \textbf{BR/BC en cours de validation pratique}. \\
\textbullet\ Secourisme : Attestation PSC1 (Croix-Rouge, 2023, Valide).

\medskip
\noindent \textbf{CENTRES D'INTÉRÊT (Valorisés)}
\textbullet\ \textbf{Volley-ball (Capitaine Équipe Lycée, 3 ans) :} Leadership, gestion du stress match, cohésion d'équipe, assiduité. \\
\textbullet\ \textbf{Bénévolat Croix-Rouge Yaoundé (Secouriste, 2 ans) :} Réactivité, sang-froid, respect protocoles, service public. \\
\textbullet\ \textbf{Veille technique :} Lecture mensuelle "Énergie \& Environnement", "Elec Magazine" (Curiosité, mise à jour normes).

\medskip
\noindent \textbf{Conclusion :} Ce CV "Maître" est dense, factuel, sans "Je", utilise l'infinitif pour les actions, quantifie (2 mois, 630 kVA, 3 points chauds, 3 ans, 16/20) et aligne les compétences sur les attendus des BTS/Entreprises (ENEO, SONARA, CDC, etc.). Il tient sur 1 page A4 bien mise en page.
\end{exoresolu}

\courssub{Adapter le CV à une offre ciblée (Personnalisation / "Tailoring")}

\begin{methode}[Adapter son CV "Maître" à une offre d'emploi spécifique]
\textbf{Objectif :} Maximiser le taux de passage aux ATS (logiciels de tri) et l'intérêt du recruteur humain en alignant \textbf{mots-clés}, \textbf{ordre des rubriques} et \textbf{lexique} sur l'offre.
\begin{enumerate}
    \item \textbf{Analyser l'offre cible} (Méthode 1) : Extraire les 10--15 mots-clés "Dur" (Diplôme, Logiciels, Machines, Normes, Permis) et "Souple" (Soft skills).
    \item \textbf{Comparer avec son CV Maître} : Surligner ce qui match, identifier les manques (compétences transférables ? formations courtes ?).
    \item \textbf{Adapter l'Accroche (Profil) :} Réécrire la 1ère phrase en reprenant \emph{l'intitulé exact du poste} et les 3 compétences "Obligatoires" n°1 de l'offre.
    \item \textbf{Réordonner / Filtrer les Expériences :} Mettre en \textbf{premier} (ou en gras) l'expérience la plus pertinente pour \emph{ce} poste. Masquer ou résumer les expériences sans lien (jobs d'été non techniques) sauf si elles prouvent une soft skill clé (ex: gestion caisse = rigueur/confiance).
    \item \textbf{Injecter les mots-clés techniques} dans les puces de missions (synonymes acceptés : \og GMAO \fg $\leftrightarrow$ \og SAP PM / Maintenance Assistée \fg). Utiliser le \emph{vocabulaire de l'entreprise} (ex: "Opérateur" vs "Technicien", "Ligne" vs "Chaîne").
    \item \textbf{Ajuster la rubrique Compétences :} Monter en haut les compétences "Obligatoires" de l'offre. Ajouter une ligne "En formation / Autodidacte" pour les compétences "Souhaitées" non encore maîtrisées (honnêteté).
    \item \textbf{Vérifier la cohérence globale :} Le CV raconte-t-il une histoire qui mène \emph{logiquement} à ce poste ? Enregistrer : \texttt{Nom\_Prenom\_CV\_[IntituléPoste]\_[Entreprise].pdf}.
\end{enumerate}
\end{methode}

\begin{exoresolu}[Adaptation du CV Maître pour une offre "Technicien Maintenance Cimenterie"]
\textbf{Énoncé :} Marie-Claire (CV Maître ci-dessus) postule chez \textbf{CIMENCAM} (Cimenteries du Cameroun) pour un poste de \og Technicien Maintenance Électrotechnicien Junior \fg.
\emph{Extrait Offre CIMENCAM :} \og BTS/DUT/Bac+2 Électrotechnique/Maintenance. Expérience stage/alternance bienvenue. \textbf{Connaissances impératives :} Variateurs de vitesse (Altivar, Sinamics), Automates (Siemens S7, Schneider M241), Lecture plans HTA/BT, Maintenance préventive/curative fours/broyeurs/convoieurs. \textbf{Habilitation BR/BC obligatoire}. Travail par roulement (équipes). Esprit sécurité, rigueur, disponibilité. \fg

\textbf{Travail demandé :} Montrez les modifications précises à apporter au CV Maître (Accroche, Expériences, Compétences, Habilitations) pour cette candidature.

\tcblower
\textbf{Solution détaillée :}

\noindent \textbf{Analyse Offre (Mots-clés Obligatoires) :} Variateurs (Altivar/Sinamics), Automates (S7/M241), Plans HTA/BT, Maintenance lourde (Fours/Broyeurs/Convoieurs), \textbf{BR/BC Obligatoire}, Roulement, Sécurité.

\noindent \textbf{Comparaison avec CV Maître :}
\begin{itemize}
    \item \textbf{Match :} Bac F3 (Électrotechnique), Stage ENEO (HTA/BT, Relais, Transfo, Thermographie, Consignation), Projet Arduino/Variateur (Altivar 12), Lecture plans, Câblage.
    \item \textbf{Ecart (Gap) :} \textbf{BR/BC non validée (seule H0B0)}, Pas d'expé "Fours/Broyeurs/Convoieurs" (mais transférable : machines tournantes, fortes puissances), Automates Siemens S7 (notions théoriques seulement), Schneider M241 (inconnu).
\end{itemize}

\noindent \textbf{Modifications Apportées (CV Ciblé CIMENCAM) :}

\medskip
\noindent \textbf{1. ACCROCHE MODIFIÉE (Alignée sur intitulé + 3 top skills) :}
\og \textbf{Technicienne Maintenance Électrotechnicienne Junior (Bac F3 AB)}, stage ENEO Maintenance Postes HTA/BT (Relais Sepam, Transfo 630kVA, Thermographie, Consignation). \textbf{Maîtrise variateurs Altivar (Projet), Automates Siemens (Notions), Lecture plans HTA/BT.} Habilitation H0B0 acquise, \textbf{BR/BC en finalisation}. Disponible immédiatement pour travail en équipes (roulement). Engagée Sécurité (PSC1). \fg

\medskip
\noindent \textbf{2. EXPÉRIENCES : REFORMULATION DES PUCES (Injection vocabulaire Cimenterie) :}
\begin{itemize}
    \item \emph{Ancienne puce :} "Assister les techniciens... maintenance préventive systématique postes HTA..."
    \item \emph{Nouvelle puce :} \textbf{Exécuter} la maintenance préventive et curative sur appareillages HTA/BT (Disjoncteurs, Relais \textbf{protection Sepam} -- équivalent process cimenterie) : \textbf{planifier} les gammes, \textbf{diagnostiquer} par thermographie, \textbf{remplacer} organes en sécurité (Consignation).
    \item \emph{Ajout :} \textbf{Participer} à la maintenance de \textbf{machines tournantes de forte puissance} (Transformateurs 630kVA, Moteurs HT) : contrôle isolement, vibration, alignement (transférable broyeurs/fours).
\end{itemize}

\medskip
\noindent \textbf{3. PROJET TECHNIQUE : REVALORISATION "VARIATEURS / AUTOMATISMES" :}
\textbullet\ \textbf{Programmer} et \textbf{paramétrer} un \textbf{variateur Schneider Altivar 12} (famille \textbf{Altivar Process} en cimenterie) pour entraînement pompe 3ph : rampes, PID, protection moteur.
\textbullet\ \textbf{Programmer} automate \textbf{Arduino} (logique échelle/Grafcet) $\rightarrow$ \textbf{Bases automatismes transférables sur Siemens S7-1200/1500 (TIA Portal) et Schneider M241 (SoMachine)}.

\medskip
\noindent \textbf{4. COMPÉTENCES TECHNIQUES : RECLASSMENT PAR PRIORITÉ OFFRE :}

\begin{center}
\begin{tabularx}{\linewidth}{|l|X|}
\hline
\rowcolor{popblueL}
\textbf{Priorité Offre} & \textbf{Rubrique Compétences CV Ciblé (Extrait)} \\
\hline
\textbf{1. Variateurs} & \textbf{Variateurs de vitesse :} \textbf{Altivar 12/312 (Maîtrise Projet)}, Notions Sinamics (Siemens). Paramétrage, diagnostic pannes, démarrage progressif. \\
\hline
\textbf{2. Automates} & \textbf{Automatismes :} \textbf{Siemens S7-1200/1500 (Notions théoriques + TIA Portal découverte)}, \textbf{Schneider M241 (Notions SoMachine)}, Arduino (Maîtrise). Langages : LAD, FBD, SCL (Notions). \\
\hline
\textbf{3. Plans HTA/BT} & \textbf{Dessin / Lecture Plans :} Schémas unifilaires/multifilaires HTA/BT (Maîtrise Stage ENEO + Cours), AutoCAD Electrical (Notions). Normes IEC 61346, NF C 15-100. \\
\hline
\textbf{4. Maintenance Process} & \textbf{Maintenance Industrielle :} Préventive (Gammes, Planning), Curative (Diagnostic, Remplacement), \textbf{Machines tournantes (Moteurs, Transfo, Pompes - Stage/Projet)}, Outils : Mégaohmmètre 5kV, Caméra Thermo, Analyseur vibra (Notions). \\
\hline
\textbf{5. Habilitation (CRITIQUE)} & \textbf{Habilitations :} \textbf{H0B0 (Validée)}. \textbf{BR/BC : Formation théorique validée, Stage pratique en cours (Finalisation Juillet 2024)}. \textcolor{red}{$\leftarrow$ HONNÊTETÉ CRUCIALE} \\
\hline
\end{tabularx}
\end{center}

\medskip
\noindent \textbf{5. CENTRES D'INTÉRÊT : CIBLAGE "ÉQUIPES / SÉCURITÉ / DISPONIBILITÉ" :}
\textbullet\ Volley-ball (Capitaine) : \textbf{Travail en équipe, roulement, gestion fatigue/stress}.
\textbullet\ Secourisme PSC1 (Croix-Rouge) : \textbf{Culture Sécurité, Réactivité urgence, Respect protocoles}.
\textbullet\ Veille technique "Énergie \& Environnement" : \textbf{Connaissance secteur ciment/énergie}.

\medskip
\noindent \textbf{Conclusion :} Le CV ciblé ne ment pas, mais \textbf{réordonne}, \textbf{reformule} avec le vocabulaire de l'offre (Altivar, Sepam, S7, BR/BC, Roulement, Sécurité) et met en avant les \textbf{compétences transférables} (HTA/BT $\rightarrow$ Process lourd, Arduino $\rightarrow$ Automates). Il signale proactivement la finalisation de la BR/BC. C'est ce CV qu'il faut envoyer à CIMENCAM.
\end{exoresolu}

\courssub{Auto-évaluation et remédiation : Grille de relecture critique (Check-list Bac)}

\begin{methode}[S'auto-évaluer et corriger son CV avant envoi (Grille de remédiation)]
\textbf{Objectif :} Garantir un CV "zéro faute" formel et "zéro incohérence" fond, prêt pour le Probatoire/Bac et le marché.
\begin{enumerate}
    \item \textbf{Check-list FORME (Élimination visuelle) :}
        \begin{itemize}
            \item [ ] Fichier PDF, nom normalisé (\texttt{Nom\_Prenom\_Poste.pdf}).
            \item [ ] 1 page (Junior) / 2 pages max (Expérimenté). Pas de page blanche page 2.
            \item [ ] Police unique pro (Calibri/Arial 10-11), Titres 13-14. Couleurs max 2 (Bleu marine + Gris).
            \item [ ] Alignements parfaits (Dates gauche, Lieux droite, Puces alignées). Pas de "veuves" (titre seul bas page).
            \item [ ] Orthographe/Grammaire : 0 faute (Correcteur + Lecture à voix haute + Tierce relecture).
            \item [ ] Coordonnées valides (Tel, Email pro, Ville). LinkedIn à jour si lien mis.
        \end{itemize}
    \item \textbf{Check-list FOND (Pertinence Bac/Recruteur) :}
        \begin{itemize}
            \item [ ] \textbf{Accroche :} Personnalisée ? Intitulé poste + 3 skills clés + Dispo ?
            \item [ ] \textbf{Expériences :} Ordre antéchronologique ? Verbes \textbf{Infinitif} ? Méthode STAR (Action + Contexte Technique + Résultat Chiffré) ? Pas de "J'ai fait" ?
            \item [ ] \textbf{Formation :} Diplômes exacts, Mentions, Projets techniques (PFEE) valorisés ?
            \item [ ] \textbf{Compétences :} Classées par familles techniques ? Niveaux honnêtes (Expert/Maîtrise/Notions) ? Mots-clés offre présents ?
            \item [ ] \textbf{Habilitations/Permis/Langues :} Validité dates ? Niveaux CECRL ? Scores ?
            \item [ ] \textbf{Cohérence globale :} Le CV raconte-t-il une progression logique vers le poste visé ? Pas de trous inexpliqués > 6 mois (voyage, formation, recherche active) ?
        \end{itemize}
    \item \textbf{Test "Recruteur 30 secondes" :} Demander à un pair/enseignant : "En 30 sec, quel poste je vise ? Quelles sont mes 3 forces ?" $\rightarrow$ Si hésitation $\rightarrow$ Remédiation Accroche / Mots-clés.
    \item \textbf{Versionning :} Garder le CV Maître. Créer une version par candidature. Archiver les réponses.
\end{enumerate}
\end{methode}

\begin{exoresolu}[Application de la grille de remédiation sur un CV "Presque parfait"]
\textbf{Énoncé :} Voici un extrait de CV d'un élève de Terminale (Électrotechnique) qui pense son CV prêt. Appliquez la grille de remédiation. Identifiez 5 erreurs (Forme ou Fond) et corrigez-les.

\emph{Extrait CV Élève :}
\begin{quote}
\textbf{COMPÉTENCES INFORMATIQUES}
\textbullet\ Word, Excel, PowerPoint (Très bien)
\textbullet\ AutoCAD (Bien)
\textbullet\ Programmation Arduino (Je sais faire)
\textbullet\ GMAO : Je connais SAP

\textbf{STAGE ENEO JUIN-JUILLET 2024}
\textbullet\ J'ai aidé pour la maintenance des postes.
\textbullet\ J'ai fait des mesures à la pince ampèremétrique.
\textbullet\ J'ai appris à lire des plans haut tension.
\textbullet\ C'était très intéressant.

\textbf{ANGLAIS} : Niveau Moyen (Lu, écrit, parlé un peu).
\textbf{PERMIS} : En cours (Théorie ok).
\end{quote}

\tcblower
\textbf{Solution détaillée : Application de la Grille}

\noindent \textbf{Erreur 1 (Forme/Fond - Verbes) :} Utilisation de la 1ère personne (\og J'ai aidé, J'ai fait, J'ai appris \fg) et verbes faibles (\og aider, faire, apprendre \fg).
\begin{itemize}
    \item \textbf{Correction :} Remplacer par \textbf{Infinitif} + Verbes d'action techniques + Contexte + Résultat.
    \item \og \textbf{Assister} les techniciens dans la maintenance préventive des postes HTA/BT (Disjoncteurs, Relais). \fg
    \item \og \textbf{Effectuer} des mesures de courant/ tension (Pince ampèremétrique, Multimètre, Mégaohmmètre) pour diagnostic. \fg
    \item \og \textbf{Exploiter} des schémas unifilaires HTA/BT pour repérage circuits et consignation. \fg
\end{itemize}

\noindent \textbf{Erreur 2 (Fond - Compétences Informatiques : Niveaux vagues) :} "Très bien", "Bien", "Je sais faire", "Je connais" $\rightarrow$ Subjectif, non standard, non vérifiable.
\begin{itemize}
    \item \textbf{Correction :} Niveaux standardisés + Détail fonctionnel.
    \item \og Pack Office : Excel \textbf{Avancé} (TCD, RechercheV/V, Formules complexes), Word \textbf{Maîtrise} (Rapports techniques), PowerPoint \textbf{Maîtrise} (Soutenances). \fg
    \item \og AutoCAD Electrical : \textbf{Notions} (Lecture schémas, Création symboles simples). \fg
    \item \og Arduino IDE : \textbf{Maîtrise} (Prog. C++, Librairies, Interface capteurs/actionneurs, Projet PFEE). \fg
    \item \og GMAO SAP PM : \textbf{Notions Utilisateur} (Création notifications, Consultation ordres, Stage ENEO). \fg
\end{itemize}

\noindent \textbf{Erreur 3 (Fond - Anglais : Niveau non certifié/flou) :} "Moyen (Lu, écrit, parlé un peu)".
\begin{itemize}
    \item \textbf{Correction :} Référentiel CECRL + Contexte technique.
    \item \og Anglais : \textbf{Niveau B1 (Seuil)} -- Compréhension technique (Notices, Datasheets) : Autonome. Expression orale : Élémentaire. \textbf{TOEIC visé 2025}. \fg
\end{itemize}

\noindent \textbf{Erreur 4 (Forme - Permis : Imprécision) :} "En cours (Théorie ok)". Le recruteur a besoin de savoir \textbf{quand} il sera obtenu.
\begin{itemize}
    \item \textbf{Correction :} \og Permis B : \textbf{Code validé (Juin 2024)}, Conduite en cours (Examen prévu Sept 2024). \fg
\end{itemize}

\noindent \textbf{Erreur 5 (Fond - Stage : Absence de résultat / "C'était intéressant") :} "C'était très intéressant" est une appréciation subjective, inutile dans un CV. Manque de résultats chiffrés.
\begin{itemize}
    \item \textbf{Correction :} Supprimer l'appréciation. Ajouter une puce résultat.
    \item \og \textbf{Rédiger} 4 rapports d'intervention hebdomadaires validés par le tuteur (Suivi KPI : 15 interventions/semaine, 0 écart sécurité). \fg
\end{itemize}

\noindent \textbf{Version Corrigée (Extrait) :}

\begin{center}
\begin{tabularx}{\linewidth}{|X|}
\hline
\rowcolor{popgreenL} \textbf{COMPÉTENCES TECHNIQUES \& INFORMATIQUES} \\
\hline
\textbullet\ \textbf{Automatismes / Électrotechnique :} Lecture plans HTA/BT (Maîtrise), Câblage armoires BT (Maîtrise), Mesures protection/isolement (Maîtrise), Variateurs Altivar (Notions Projet), Automates Siemens/Arduino (Notions/Projet). \\
\textbullet\ \textbf{Informatique Industrielle :} AutoCAD Electrical (Notions), GMAO SAP PM (Notions Utilisateur Stage). \\
\textbullet\ \textbf{Bureautique :} Excel Avancé (TCD, RechercheV, Macros bases), Word/PowerPoint (Maîtrise). \\
\hline
\rowcolor{poporangeL} \textbf{EXPÉRIENCE PROFESSIONNELLE} \\
\hline
\textbf{06/2024 -- 07/2024} \hfill \textbf{Stagiaire Maintenance Postes HTA/BT} \hfill \textit{ENEO, Yaoundé} \\
\textbullet\ \textbf{Assister} la maintenance préventive/curative postes HTA (Disjoncteurs SF6, Relais Sepam). \\
\textbullet\ \textbf{Effectuer} mesures thermographiques, isolement (Mégaohmmètre 5kV), courant (Pince). \\
\textbullet\ \textbf{Exploiter} schémas unifilaires HTA/BT pour consignation et repérage pannes. \\
\textbullet\ \textbf{Rédiger} 4 rapports hebdo (Excel TCD) : 15 interv/sem, 0 écart sécurité, 3 points chauds détectés. \\
\hline
\rowcolor{popblueL} \textbf{LANGUES \& PERMIS} \\
\hline
\textbullet\ Français : Maternel \quad \textbullet\ Anglais : B1 Technique (Datasheets autonomes) -- TOEIC prévu 2025. \\
\textbullet\ Permis B : Code validé (06/2024) -- Examen conduite 09/2024. \quad \textbullet\ Habilitation : H0B0 (Validée), BR/BC (En cours). \\
\hline
\end{tabularx}
\end{center}

\noindent \textbf{Conclusion :} La remédiation transforme un CV "élève" (subjectif, passif, vague) en CV "technicien junior" (objectif, actif, normé, vérifiable). C'est cette rigueur qui fait la différence au Bac et en entreprise.
\end{exoresolu}

\courssub{Les 5 erreurs qui coûtent des points au Bac (et en recrutement)}

\begin{attention}[Les 5 erreurs fatales au Probatoire / Bac Technique -- Production CV]
\begin{enumerate}
    \item \textbf{Utiliser le "Je" et les temps conjugués (Passé Composé / Imparfait) au lieu de l'Infinitif.}
    \og J'ai réparé... \fg $\rightarrow$ \textbf{Zéro point style pro.} \textbf{Reflexe :} \cle{Infinitif} (Dépanner, Réaliser, Programmer). C'est la norme du CV moderne et l'application directe du cours de Langue (Valeurs de l'infinitif : consigne/injonction/nominalisation).
    \item \textbf{Lister des "tâches" sans Résultat ni Contexte Technique (Méthode STAR absente).}
    \og Maintenance machines \fg $\rightarrow$ \textbf{Vide.} \textbf{Reflexe :} \cle{Verbe + Machine/Logiciel/Norme + Chiffre/Preuve.} (Ex: \og Maintenance préventive ligne Krones (SAP PM) : 98\% dispo \fg).
    \item \textbf{Niveaux de compétences subjectifs ("Bien", "Moyen", "Bon") sans référentiel.}
    \og Anglais : Bien \fg $\rightarrow$ \textbf{Non professionnel.} \textbf{Reflexe :} \cle{CECRL (A1--C2)} + Contexte (Technique/Oral) + Score si existant (TOEIC 750).
    \item \textbf{CV "Unique" non adapté à l'offre (Mots-clés absents, Accroche générique).}
    Envoyer le même CV pour "Maintenance Cimenterie" et "Helpdesk Informatique" $\rightarrow$ \textbf{Rejet ATS / Recruteur.} \textbf{Reflexe :} \cle{Tailoring} (Injection mots-clés offre, Réordre rubriques, Accroche ciblée).
    \item \textbf{Fautes d'orthographe, incohérences de dates, mise en page brouillonne (Police fantaisie, marges 0.5cm, photo plage).}
    \textbf{Élimination immédiate} (Forme = Fiabilité). \textbf{Reflexe :} \cle{Grille Remédiation} (PDF, 1 page, Police pro, Alignements, Relecture J+1 + Tierce personne).
\end{enumerate}
\end{attention}

\newpage

\coursec{Exercices}

\courssub{Je vérifie mes connaissances}

\begin{exercice}[QCM : les rubriques du CV]
Pour chaque question, choisis la bonne réponse.
\begin{enumerate}
\item Dans un CV, l'ordre \cle{antéchronologique} consiste à présenter les informations :
\begin{enumerate}
\item de la plus ancienne à la plus récente ;
\item de la plus récente à la plus ancienne ;
\item dans un ordre libre.
\end{enumerate}
\item La rubrique « État civil » (ou « Informations personnelles ») contient obligatoirement :
\begin{enumerate}
\item le nom, les prénoms, les contacts du candidat ;
\item les loisirs du candidat ;
\item le salaire souhaité.
\end{enumerate}
\item Dans la rubrique « Compétences », les verbes sont le plus souvent employés :
\begin{enumerate}
\item au passé composé ;
\item à l'infinitif ;
\item au futur simple.
\end{enumerate}
\item Un CV destiné à une offre d'emploi doit idéalement tenir sur :
\begin{enumerate}
\item trois pages minimum ;
\item une page (deux au maximum) ;
\item autant de pages que nécessaire.
\end{enumerate}
\end{enumerate}
\end{exercice}

\begin{exercice}[Vrai ou faux ? Justifie ta réponse]
Réponds par vrai ou faux, puis justifie en une phrase.
\begin{enumerate}
\item Il est conseillé d'adapter son CV à chaque offre d'emploi.
\item La photo est toujours obligatoire sur un CV au Cameroun.
\item Dans un CV, on peut rédiger toutes les rubriques en longues phrases rédigées.
\item L'impératif dans un avis de recrutement exprime toujours un souhait.
\end{enumerate}
\end{exercice}

\begin{exercice}[Définitions]
Définis en deux ou trois phrases chacun des termes suivants :
\begin{enumerate}
\item le \cle{CV} (curriculum vitae) ;
\item la \cle{structure externe} d'un CV ;
\item la \cle{structure interne} d'un CV ;
\item un \cle{texte fonctionnel}.
\end{enumerate}
\end{exercice}

\begin{exercice}[Repérage des modes]
Dans chacune des phrases suivantes, extraites d'avis de recrutement, indique le mode du verbe souligné (impératif ou infinitif) et précise sa valeur (ordre, conseil, consigne, énumération).
\begin{enumerate}
\item \emph{Envoyez} votre dossier avant le 30 juin.
\item \emph{Maîtriser} les outils informatiques est exigé.
\item \emph{Présentez}-vous à l'accueil muni de votre pièce d'identité.
\item Le candidat devra \emph{assurer} la maintenance des machines et \emph{rédiger} les rapports.
\end{enumerate}
\end{exercice}

\courssub{Je m'entraîne}

\begin{exercice}[Remise en ordre d'un CV]
Les huit rubriques du CV de Mme ESSOMBA ont été mélangées. Reconstitue le CV en les remettant dans l'ordre logique : \emph{Expérience professionnelle ; Compétences ; Formation ; État civil ; Centres d'intérêt ; Langues ; Références ; Titre du poste recherché}. Justifie l'ordre choisi.
\end{exercice}

\begin{exercice}[Correction d'un extrait de CV]
L'extrait suivant contient cinq fautes d'orthographe et deux incohérences de dates. Relève-les et corrige-les.

\emph{« Formation : 2020-2021 : BEPC, collège bilingue de Edéa. 2021-2024 : Baccalauréat technique série F4, lycée technique de Douala-Bassa. Expérience : juillet 2023 à janvier 2022 : stage d'ouvrier qualifié à la SODECOTON. J'ai travailler à la maintanance des machines et j'était chargé de l'acceuil des clients. »}
\end{exercice}

\begin{exercice}[De la phrase à l'énumération]
Transforme le paragraphe suivant, rédigé en phrases longues, en une rubrique « Compétences » rédigée en énumérations concises à l'infinitif.

\emph{« Pendant mon stage, j'ai appris à conduire les engins de chantier et je savais aussi lire les plans de construction. En plus, je réparais les installations électriques et je tenais le registre du matériel. »}
\end{exercice}

\begin{exercice}[Impératif et infinitif dans un avis de recrutement]
Voici un extrait d'avis de recrutement publié par une entreprise de Yaoundé :

\emph{« La société CAMBATIMENT recrute un maçon qualifié. Effectuer les travaux de gros œuvre ; respecter les normes de sécurité ; rendre compte au chef de chantier. Les candidats intéressés déposeront leurs dossiers au secrétariat. Contactez-nous au 6 90 00 00 00. Ne vous présentez pas après la date limite. »}

\begin{enumerate}
\item Relève cinq verbes à l'infinitif et deux verbes à l'impératif.
\item Justifie la valeur de chaque verbe relevé (consigne, énumération de tâches, ordre, interdiction).
\end{enumerate}
\end{exercice}

\courssub{Je me prépare au Bac}

\begin{exercice}[Sujet type Baccalauréat technique : rédaction complète d'un CV]
\textbf{Situation.} L'entreprise ELECTROCAM, située à Douala-Bonabéri, publie l'offre suivante : \emph{« Recrutement d'un(e) technicien(ne) en électronique. Profil : Baccalauréat technique série F3 ou équivalent ; maîtriser la lecture de schémas électroniques ; assurer la maintenance des équipements ; rédiger les rapports d'intervention. Dossier : CV, lettre de motivation, photocopie des diplômes. Date limite : 15 septembre. »}

Tu es candidat(e) à ce poste. Rédige ton CV complet en respectant la structure externe et la structure interne étudiées en classe. Tu inventeras les informations personnelles, la formation (ordre antéchronologique), l'expérience et les compétences, en les adaptant à l'offre.

\textbf{Barème indicatif :} prise en compte de la situation de communication : 4 pts ; structure du CV (rubriques, ordre, cohérence) : 6 pts ; correction de la langue (modes, orthographe) : 6 pts ; présentation matérielle : 4 pts.
\end{exercice}

\begin{exercice}[Adaptation du CV à deux situations d'embauche]
\textbf{Situation.} M. NGONO, 22 ans, titulaire d'un Baccalauréat technique série G2 (comptabilité), a effectué un stage de trois mois dans une PME de Bafoussam (tenue des livres de comptes) et un emploi de vacances de deux mois comme caissier dans un supermarché. Il postule à deux offres :
\begin{enumerate}
\item \emph{Offre A} : « PME recherche un comptable : tenir la comptabilité courante ; établir les factures ; maîtriser un logiciel de gestion. »
\item \emph{Offre B} : « Supermarché recherche un vendeur : accueillir les clients ; conseiller la clientèle ; gérer la caisse et les rayons. »
\end{enumerate}

\begin{enumerate}
\item Rédige, pour chacune des deux offres, les rubriques « Compétences » et « Expérience professionnelle » du CV de M. NGONO, en sélectionnant et en ordonnant les informations pertinentes.
\item Justifie en cinq à huit lignes les choix opérés : pourquoi certaines informations sont mises en avant dans une version et non dans l'autre ?
\end{enumerate}
\end{exercice}

\begin{exercice}[Sujet type Probatoire : rubriques à partir d'une fiche d'identité professionnelle]
\textbf{Situation.} Voici la fiche d'identité professionnelle de Mlle FOTSO, candidate à un poste de secrétaire de direction à Bafoussam :

\begin{center}
\begin{tabularx}{\linewidth}{|l|X|}\hline
\rowcolor{popblueL} \textbf{Rubrique} & \textbf{Informations} \\ \hline
Identité & FOTSO Léa, née le 12 mars 2003 à Dschang ; tél. 6 99 11 22 33 ; lea.fotso@mail.com \\ \hline
Études & 2019-2020 : Probatoire G1, lycée de Dschang. 2020-2021 : Baccalauréat G1, lycée de Dschang. 2021-2023 : BTS en secrétariat de direction, IST de Bafoussam. \\ \hline
Expérience & Sept. 2022 - déc. 2022 : stage de secrétaire à la mairie de Bafoussam. Janv. 2023 - juin 2024 : secrétaire comptable, cabinet FOKAM et Associés. \\ \hline
Divers & Français : courant ; anglais : niveau scolaire ; informatique : Word, Excel. \\ \hline
\end{tabularx}
\end{center}

\begin{enumerate}
\item Rédige la rubrique « Formation » du CV de Mlle FOTSO en respectant l'ordre antéchronologique.
\item Rédige la rubrique « Expérience professionnelle » en employant des énumérations à l'infinitif pour décrire les tâches (tu préciseras deux tâches plausibles par poste).
\item Rédige la rubrique « Langues et informatique » à partir de la rubrique « Divers » de la fiche.
\item Relis ta production : signale deux règles de présentation matérielle du CV que tu as respectées.
\end{enumerate}
\end{exercice}



\newpage

\coursec{Corrigés des exercices}

\courssub{Je vérifie mes connaissances}

\begin{corrige}[Exercice 1 — QCM : les rubriques du CV]
\begin{enumerate}
\item Réponse b : l'ordre \cle{antéchronologique} présente les informations de la plus récente à la plus ancienne, afin que l'employeur lise d'abord l'expérience ou le diplôme le plus actuel.
\item Réponse a : l'état civil (ou « Informations personnelles ») donne obligatoirement l'identité et les coordonnées du candidat (nom, prénoms, téléphone, adresse électronique) ; les loisirs relèvent des centres d'intérêt et le salaire n'apparaît jamais dans un CV.
\item Réponse b : les compétences s'énoncent en énumérations brèves avec des verbes à l'\cle{infinitif} (\emph{conduire, rédiger, assurer}), mode de la consigne et de l'énumération.
\item Réponse b : un CV efficace tient sur une page, deux au maximum, car l'employeur consacre peu de temps à chaque dossier.
\end{enumerate}
\end{corrige}

\begin{corrige}[Exercice 2 — Vrai ou faux ? Justifie ta réponse]
\begin{enumerate}
\item Vrai : le candidat adapte son CV à chaque offre en mettant en avant les compétences et les expériences qui correspondent au profil recherché par l'employeur.
\item Faux : la photo est facultative ; elle n'est jointe que si l'offre la demande explicitement.
\item Faux : le CV privilégie les énumérations brèves, les verbes à l'infinitif et les mots-clés, pour permettre une lecture rapide ; les longues phrases rédigées sont à proscrire.
\item Faux : dans un avis de recrutement, l'impératif exprime le plus souvent un ordre ou une consigne (\emph{Envoyez, Présentez-vous}), et parfois une interdiction (\emph{Ne vous présentez pas}), non un souhait.
\end{enumerate}
\end{corrige}

\begin{corrige}[Exercice 3 — Définitions]
\begin{enumerate}
\item Le \cle{CV} (curriculum vitae, « déroulement de la vie » en latin) est un texte fonctionnel qui présente de manière synthétique l'identité, la formation, l'expérience professionnelle et les compétences d'un candidat à un emploi. Il vise à obtenir un entretien d'embauche.
\item La \cle{structure externe} d'un CV désigne son organisation matérielle visible : le titre, la disposition des rubriques en blocs, les alinéas, les puces, la mise en page aérée sur une page.
\item La \cle{structure interne} d'un CV désigne l'organisation du contenu : le choix et l'ordre des rubriques (état civil, formation, expérience, compétences, langues, centres d'intérêt), l'ordre antéchronologique des informations et la cohérence avec le poste visé.
\item Un \cle{texte fonctionnel} est un texte utilitaire de la vie courante (CV, lettre, avis de recrutement, formulaire) qui vise un but pratique précis : informer, demander, obtenir un emploi, etc.
\end{enumerate}
\end{corrige}

\begin{corrige}[Exercice 4 — Repérage des modes]
\begin{enumerate}
\item \emph{Envoyez} : \cle{impératif}, valeur d'ordre (consigne adressée au candidat).
\item \emph{Maîtriser} : \cle{infinitif}, valeur de consigne (compétence exigée ; l'infinitif est ici sujet du verbe « est exigé »).
\item \emph{Présentez} : \cle{impératif}, valeur d'ordre.
\item \emph{assurer} et \emph{rédiger} : \cle{infinitifs}, valeur d'énumération des tâches du poste.
\end{enumerate}
Point de vigilance : ne pas confondre l'infinitif d'énumération (liste de tâches) avec l'impératif d'ordre (action exigée du lecteur) ; seul l'impératif s'adresse directement au candidat.
\end{corrige}

\courssub{Je m'entraîne}

\begin{corrige}[Exercice 5 — Remise en ordre d'un CV]
Ordre logique : 1. État civil ; 2. Titre du poste recherché ; 3. Formation ; 4. Expérience professionnelle ; 5. Compétences ; 6. Langues ; 7. Centres d'intérêt ; 8. Références.

Justification : le CV s'ouvre sur l'identification du candidat, puis annonce l'objectif (le poste visé). Viennent ensuite les informations essentielles pour l'employeur (formation, expérience, compétences), puis les informations complémentaires (langues, centres d'intérêt) et enfin les références, qui ferment le document en garantissant la véracité des informations.
\end{corrige}

\begin{corrige}[Exercice 6 — Correction d'un extrait de CV]
Fautes d'orthographe :
\begin{enumerate}
\item \emph{de Edéa} $\rightarrow$ \emph{d'Édéa} (élision devant voyelle et majuscule accentuée) ;
\item \emph{j'ai travailler} $\rightarrow$ \emph{j'ai travaillé} (participe passé après l'auxiliaire \emph{avoir}) ;
\item \emph{maintanance} $\rightarrow$ \emph{maintenance} ;
\item \emph{j'était} $\rightarrow$ \emph{j'étais} (imparfait, première personne) ;
\item \emph{acceuil} $\rightarrow$ \emph{accueil}.
\end{enumerate}
Incohérences de dates :
\begin{enumerate}
\item « juillet 2023 à janvier 2022 » : la date de fin est antérieure à la date de début ; corriger par exemple en « juillet 2022 à janvier 2023 » ;
\item le stage (2022-2023) se déroule pendant la période du Baccalauréat (2021-2024) : il faut préciser qu'il s'agit d'un stage de vacances ou ajuster les dates pour garantir la cohérence du parcours.
\end{enumerate}
Point de vigilance : dans un CV, toute incohérence de dates discrédite le candidat ; la relecture des dates est aussi importante que la correction orthographique.
\end{corrige}

\begin{corrige}[Exercice 7 — De la phrase à l'énumération]
\textbf{Compétences}
\begin{itemize}
\item Conduire les engins de chantier ;
\item Lire les plans de construction ;
\item Réparer les installations électriques ;
\item Tenir le registre du matériel.
\end{itemize}
Chaque compétence est exprimée par un verbe à l'\cle{infinitif} suivi de son complément : la formulation est brève, directe et facile à repérer pour l'employeur. On supprime les marques de personne (\emph{je}) et les connecteurs (\emph{en plus}), inutiles dans un CV.
\end{corrige}

\begin{corrige}[Exercice 8 — Impératif et infinitif dans un avis de recrutement]
\begin{enumerate}
\item Infinitifs : \emph{effectuer, respecter, rendre} (énumération des tâches) ; on peut ajouter \emph{déposer}, dont le futur simple « déposeront » exprime une consigne proche de l'infinitif. Impératifs : \emph{Contactez} ; \emph{présentez} (dans « Ne vous présentez pas »).
\item Valeurs : \emph{effectuer, respecter, rendre} : infinitifs d'énumération des tâches du poste ; \emph{Contactez} : impératif d'ordre (invitation pressante à agir) ; \emph{Ne vous présentez pas} : impératif exprimant l'interdiction.
\end{enumerate}
Point de vigilance : l'impératif n'exprime pas toujours un ordre positif ; à la forme négative, il peut exprimer une interdiction, valeur à citer explicitement.
\end{corrige}

\courssub{Je me prépare au Bac}

\begin{corrige}[Exercice 9 — Sujet type Baccalauréat technique : rédaction complète d'un CV]
\textbf{État civil} : MBALLA Arnaud, né le 5 mai 2004 à Douala ; tél. 6 91 22 33 44 ; arnaud.mballa@mail.com ; Douala-Bonabéri.

\textbf{Poste recherché} : technicien en électronique.

\textbf{Formation} (ordre antéchronologique) :
\begin{itemize}
\item 2023-2024 : Baccalauréat technique série F3 (électronique), lycée technique de Douala-Bassa ;
\item 2020-2021 : Probatoire technique F3, lycée technique de Douala-Bassa ;
\item 2019-2020 : BEPC, collège de la Cité des Palmiers.
\end{itemize}

\textbf{Expérience professionnelle} :
\begin{itemize}
\item Juillet-septembre 2023 : stage de technicien, atelier SORELEC (Douala) : dépanner des équipements électroniques ; rédiger les rapports d'intervention.
\end{itemize}

\textbf{Compétences} :
\begin{itemize}
\item Lire et interpréter des schémas électroniques ;
\item Assurer la maintenance des équipements ;
\item Rédiger des rapports d'intervention clairs ;
\item Utiliser les appareils de mesure (multimètre, oscilloscope).
\end{itemize}

\textbf{Langues} : français (courant), anglais (niveau scolaire).

\textbf{Centres d'intérêt} : robotique, football.

\textbf{Références} : M. KAMGA, chef d'atelier SORELEC, tél. 6 99 88 77 66.

\textbf{Barème indicatif} : prise en compte de la situation de communication (poste visé, adaptation à l'offre) : 4 pts ; structure du CV (rubriques complètes, ordre logique, dates antéchronologiques et cohérentes) : 6 pts ; correction de la langue (infinitifs d'énumération, orthographe, accords) : 6 pts ; présentation matérielle (blocs, puces, une page) : 4 pts.

\textbf{Points de vigilance} : reprendre dans les compétences les mots-clés de l'offre (\emph{schémas électroniques, maintenance, rapports d'intervention}) ; vérifier que les dates de formation précèdent l'expérience ; employer l'infinitif, jamais le passé composé, dans les énumérations.
\end{corrige}

\begin{corrige}[Exercice 10 — Adaptation du CV à deux situations d'embauche]
\textbf{Version pour l'offre A (comptable).}

Expérience professionnelle :
\begin{itemize}
\item Stage de trois mois, PME de Bafoussam : tenir les livres de comptes ; établir les factures ; classer les pièces comptables.
\item Emploi de vacances (deux mois), caissier : gérer la caisse ; vérifier les recettes journalières.
\end{itemize}
Compétences : tenir la comptabilité courante ; établir les factures ; utiliser un logiciel de gestion ; maîtriser les calculs commerciaux.

\textbf{Version pour l'offre B (vendeur).}

Expérience professionnelle :
\begin{itemize}
\item Emploi de vacances (deux mois), caissier : accueillir les clients ; encaisser les ventes ; ranger les rayons.
\item Stage de trois mois, PME de Bafoussam : gérer les stocks de fournitures ; renseigner les clients au téléphone.
\end{itemize}
Compétences : accueillir et conseiller la clientèle ; gérer la caisse ; organiser les rayons ; communiquer avec aisance.

\textbf{Justification} : un CV efficace sélectionne les informations en fonction du poste visé. Pour l'offre A, le stage comptable est placé en premier car il correspond directement au profil recherché ; l'emploi de caissier n'est conservé que pour ses aspects financiers (gestion de caisse). Pour l'offre B, c'est l'expérience de caissier qui devient centrale, car elle prouve le contact avec la clientèle ; le stage est reformulé pour valoriser l'accueil et la gestion des stocks. L'ordre des rubriques et le choix des verbes à l'infinitif reprennent les mots-clés de chaque offre, ce qui montre à l'employeur que le candidat correspond au poste.

\textbf{Points de vigilance} : ne jamais inventer d'expérience inexistante, mais reformuler les tâches réelles selon l'angle utile ; conserver les mêmes dates dans les deux versions (seul l'ordre et la formulation changent).
\end{corrige}

\begin{corrige}[Exercice 11 — Sujet type Probatoire : rubriques à partir d'une fiche d'identité professionnelle]
\begin{enumerate}
\item \textbf{Formation} (ordre antéchronologique) :
\begin{itemize}
\item 2021-2023 : BTS en secrétariat de direction, IST de Bafoussam ;
\item 2020-2021 : Baccalauréat série G1, lycée de Dschang ;
\item 2019-2020 : Probatoire série G1, lycée de Dschang.
\end{itemize}
\item \textbf{Expérience professionnelle} :
\begin{itemize}
\item Janvier 2023 - juin 2024 : secrétaire comptable, cabinet FOKAM et Associés (Bafoussam) : rédiger le courrier administratif ; tenir la comptabilité courante du cabinet.
\item Septembre 2022 - décembre 2022 : stage de secrétaire, mairie de Bafoussam : accueillir et orienter les usagers ; classer et archiver les dossiers.
\end{itemize}
\item \textbf{Langues et informatique} : français : courant ; anglais : niveau scolaire ; informatique : traitement de texte (Word), tableur (Excel).
\item Règles de présentation respectées (exemples) : les informations sont présentées en énumérations brèves et non en longues phrases ; les dates suivent l'ordre antéchronologique et sont alignées de façon homogène ; chaque rubrique porte un titre visible et l'ensemble tient sur une page.
\end{enumerate}
\textbf{Points de vigilance} : dans la fiche, la formation est donnée de la plus ancienne à la plus récente : il faut impérativement l'inverser pour respecter l'ordre antéchronologique ; les tâches inventées doivent rester plausibles et cohérentes avec l'intitulé du poste.
\end{corrige}

\newpage

\coursec{Fiche bilan}

\courssub{Synthèse visuelle de la leçon}

\begin{popfigure}
\begin{tikzpicture}[
    mindmap,
    concept/.style={circle, minimum size=3.2cm, text width=2.8cm, align=center, font=\bfseries\footnotesize, inner sep=4pt, execute at begin node=\hskip0pt},
    grow cyclic,
    level 1/.style={level distance=5.5cm, sibling angle=360/7},
    concept color=popblue,
    every node/.style={concept},
    root/.style={concept, concept color=popdark, minimum size=4cm, font=\bfseries\large},
    branch1/.style={concept color=popblue},
    branch2/.style={concept color=poporange},
    branch3/.style={concept color=popgreen},
    branch4/.style={concept color=poppurple},
    branch5/.style={concept color=poppink},
    branch6/.style={concept color=popgold},
    branch7/.style={concept color=popmuted}
]
\node[root] {Le CV}
    child[branch1] {node[concept, align=center]{Structure\\externe}}
    child[branch2] {node[concept, align=center]{Structure\\interne}}
    child[branch3] {node[concept, align=center]{Langue :\\impératif / infinitif}}
    child[branch4] {node[concept] {Rédaction et amélioration}}
    child[branch5] {node[concept] {Adaptation à l'offre}}
    child[branch6] {node[concept] {Textes supports}}
    child[branch7] {node[concept] {Évaluation et remédiation}};
\end{tikzpicture}
\legende{Carte mentale récapitulative de la leçon 12 : le CV, de la structure à l'adaptation à une situation d'embauche.}
\end{popfigure}

\courssub{Bilan des savoirs et savoir-faire}

\textbf{1. Définitions clés}\par
\begin{itemize}
    \item \cle{CV (Curriculum Vitae)} : document synthétique qui présente le parcours, la formation et les compétences d'un candidat en vue d'une \emph{situation d'embauche} donnée. C'est un \cle{texte fonctionnel} : il vise une action précise (obtenir un entretien).
    \item \cle{Structure externe} : tout ce qui concerne la présentation matérielle du document : mise en page, police, marges, longueur (idéalement une page), format du fichier, lisibilité visuelle.
    \item \cle{Structure interne} : l'organisation du contenu en rubriques logiques et ordonnées (état civil et coordonnées, accroche, expériences professionnelles, formation, compétences, centres d'intérêt).
    \item \cle{Adaptation} : personnalisation du CV en fonction des \cle{mots-clés} de l'offre d'emploi (compétences exigées, diplômes demandés, qualités recherchées).
\end{itemize}

\textbf{2. La structure externe (check-list visuelle)}\par
\begin{itemize}
    \item \textbf{Format} : feuille A4, orientation portrait, police classique et lisible (Arial, Calibri, Times New Roman) en corps 10 à 12.
    \item \textbf{Marges} : régulières (environ 1,5 à 2 cm) ; alignement à gauche ou justifié ; puces simples et identiques dans tout le document.
    \item \textbf{En-tête} : nom et prénom, intitulé du poste visé, coordonnées (téléphone, adresse électronique professionnelle, ville de résidence), photo d'identité (facultative, de bonne qualité).
    \item \textbf{Longueur} : une page pour un jeune candidat ; deux pages au maximum pour un candidat expérimenté.
    \item \textbf{Aération} : espaces entre les rubriques, titres de rubriques en gras, aucun paragraphe dense.
\end{itemize}

\textbf{3. La structure interne (rubriques et ordre)}\par
\begin{center}
\begin{tabularx}{\linewidth}{|l|X|}\hline
\rowcolor{popblueL}
\textbf{Rubrique} & \textbf{Contenu et ordre (ordre antéchronologique : du plus récent au plus ancien)} \\\hline
\cle{État civil et coordonnées} & Nom, prénom, date et lieu de naissance (facultatifs), téléphone, adresse électronique, ville. \\\hline
\cle{Accroche / objectif} & Deux ou trois lignes : poste visé, compétence principale, valeur ajoutée pour l'employeur. \\\hline
\cle{Expériences professionnelles} & Poste occupé, entreprise, dates, lieu, puis trois à cinq missions ou résultats rédigés à l'\textbf{infinitif} (style télégraphique). \\\hline
\cle{Formation} & Intitulé exact du diplôme, établissement, année d'obtention, spécialité ou mention. \\\hline
\cle{Compétences} & Compétences \textbf{techniques} (logiciels, machines, procédés, normes) et compétences \textbf{transversales} (organisation, communication, langues). \\\hline
\cle{Centres d'intérêt} & Choisis avec soin : sport collectif (esprit d'équipe), bénévolat (sens de l'engagement), lecture technique (curiosité professionnelle). \\\hline
\end{tabularx}
\end{center}

\textbf{4. Langue française : les valeurs de l'impératif et de l'infinitif (morphosyntaxe)}\par

\begin{definition}[Rappels grammaticaux]
L'\cle{impératif} est le mode de l'ordre, du conseil ou de la défense ; il ne possède que trois personnes (2\up{e} personne du singulier, 1\up{re} et 2\up{e} personnes du pluriel) et se construit sans sujet exprimé. L'\cle{infinitif} est un mode impersonnel : il nomme l'action sans l'attacher à une personne ni à un temps précis.
\end{definition}

\begin{center}
\begin{tabularx}{\linewidth}{|l|X|X|}\hline
\rowcolor{popgreenL}
\textbf{Mode} & \textbf{Valeurs principales} & \textbf{Exemples} \\\hline
\cle{Impératif} & Ordre, conseil, instruction (mode d'emploi, consignes). Fréquent dans les offres d'emploi et les lettres, rare dans le CV lui-même. & \emph{« Envoyez votre CV avant le 30 juin. »} (consigne d'une annonce) \\\hline
\cle{Infinitif} & \textbf{Valeur nominale} : l'infinitif tient lieu de nom (titre de rubrique). \textbf{Valeur verbale} : après une préposition (\emph{pour}, \emph{sans}, \emph{après avoir}). \textbf{Style énumératif} : c'est la norme du CV pour la liste des missions. & \emph{Expériences professionnelles} (titre) ; \emph{Pour obtenir un entretien} ; \emph{Gérer une équipe de cinq personnes} ; \emph{Après avoir validé le BTS} \\\hline
\end{tabularx}
\end{center}

\begin{attention}[Piège fréquent]
Ne pas mélanger les modes dans une même liste de missions. Pour la cohérence, choisir \textbf{l'infinitif} (style télégraphique) : \emph{Gérer, Organiser, Développer, Assurer}. Éviter le passé composé (\emph{J'ai géré}) et l'impératif (\emph{Gérez}) dans les rubriques du CV.
\end{attention}

\textbf{5. Méthode : rédiger et améliorer son CV en cinq étapes}\par

\begin{methode}[De l'offre d'emploi au CV finalisé]
\begin{enumerate}
    \item \textbf{Analyser l'offre} : souligner les mots-clés (compétences exigées, diplômes, qualités recherchées).
    \item \textbf{Collecter et trier} : lister toutes ses expériences et formations, puis ne retenir que ce qui correspond à l'offre.
    \item \textbf{Rédiger les rubriques} : employer l'infinitif pour les missions ; choisir des verbes d'action précis (\cle{Concevoir}, \cle{Organiser}, \cle{Réparer}, \cle{Installer} plutôt que \emph{faire}, \emph{aider}) ; chiffrer les résultats quand c'est possible (\emph{réduction des pannes de 15 \%}).
    \item \textbf{Mettre en forme} : appliquer la structure externe (alignement des dates, intitulés en gras, espaces entre rubriques).
    \item \textbf{Relire et valider} : corriger l'orthographe, vérifier la cohérence des dates, nommer le fichier clairement (\emph{Nom\_Prenom\_Poste.pdf}) et le faire relire par une autre personne.
\end{enumerate}
\end{methode}

\textbf{6. Adapter le CV à une situation d'embauche (spécialités techniques)}\par
\begin{itemize}
    \item \textbf{Demande de stage} : placer la \cle{Formation} et les \cle{travaux pratiques} avant les expériences ; détailler les matières, les ateliers et les logiciels étudiés.
    \item \textbf{Premier emploi de technicien} : valoriser les \cle{stages} en entreprise, le \cle{projet de fin d'études} et les interventions réelles effectuées (maintenance, installation, câblage).
    \item \textbf{Candidature spontanée} : rédiger une accroche ciblée sur le secteur d'activité de l'entreprise et mettre en avant les réalisations concrètes personnelles.
\end{itemize}

\begin{exemplebox}[Extrait de rubrique « Expériences » bien rédigée]
\textbf{Stagiaire technicien de maintenance} — Société X, Douala, juillet--septembre 2025
\begin{itemize}
    \item Assurer la maintenance préventive des machines de l'atelier ;
    \item Diagnostiquer les pannes électriques et proposer des solutions ;
    \item Rédiger les rapports d'intervention quotidiens.
\end{itemize}
Remarque : les trois puces sont à l'infinitif, les verbes sont précis, les dates sont indiquées.
\end{exemplebox}

\courssub{Checklist avant le Probatoire / Baccalauréat technique}

\begin{aretenir}[Checklist : mon CV est-il prêt ?]
\begin{itemize}
    \item $\square$ \textbf{Titre clair} : poste visé et spécialité (ex. \emph{Technicien de maintenance industrielle}).
    \item $\square$ \textbf{Coordonnées professionnelles} : adresse électronique sérieuse (\emph{prenom.nom@...}), téléphone, ville.
    \item $\square$ \textbf{Accroche ciblée} : deux lignes qui répondent directement au besoin de l'employeur.
    \item $\square$ \textbf{Expériences antéchronologiques} : dates indiquées, verbes à \textbf{l'infinitif}, résultats concrets.
    \item $\square$ \textbf{Formation lisible} : intitulé exact du diplôme, établissement, année, spécialité.
    \item $\square$ \textbf{Compétences techniques séparées} : logiciels, machines, procédés, avec le niveau (maîtrise, notions).
    \item $\square$ \textbf{Langues et attestations} : niveau précis (bon niveau, courant), attestations professionnelles obtenues.
    \item $\square$ \textbf{Mise en page impeccable} : une page, une seule police, alignements réguliers, aucun espace vide disgracieux.
    \item $\square$ \textbf{Zéro faute} : relecture à voix haute, accords vérifiés, homonymes contrôlés.
    \item $\square$ \textbf{Cohérence CV / lettre de motivation} : mêmes informations, mêmes dates, même présentation.
    \item $\square$ \textbf{Nom de fichier professionnel} : \emph{Nom\_Prenom\_Poste.pdf} (et non \emph{CV\_final\_v2.pdf}).
    \item $\square$ \textbf{Adaptation finale} : les mots-clés de l'offre visée sont intégrés naturellement dans le CV.
\end{itemize}
\end{aretenir}

\findecours

\end{document}
